% =====================================================================
%  CHƯƠNG 3. THIẾT KẾ BACKEND
%  Nội dung chia thành các tệp con trong chapters/c3/ để dễ cộng tác.
% =====================================================================
\chapter{THIẾT KẾ BACKEND}\label{chap:backend}

Chương~\ref{chap:dac-ta} đã xác định hệ thống cần làm gì thông qua 25 use case, 23 quy tắc nghiệp vụ và 23 yêu cầu phi chức năng. Chương này trình bày cách phần máy chủ (backend) hiện thực các yêu cầu đó: lựa chọn công nghệ, kiến trúc phần mềm, thiết kế chi tiết các thành phần, cơ sở dữ liệu, giao diện lập trình ứng dụng (API), bảo mật và phương án triển khai. Ứng dụng di động được xây dựng bằng Flutter và chạy trên cả Android lẫn iOS (\nfr{NFR19}), vì vậy các quyết định ở phía máy chủ đều tính tới đặc thù của hai nền tảng: thiết bị thường xuyên mất kết nối khi người dùng đi chợ, hệ điều hành hạn chế ứng dụng chạy nền, và phiên bản ứng dụng trên kho ứng dụng được cập nhật chậm hơn máy chủ.

Hoạt động ngoại tuyến và đồng bộ thời gian thực (\nfr{NFR10}, \nfr{NFR11}) là định hướng phát triển: các mục tương ứng được đánh dấu \emph{Phạm vi} và không thuộc tiêu chí chấp nhận. Quét mã vạch không thuộc phạm vi nên không có trong thiết kế.

Thiết kế được trình bày theo hướng từ tổng quát đến chi tiết. Mục~\ref{sec:yeu-cau-kien-truc} tổng hợp các yếu tố dẫn dắt kiến trúc rút ra từ yêu cầu phi chức năng. Mục~\ref{sec:cong-nghe} giới thiệu nền tảng công nghệ được lựa chọn. Mục~\ref{sec:kien-truc} trình bày kiến trúc tổng thể theo mô hình C4~\cite{c4model} ở mức ngữ cảnh, mức vùng chứa và mức thành phần; mục~\ref{sec:tich-hop-di-dong} đối chiếu thiết kế với kiến trúc tham chiếu chính thức của Android, iOS và Flutter để bảo đảm backend phục vụ đúng chuẩn ứng dụng di động; mục~\ref{sec:thiet-ke-chi-tiet} đi xuống mức lớp, luồng xử lý, máy trạng thái, các kỹ thuật cốt lõi và cách hiện thực hóa từng use case trọng tâm. Ba mục tiếp theo đặc tả cơ sở dữ liệu, API và bảo mật. Mục~\ref{sec:trien-khai} mô tả hạ tầng triển khai, khả năng sẵn sàng và lộ trình mở rộng. Mục~\ref{sec:truy-vet-thiet-ke} khép lại chương bằng ma trận truy vết từ yêu cầu tới quyết định thiết kế và kế hoạch kiểm chứng. Các con số về quy mô trong chương là giả định thiết kế, dùng để định cỡ hệ thống, không phải số liệu đo thực tế.

% Sơ đồ của chương này cao: dùng \usetalldiagrams (khai báo trong config/preamble.tex).
\begingroup
\usetalldiagrams

% ---------------------------------------------------------------------
%  3.1. Yêu cầu kiến trúc và định hướng thiết kế
% ---------------------------------------------------------------------
\section{Yêu cầu kiến trúc và định hướng thiết kế}\label{sec:yeu-cau-kien-truc}

Mục này xác định các yêu cầu quyết định hình dạng của backend, mức tải phải chịu và các nguyên tắc dùng làm thước đo cho những lựa chọn ở các mục sau.

\subsection{Các yếu tố dẫn dắt kiến trúc}\label{ssec:yeu-to-dan-dat}

Theo Bass và cộng sự~\cite{bass}, yêu cầu có ý nghĩa kiến trúc (architecturally significant requirement – ASR) là yêu cầu mà khi thay đổi buộc phải đổi cách chia thành phần, lưu dữ liệu hay giao tiếp; phần lớn chúng là thuộc tính chất lượng chứ không phải chức năng. “Đánh dấu đã mua” hiện thực được bằng nhiều kiến trúc, nhưng đòi hỏi thao tác đó chạy khi mất mạng và hiện trên máy người khác trong 3 giây thì loại bỏ phần lớn kiến trúc đơn giản. Mỗi yêu cầu chất lượng được diễn đạt thành kịch bản thuộc tính chất lượng (quality attribute scenario) gồm sáu phần: nguồn kích thích, kích thích, đối tượng chịu tác động, môi trường, phản ứng và thước đo phản ứng; mỗi kịch bản được đáp ứng bằng các chiến thuật (tactic), tức quyết định thiết kế đã được kiểm nghiệm.

Bảng~\ref{tab:c3-kich-ban-chat-luong} chuyển các yêu cầu phi chức năng ở mục~\ref{sec:phi-chuc-nang} thành mười kịch bản KB01–KB10; nguồn và đối tượng chịu tác động được gộp vào cột kích thích, thước đo lấy theo Bảng~\ref{tab:tong-hop-nfr}. Các yêu cầu \nfr{NFR01}--\nfr{NFR08} thuộc về giao diện nên không có kịch bản riêng.

{\small
\begin{longtable}{|C{1.0cm}|L{2.2cm}|L{3.6cm}|L{3.3cm}|L{\dimexpr\textwidth-10\tabcolsep-6\arrayrulewidth-1.0cm-2.2cm-3.6cm-3.3cm\relax}|}
\caption{Kịch bản thuộc tính chất lượng rút ra từ yêu cầu phi chức năng}\label{tab:c3-kich-ban-chat-luong}\\
\hline
\hd{Mã} & \hd{Thuộc tính} & \hd{Kích thích, môi trường} & \hd{Phản ứng, thước đo} & \hd{Chiến thuật (mục)}\\ \hline
\endfirsthead
\multicolumn{5}{l}{\footnotesize\itshape Bảng \thetable\ (tiếp theo)}\\ \hline
\hd{Mã} & \hd{Thuộc tính} & \hd{Kích thích, môi trường} & \hd{Phản ứng, thước đo} & \hd{Chiến thuật (mục)}\\ \hline
\endhead
KB01 & Hiệu năng (\nfr{NFR09}) & Giờ cao điểm, tải đỉnh 50 yêu cầu/giây & 95\% API $\le$ 1 s & Chỉ mục, keyset, cache Redis, worker (\ref{sec:csdl})\\ \hline
KB02 & Hiệu năng (\nfr{NFR09}) & Gõ tên không dấu, sai chính tả; xin gợi ý món & Gợi ý tên $\le$ 500 ms; gợi ý món $\le$ 2 s & Elasticsearch, edge n-gram (\ref{sec:cong-nghe})\\ \hline
KB03 & Thời gian thực (\nfr{NFR10}) & A đánh dấu đã mua khi B đang mở cùng danh sách & B ở tiền cảnh thấy thay đổi $\le$ 3 s & Pub/Sub, WebSocket, kéo phần thay đổi; FCM khi ở nền (\ref{sec:thiet-ke-chi-tiet})\\ \hline
KB04 & Ngoại tuyến (\nfr{NFR11}) & Mất sóng khi đi chợ; có mạng lại, yêu cầu bị gửi lặp & Áp dụng đúng một lần; xung đột theo quy tắc & ID do thiết bị sinh, khóa lũy đẳng (\ref{sec:thiet-ke-chi-tiet})\\ \hline
KB05 & Toàn vẹn (\nfr{NFR11}, \nfr{NFR13}) & Hai người cùng xác nhận nấu (\uc{UC19}) hoặc nhập tủ (\uc{UC11}) & Không lô âm, không trừ hay nhập trùng; nguyên tử & Một giao dịch, khóa dòng theo thứ tự cố định (\ref{sec:kien-truc})\\ \hline
KB06 & Nhắc nhở (\nfr{NFR12}) & Đến giờ nhắc; worker khởi động lại, FCM lỗi tạm thời & Tối đa một thông báo nhắc mỗi người mỗi ngày; SLI tạo: $\ge$ 99\% thông báo được tạo trong 15 phút sau giờ nhắc; SLI gửi: tỉ lệ giao tới FCM, thử lại 3 lần & Tác vụ định kỳ mỗi phút, tác vụ duy nhất theo người và ngày, khóa chống trùng, backoff (\ref{sec:thiet-ke-chi-tiet})\\ \hline
KB07 & Sẵn sàng (\nfr{NFR13}) & Triển khai bản mới, API lỗi, mất dữ liệu & $\ge$ 99\% lúc 06:00–23:00; RPO $\le$ 5 phút, RTO $\le$ 1 giờ & Hai bản API, cuốn chiếu, WAL liên tục (\ref{sec:trien-khai})\\ \hline
KB08 & Bảo mật (\nfr{NFR14}--\nfr{NFR17}) & Sửa mã không gian trong URL; dò mật khẩu, OTP; dùng lại token & Từ chối sai phạm vi; khóa tạm 15 phút; thu hồi họ phiên & Kiểm tra thành viên, Argon2id, giới hạn tần suất (\ref{sec:bao-mat})\\ \hline
KB09 & Riêng tư (\nfr{NFR18}) & Người dùng yêu cầu xóa tài khoản & Thu hồi phiên; xóa dữ liệu cá nhân, ảnh & Tác vụ nền dọn dữ liệu (\ref{sec:bao-mat})\\ \hline
KB10 & Tương thích, bảo trì (\nfr{NFR19}--\nfr{NFR21}) & Đổi API khi bản ứng dụng cũ còn chạy; đổi tham số & Bản quá cũ nhận HTTP 426; tham số áp dụng ngay & Phiên bản trong URL, OpenAPI, cấu hình trong CSDL (\ref{sec:api})\\ \hline
\end{longtable}
}

Theo cách xếp hạng trên cây tiện ích (utility tree) của Bass, KB03, KB04 và KB05 cao cả về tầm quan trọng lẫn rủi ro kỹ thuật nên chi phối lựa chọn kiến trúc ở mục~\ref{sec:kien-truc}. Các kịch bản có lúc kéo ngược nhau: khóa dòng của KB05 giảm mức song song cho KB01, còn ghi ngoại tuyến của KB04 buộc máy chủ đôi khi từ chối thao tác người dùng tưởng đã xong; nhóm giữ phạm vi khóa hẹp và trả lỗi xung đột có cấu trúc. Mục tiêu 99\% của KB07 trong khung 17 giờ mỗi ngày cho phép khoảng 5 giờ gián đoạn mỗi tháng, nên một máy chủ duy nhất có thời gian khôi phục một giờ vẫn đáp ứng được ở giai đoạn đầu. Thước đo của KB06 được tách thành hai chỉ số mức dịch vụ (service level indicator – SLI)~\cite{sre}: việc tạo thông báo trong hộp thư nằm trọn trong tầm kiểm soát của hệ thống nên gắn với ngưỡng 99\% trong 15 phút của \nfr{NFR12}, còn việc giao tới thiết bị phụ thuộc FCM, APNs và trạng thái máy nên được đo, cảnh báo riêng để sự cố của bên ngoài không làm sai lệch đánh giá thiết kế. \nfr{NFR21}--\nfr{NFR23} là ràng buộc (constraint) đã định sẵn. Ngoài ra, nền tảng di động áp lên backend ba điều kiện ít gặp ở ứng dụng web: mạng chập chờn, hạn chế chạy nền và cập nhật ứng dụng chậm (mục~\ref{ssec:cn-di-dong}).

\subsection{Giả định quy mô và ước lượng tải}\label{ssec:uoc-luong-tai}

Hệ thống chưa có người dùng thật nên các con số trong Bảng~\ref{tab:c3-uoc-luong-tai} là giả định thiết kế, không phải số liệu đo. Năm đầu có 5.000 tài khoản (khoảng 2.000 không gian gia đình), 30\% hoạt động mỗi ngày với khoảng 60 yêu cầu API mỗi người. Tải dồn vào khung 17:00–19:00 khi các gia đình đi chợ, nấu ăn; giờ bận nhất chiếm khoảng 20\% lượng yêu cầu cả ngày, và hệ số an toàn 10 bù cho đột biến như nhiều thiết bị cùng gửi hàng đợi ngoại tuyến. Thông báo nhắc hạn dùng được gộp mọi không gian thành tối đa một thông báo cho mỗi người mỗi ngày (\br{BR13}), nên số thông báo nhắc bị chặn trên bởi số tài khoản. Định luật Little $L = \lambda W$ với $W = 1$ s cho số yêu cầu đồng thời để định cỡ bể kết nối CSDL; dung lượng PostgreSQL tính với khoảng mười bản ghi lâu dài, mỗi bản ghi khoảng 300 byte kể cả chỉ mục, cho mỗi người dùng hoạt động mỗi ngày.

{\small
\begin{longtable}{|L{4.6cm}|L{\dimexpr\textwidth-6\tabcolsep-4\arrayrulewidth-4.6cm-3.6cm\relax}|L{3.6cm}|}
\caption{Giả định thiết kế và ước lượng tải năm đầu}\label{tab:c3-uoc-luong-tai}\\
\hline
\hd{Đại lượng} & \hd{Giả định hoặc công thức} & \hd{Giá trị}\\ \hline
\endfirsthead
\multicolumn{3}{l}{\footnotesize\itshape Bảng \thetable\ (tiếp theo)}\\ \hline
\hd{Đại lượng} & \hd{Giả định hoặc công thức} & \hd{Giá trị}\\ \hline
\endhead
Người dùng hoạt động/ngày & 30\% × 5.000 tài khoản & 1.500\\ \hline
Yêu cầu API mỗi ngày & 1.500 × 60 & Khoảng 90.000\\ \hline
Tải giờ cao điểm & 0,2 × 90.000 / 3.600 s & 5 yêu cầu/giây\\ \hline
Tải đỉnh thiết kế & 5 × 10 (hệ số an toàn) & 50 yêu cầu/giây\\ \hline
Yêu cầu xử lý đồng thời & \emph{L} = 50 × 1 s (định luật Little) & $\le$ 50\\ \hline
Kết nối WebSocket & Cận trên: số người dùng hoạt động & $\le$ 1.500\\ \hline
Thông báo nhắc mỗi ngày & 5.000 tài khoản × tối đa 1 thông báo gộp mọi không gian (\br{BR13}) & $\le$ 5.000\\ \hline
Tốc độ tạo thông báo nhắc & 5.000 / 900 s (15 phút, \nfr{NFR12}) & Khoảng 6 thông báo/giây\\ \hline
Dữ liệu PostgreSQL/năm & 1.500 × 10 × 300 B × 365 & Khoảng 1,6 GB (cận trên 2 GB)\\ \hline
Ảnh mỗi năm & 20.000 ảnh × 200 KB sau nén & Khoảng 4 GB\\ \hline
\end{longtable}
}

Tải này nằm gọn trong một máy chủ ảo nhiều nhân nên giai đoạn đầu chỉ cần một máy chủ (mục~\ref{sec:trien-khai}). Vì giả định có độ bất định cao, kiến trúc vẫn giữ sẵn đường mở rộng ngang với chi phí gần như bằng không: tiến trình API không giữ trạng thái phiên, phát tán thay đổi giữa các bản API đi qua Redis, tác vụ nền chạy ở worker riêng và truy vấn báo cáo có thể chuyển sang bản sao chỉ đọc.

\subsection{Nguyên tắc thiết kế}\label{ssec:nguyen-tac}

Bảy nguyên tắc sau là tiêu chí khi phải chọn giữa các phương án ở những mục tiếp theo. \textbf{PostgreSQL là nguồn sự thật duy nhất.} Dữ liệu nghiệp vụ được ghi vào PostgreSQL trong giao dịch; chỉ mục Elasticsearch, cache Redis và CSDL trên thiết bị là dẫn xuất dựng lại được, ảnh trong kho đối tượng có siêu dữ liệu do PostgreSQL quản lý, nên sự cố ở kho phụ chỉ làm suy giảm hiệu năng chứ không làm sai dữ liệu. \textbf{Máy chủ quyết định quyền và quy tắc nghiệp vụ.} Ứng dụng kiểm tra dữ liệu để phản hồi nhanh, máy chủ kiểm tra lại toàn bộ; ẩn nút trên giao diện không thay thế phân quyền (\br{BR23}, \nfr{NFR16}). \textbf{Mọi thao tác ghi lũy đẳng.} Khóa lũy đẳng ở API, ràng buộc duy nhất trong CSDL và khóa chống trùng ở tác vụ nền khiến cơ chế giao ít nhất một lần (at-least-once) cho hiệu quả như xử lý đúng một lần.

\textbf{Thiết kế cho mạng chập chờn.} Thiết bị di động là khách hàng chính (mobile-first): tạo được dữ liệu khi ngoại tuyến, chỉ kéo phần thay đổi; mọi lời gọi giữa các thành phần có thời hạn chờ. \textbf{Đơn giản trước, ranh giới rõ.} Chỉ thêm thành phần khi có yêu cầu cụ thể; hệ thống bắt đầu bằng một khối triển khai duy nhất~\cite{fowlermonolith} nhưng chia module với giao diện rõ ràng để có thể tách dịch vụ về sau.

\textbf{Quan sát được.} Mỗi yêu cầu mang một mã định danh xuyên suốt nhật ký, số đo và trace; mỗi NFR có số đo đi kèm, như thời gian phản hồi theo route cho \nfr{NFR09}. \textbf{Cấu hình theo Twelve-Factor.} Cấu hình triển khai đọc từ biến môi trường, tiến trình không trạng thái, một image cho mọi môi trường~\cite{twelvefactor}; tham số nghiệp vụ nằm trong CSDL, do quản trị viên đổi qua \uc{UC25} (Bảng~\ref{tab:tham-so-he-thong}).

% ---------------------------------------------------------------------
%  3.2. Nền tảng công nghệ sử dụng
% ---------------------------------------------------------------------
\section{Nền tảng công nghệ sử dụng}\label{sec:cong-nghe}

Mục này trình bày công nghệ được chọn từ thiết bị tới hạ tầng, tập trung vào đặc tính phục vụ các kịch bản ở mục~\ref{ssec:yeu-to-dan-dat}; mục~\ref{ssec:cn-so-sanh} so sánh các phương án và tổng hợp lựa chọn.

\subsection{Ứng dụng di động đa nền tảng và các yêu cầu đặt lên backend}\label{ssec:cn-di-dong}

Ứng dụng dùng Flutter~\cite{flutter} 3 (Dart 3) cùng Riverpod, go\_router và dio, biên dịch trước thành mã máy và tự vẽ giao diện nên một mã nguồn chạy nhất quán trên Android từ 8.0 và iOS từ 13 (\nfr{NFR19}). Dữ liệu cục bộ nằm trong SQLite qua drift, vừa là bản sao các không gian vừa là hàng đợi lệnh ghi chờ gửi (\nfr{NFR11}, \nfr{NFR23}); refresh token được flutter\_secure\_storage cất trong Keychain của iOS và Keystore của Android (\nfr{NFR15}). Thông báo đẩy dùng Firebase Cloud Messaging HTTP v1~\cite{fcm}: máy chủ gọi một API cho cả hai nền tảng, FCM tự chuyển tiếp qua APNs~\cite{apns} tới thiết bị iOS.

Đặc thù của thiết bị di động đặt lên backend bốn yêu cầu. Thứ nhất, iOS tạm dừng ứng dụng vài giây sau khi chuyển nền và Android hoãn truy cập mạng bằng chế độ Doze, nên WebSocket~\cite{rfc6455} chỉ dùng ở tiền cảnh; khi ứng dụng ở nền, máy chủ đẩy qua FCM thay vì giữ kết nối hay chờ thiết bị thăm dò (polling). Thứ hai, bản cập nhật phải qua xét duyệt của kho ứng dụng và tùy người dùng cài đặt, nên API được phiên bản hóa trong đường dẫn (\path{/it4788/v1}) và chỉ thay đổi theo hướng bổ sung; ứng dụng có \texttt{X-App-Version} thấp hơn mức tối thiểu nhận HTTP 426, mã 4260. Thứ ba, mạng chập chờn làm mất phản hồi của lệnh đã thành công, nên mọi lệnh ghi của người dùng đã đăng nhập mang \texttt{Idempotency-Key}~\cite{idempotencykey} và thực thể dùng UUIDv7~\cite{rfc9562} do thiết bị sinh; lệnh phát lại sau thời gian ngoại tuyến còn kèm \texttt{X-Client-Time} để máy chủ kiểm quy tắc ngày theo thời điểm thao tác thật. Thứ tư, để tiết kiệm pin và dữ liệu, thiết bị chỉ kéo phần thay đổi (delta) qua \path{/spaces/{sid}/changes?since=}.

\subsection{Ngôn ngữ Go và hệ thư viện phía máy chủ}\label{ssec:cn-go}

Backend viết bằng Go~\cite{golang} 1.23, biên dịch thành một tệp thực thi tĩnh chạy ở hai chế độ \texttt{api} và \texttt{worker} trên image distroless nhỏ, khởi động nhanh. Mỗi yêu cầu hay kết nối WebSocket chạy trong một goroutine, luồng nhẹ ánh xạ lên số ít luồng hệ điều hành, nên hàng nghìn kết nối không cần lập trình bất đồng bộ phức tạp; gói \texttt{context} truyền thời hạn và tín hiệu hủy, nên khi thiết bị ngắt kết nối thì truy vấn CSDL và lời gọi ra ngoài cũng dừng theo.

Các thư viện được chọn theo hướng bám sát \texttt{net/http} và sinh mã lúc biên dịch thay vì phản chiếu (reflection) lúc chạy: chi v5 định tuyến, coder/websocket xử lý WebSocket, pgx v5 truy cập PostgreSQL với giao dịch tường minh, \texttt{log/slog} ghi nhật ký JSON. Thay cho ORM, sqlc~\cite{sqlc} kiểm tra các câu SQL viết tay với lược đồ rồi sinh hàm Go an toàn kiểu; goose quản lý di trú bằng tệp SQL có phiên bản theo chiến lược mở rộng–thu hẹp (expand–contract); từ hợp đồng OpenAPI 3.1~\cite{openapi}, oapi-codegen sinh giao diện máy chủ và openapi-generator sinh client Dart, nên lệch hợp đồng thành lỗi biên dịch (\nfr{NFR20}). Mã nguồn~\ref{code:c3-sqlc-vi-du} minh họa sqlc với truy vấn khóa lô theo FEFO (\br{BR15}) của \uc{UC13} và \uc{UC19}.

\codecaption{Truy vấn sqlc khóa các lô theo FEFO, hàm Go được sinh và lời gọi truyền ngày nghiệp vụ}{code:c3-sqlc-vi-du}
\begin{Verbatim}
-- name: LockActiveLotsForFood :many
SELECT id, remaining_qty, unit_id, unit_factor, base_remaining,
       expires_on, version
FROM   pantry.food_lots
WHERE  space_id = @space_id AND food_id = @food_id
  AND  lifecycle = 'active' AND remaining_qty > 0
  AND  expires_on >= @today::date   -- ngày VN do platform/clock cấp
  AND  deleted_at IS NULL
ORDER  BY expires_on, id      -- thứ tự FEFO, cũng là thứ tự lấy khóa
FOR UPDATE;

// Do sqlc sinh ra; Params gồm SpaceID, FoodID (uuid.UUID), Today
func (q *Queries) LockActiveLotsForFood(ctx context.Context,
    arg LockActiveLotsForFoodParams) ([]LockActiveLotsForFoodRow, error)

// Tầng ứng dụng truyền ngày nghiệp vụ, không dùng current_date
lots, err := q.LockActiveLotsForFood(ctx, LockActiveLotsForFoodParams{
    SpaceID: sid, FoodID: fid, Today: clock.Today(), // Asia/Ho_Chi_Minh
})
\end{Verbatim}

Chú thích \texttt{-- name: ... :many} đặt tên hàm và báo truy vấn trả nhiều dòng; tham số có tên (\texttt{@space\_id}, \texttt{@food\_id}, \texttt{@today}) trở thành trường của kiểu tham số, sqlc sinh thêm kiểu dòng kết quả và phương thức trên \texttt{Queries}, và bước kiểm tra trong CI phát hiện truy vấn lệch lược đồ trước khi hợp nhất mã. Ngày so sánh với \texttt{expires\_on} được truyền vào thay vì dùng \texttt{current\_date}: phiên CSDL chạy theo UTC, nên từ 00:00 đến 07:00 giờ Việt Nam \texttt{current\_date} vẫn là ngày hôm trước và một lô vừa quá hạn sẽ bị trừ kho, trái \br{BR15}. Vì vậy mọi phép so sánh ngày nghiệp vụ trong hệ thống dùng tham số \texttt{\$today} do gói \texttt{platform/clock} tính theo múi giờ Asia/Ho\_Chi\_Minh, và kiểm thử có thể cố định đồng hồ để tái hiện các mốc nửa đêm. Mệnh đề \texttt{FOR UPDATE}, thứ ORM thường che khuất, được giữ nguyên văn; thứ tự \texttt{expires\_on, id} cố định giúp tránh khóa chết (deadlock).

\subsection{Hệ quản trị cơ sở dữ liệu PostgreSQL}\label{ssec:cn-postgresql}

PostgreSQL~\cite{postgresql} 16 được chọn trước hết vì giao dịch ACID: \uc{UC11} và \uc{UC19} ghi vào bảng của nhiều module và chỉ đúng khi mọi bước cùng thành công (\nfr{NFR13}). Cơ chế đa phiên bản (MVCC) cho phép truy vấn đọc, kể cả báo cáo dài, không chặn lệnh ghi; chỗ cần tuần tự hóa dùng khóa dòng tường minh. Chỉ mục từng phần (partial index) vừa thu nhỏ chỉ mục FEFO vừa biến quy tắc thành ràng buộc, chẳng hạn chỉ mục duy nhất có điều kiện bảo đảm không nhóm nào có hai trưởng nhóm (\br{BR05}) kể cả khi hai yêu cầu chạy đồng thời. Cột sinh (generated column) tính sẵn lượng quy về đơn vị cơ sở, JSONB chứa dữ liệu bán cấu trúc như các bước công thức, extension citext, unaccent và pg\_trgm hỗ trợ so khớp email và tìm kiếm không dấu, mỗi module sở hữu một schema riêng. Cuối cùng, \texttt{FOR UPDATE SKIP LOCKED} cho nhiều tiến trình lấy các dòng chưa bị khóa mà không chờ nhau, nền tảng để dùng chính PostgreSQL làm hàng đợi.

\subsection{Bộ nhớ đệm Redis}\label{ssec:cn-redis}

Redis~\cite{redis} 7 là kho khóa–giá trị trong bộ nhớ, độ trễ dưới mili giây, chỉ giữ dữ liệu tạm có thời hạn sống (TTL) và không là nguồn sự thật: mất Redis chỉ làm hệ thống chậm đi và tạm nới giới hạn tần suất. Redis có bốn vai trò: cache dữ liệu đọc nhiều, ít đổi (cấu hình, danh mục, vai trò thành viên, kết quả gợi ý món); giới hạn tần suất bằng kịch bản Lua chạy nguyên tử nên nhiều bản API không thể cùng để lọt yêu cầu vượt ngưỡng; bộ đếm thời gian chờ và số lần gửi lại OTP (\br{BR02}); và Pub/Sub phát tán thay đổi tới mọi bản API đang giữ WebSocket. Pub/Sub chỉ giao nhiều nhất một lần, chấp nhận được vì thông điệp chỉ là gợi ý để thiết bị kéo dữ liệu.

\subsection{Máy tìm kiếm Elasticsearch}\label{ssec:cn-elasticsearch}

Elasticsearch~\cite{elasticsearch} 8.15 lưu chỉ mục đảo (inverted index) ánh xạ mỗi từ tố tới tài liệu chứa nó. Bộ phân tích (analyzer) dùng asciifolding để bỏ dấu, đổi “đ” thành “d” cho “ca rot” khớp “cà rốt”, bộ lọc đồng nghĩa ghép “heo” với “lợn”, trường con edge n-gram lưu sẵn tiền tố để gợi ý khi đang gõ; fuzziness chấp nhận sai một, hai ký tự và BM25 xếp hạng độ liên quan. PostgreSQL với pg\_trgm cũng tìm không dấu được, nhưng khó cùng lúc chịu lỗi gõ theo tiền tố, xếp hạng có từ đồng nghĩa và đạt gợi ý tên $\le$ 500 ms (\nfr{NFR09}) mà không tăng tải lên CSDL giao dịch. Chỉ mục chứa cả danh mục chung lẫn thực phẩm riêng và công thức của từng không gian, vì người dùng cần gợi ý tên cho cả hai loại; do đó mọi truy vấn bắt buộc kèm bộ lọc “bản chung hoặc thuộc không gian đang thao tác”, tài liệu của nhóm bị xóa khi nhóm giải tán hay chủ tài khoản xóa tài khoản (\nfr{NFR18}). Ngược lại, lô và mặt hàng trong không gian vẫn tìm bằng PostgreSQL vì mỗi không gian chỉ có vài trăm bản ghi đổi liên tục và phải thấy ngay sau khi ghi, còn Elasticsearch chỉ nhất quán sau một khoảng trễ. Chỉ mục là dẫn xuất dựng lại được; khi Elasticsearch lỗi, API chuyển sang truy vấn trigram trên PostgreSQL.

\subsection{Lưu trữ đối tượng MinIO}\label{ssec:cn-minio}

MinIO~\cite{minio} là kho đối tượng mã nguồn mở tương thích API Amazon S3, nên có thể chuyển sang dịch vụ S3 khác mà không sửa mã. Máy chủ cấp URL ký sẵn (presigned URL) có thời hạn, giới hạn kích thước và kiểu nội dung, để thiết bị tải ảnh thẳng lên bucket riêng tư thay vì đi qua API; quy tắc vòng đời (lifecycle) tự xóa tệp báo cáo xuất sau 7 ngày. Ảnh không lưu trong CSDL vì chiếm dung lượng gấp đôi dữ liệu nghiệp vụ (Bảng~\ref{tab:c3-uoc-luong-tai}), sẽ làm phình bản sao lưu, nhật ký WAL và thời gian khôi phục; \nfr{NFR22} cũng yêu cầu dịch vụ lưu tệp riêng.

\subsection{Hàng đợi tác vụ River}\label{ssec:cn-river}

River~\cite{river} là thư viện hàng đợi viết bằng Go, lưu tác vụ trong bảng PostgreSQL và cho phép chèn tác vụ trong cùng giao dịch với dữ liệu nghiệp vụ: commit thì tác vụ chắc chắn tồn tại, rollback thì tác vụ biến mất. Đó là mẫu hộp thư đi giao dịch (transactional outbox)~\cite{richardson} mà không cần bảng outbox và tiến trình chuyển tiếp riêng. River còn có tác vụ duy nhất (unique job) để mỗi người dùng chỉ có một tác vụ nhắc mỗi ngày (\br{BR13}, \nfr{NFR12}), tác vụ định kỳ, thử lại với backoff lũy thừa và bầu chọn tiến trình dẫn đầu (leader election) để chỉ một worker lập lịch tác vụ định kỳ. Hàng đợi trên Redis như Asynq hay BullMQ nhanh, nhưng ghi CSDL và đẩy tác vụ là hai thao tác trên hai hệ thống, không thể nguyên tử: commit xong mà đẩy thất bại thì mất sự kiện, đẩy trước commit thì worker có thể đọc dữ liệu sẽ bị hủy. Kafka hợp với thông lượng rất lớn, nhưng tải giả định chỉ vài chục sự kiện mỗi giây trong khi Kafka đòi vận hành cụm broker và vẫn cần outbox. Đánh đổi của River là thêm tải lên CSDL chính; khi vượt ngưỡng, module notification sẽ được tách ra cùng một message broker (mục~\ref{sec:trien-khai}).

\subsection{Hạ tầng triển khai và vận hành}\label{ssec:cn-ha-tang}

Hệ thống đóng gói bằng Docker, mô tả bằng Docker Compose~\cite{docker} với một tệp khai báo dịch vụ, mạng, volume, kiểm tra sức khỏe và giới hạn tài nguyên. Caddy~\cite{caddy} 2 làm reverse proxy ở biên: tự cấp, gia hạn chứng chỉ Let's Encrypt, chỉ nhận TLS từ 1.2 (\nfr{NFR14}), hỗ trợ HTTP/2, WebSocket và cân bằng tải giữa hai bản API với cấu hình ngắn hơn Nginx. Mã nguồn dự án nằm trên GitLab nên GitLab CI~\cite{gitlabci} đảm nhận kiểm tra, kiểm thử, xây dựng, quét lỗ hổng và triển khai. Prometheus~\cite{prometheus} thu số đo RED (tốc độ, lỗi, thời gian phản hồi theo route) và độ sâu hàng đợi, Grafana hiển thị và cảnh báo, Loki lưu nhật ký do Grafana Alloy thu thập. OpenTelemetry~\cite{opentelemetry} tạo span cho yêu cầu HTTP, truy vấn SQL, lời gọi ra ngoài và gắn mã trace vào nhật ký; vì chuẩn này trung lập nhà cung cấp, thêm kho lưu trace về sau chỉ là đổi cấu hình. pgBackRest~\cite{pgbackrest} sao lưu đầy đủ hằng ngày, giữ 7 bản (\nfr{NFR13}) và lưu liên tục WAL ra bucket ngoài máy chủ để khôi phục về thời điểm bất kỳ.

\subsection{So sánh và tổng hợp lựa chọn công nghệ}\label{ssec:cn-so-sanh}

Bảng~\ref{tab:c3-so-sanh-ngon-ngu} so sánh bốn phương án phía máy chủ theo tám tiêu chí rút từ các kịch bản ở Bảng~\ref{tab:c3-kich-ban-chat-luong}; BaaS (Backend as a Service) gồm Firebase và Supabase.

{\small
\begin{longtable}{|L{2.1cm}|L{2.95cm}|L{2.95cm}|L{2.95cm}|L{\dimexpr\textwidth-10\tabcolsep-6\arrayrulewidth-2.1cm-8.85cm\relax}|}
\caption{So sánh các phương án nền tảng phía máy chủ}\label{tab:c3-so-sanh-ngon-ngu}\\
\hline
\hd{Tiêu chí} & \hd{Go} & \hd{NestJS (TypeScript)} & \hd{Spring Boot (Java)} & \hd{BaaS}\\ \hline
\endfirsthead
\multicolumn{5}{l}{\footnotesize\itshape Bảng \thetable\ (tiếp theo)}\\ \hline
\hd{Tiêu chí} & \hd{Go} & \hd{NestJS (TypeScript)} & \hd{Spring Boot (Java)} & \hd{BaaS}\\ \hline
\endhead
Hiệu năng, tài nguyên & Mã máy, ít bộ nhớ; goroutine rẻ & Tốt cho I/O; việc nặng CPU chặn vòng lặp sự kiện & Thông lượng cao; tốn bộ nhớ, khởi động chậm & Tự co giãn; ít kiểm soát độ trễ\\ \hline
Giao dịch, khóa dòng & SQL tường minh, \texttt{FOR UPDATE} trực tiếp & ORM; khóa dòng cần SQL thô & JPA hoàn thiện; ORM che SQL & Firestore không ràng buộc quan hệ; Supabase qua hàm CSDL\\ \hline
OpenAPI & Sinh giao diện máy chủ từ hợp đồng & Chủ yếu sinh tài liệu từ mã & Sinh khung máy chủ hoàn thiện & Gắn với lược đồ hoặc SDK\\ \hline
Tiến trình nền & River, chèn tác vụ trong giao dịch & BullMQ trên Redis, ghi kép & Quartz, JobRunr; cấu hình nặng & Hàm đám mây; khó chống trùng\\ \hline
Độ phức tạp & Tường minh; nhiều mã lặp & Trung bình & Khung lớn, nhiều cơ chế ngầm & Thấp lúc đầu; khó diễn đạt FEFO, \br{BR20}\\ \hline
Chi phí hạ tầng & Vừa một máy chủ nhỏ & Vừa một máy chủ nhỏ & Cần nhiều bộ nhớ hơn & Theo lượt đọc ghi, khó dự báo\\ \hline
Phụ thuộc nhà cung cấp & \multicolumn{3}{L{\dimexpr8.85cm+4\tabcolsep+2\arrayrulewidth\relax}|}{Thấp: mã nguồn mở, tự vận hành bằng Docker} & Firebase cao; Supabase vừa\\ \hline
Kiểm thử & Gói \texttt{testing} chuẩn, phát hiện tranh chấp & Jest, dễ giả lập & JUnit, Testcontainers; khởi tạo chậm & Phụ thuộc bộ giả lập\\ \hline
\end{longtable}
}

Nhóm chọn Go vì ưu điểm của nó khớp các kịch bản rủi ro cao nhất: giao dịch và khóa dòng tường minh qua pgx, sqlc cho KB05; goroutine giữ nhiều kết nối với chi phí thấp cho KB03; River gắn tác vụ nền nguyên tử với giao dịch cho KB04, KB06; tệp thực thi nhỏ giúp cả hệ thống vừa một máy chủ. Spring Boot mạnh về giao dịch nhưng nặng hơn mức cần, hàng đợi phổ biến của NestJS vướng ghi kép, còn BaaS đẩy FEFO, \br{BR20} và nhắc nhở chống trùng vào các hàm đám mây rời rạc, kèm phụ thuộc nhà cung cấp. Cái giá của Go là mã dài hơn và ít tiện ích dựng sẵn, được bù bằng sinh mã và quy ước module thống nhất (mục~\ref{sec:thiet-ke-chi-tiet}).

Lưu trữ đa CSDL (polyglot persistence) ở đây theo phân biệt của Kleppmann giữa hệ thống ghi nhận (system of record) và dữ liệu dẫn xuất~\cite{kleppmann}: PostgreSQL là hệ thống ghi nhận duy nhất, mỗi kho khác chỉ làm việc PostgreSQL làm kém, gồm Redis cho bộ đếm và phát tán sự kiện, Elasticsearch cho tìm kiếm có độ liên quan, MinIO cho tệp nhị phân lớn; phân bổ dữ liệu chi tiết nằm ở mục~\ref{sec:csdl}. Vì mỗi kho thêm vào kéo theo chi phí vận hành, nhóm không dùng MongoDB (dữ liệu có tính quan hệ, cần giao dịch nhiều thực thể, phần bán cấu trúc đã có JSONB), ClickHouse (báo cáo chỉ quét một không gian trong tối đa 12 tháng theo \br{BR21}, view vật chất hóa là đủ), Kafka (mục~\ref{ssec:cn-river}) và Kubernetes (khoảng mười lăm container trên một máy chủ được Docker Compose quản lý đủ; xem lộ trình ở mục~\ref{sec:trien-khai}). Bảng~\ref{tab:c3-tong-hop-cong-nghe} tổng hợp các lựa chọn.

{\small
\begin{longtable}{|L{2.5cm}|L{4.4cm}|L{\dimexpr\textwidth-8\tabcolsep-5\arrayrulewidth-2.5cm-4.4cm-3.1cm\relax}|L{3.1cm}|}
\caption{Tổng hợp công nghệ sử dụng}\label{tab:c3-tong-hop-cong-nghe}\\
\hline
\hd{Thành phần} & \hd{Công nghệ, phiên bản} & \hd{Vai trò} & \hd{NFR đáp ứng}\\ \hline
\endfirsthead
\multicolumn{4}{l}{\footnotesize\itshape Bảng \thetable\ (tiếp theo)}\\ \hline
\hd{Thành phần} & \hd{Công nghệ, phiên bản} & \hd{Vai trò} & \hd{NFR đáp ứng}\\ \hline
\endhead
Thiết bị & Flutter 3, drift (SQLite), flutter\_secure\_storage & Đa nền tảng, ngoại tuyến & \nfr{NFR11}, \nfr{NFR15}, \nfr{NFR19}, \nfr{NFR23}\\ \hline
Thông báo đẩy & FCM HTTP v1, APNs & Nhắc nhở, phân công & \nfr{NFR10}, \nfr{NFR12}, \nfr{NFR22}\\ \hline
Máy chủ & Go 1.23, chi v5, coder/websocket & API, WebSocket, worker & \nfr{NFR09}, \nfr{NFR10}, \nfr{NFR21}\\ \hline
Thư viện Go & pgx v5, sqlc, goose, oapi-codegen & Giao dịch, di trú, sinh mã & \nfr{NFR13}, \nfr{NFR20}, \nfr{NFR21}\\ \hline
Hàng đợi & River & Outbox, định kỳ, thử lại & \nfr{NFR12}, \nfr{NFR22}\\ \hline
CSDL chính & PostgreSQL 16 (citext, pg\_trgm, unaccent) & Nguồn sự thật, ACID & \nfr{NFR13}, \nfr{NFR16}, \nfr{NFR22}\\ \hline
Bộ nhớ đệm & Redis 7 & Cache, giới hạn, Pub/Sub & \nfr{NFR09}, \nfr{NFR10}, \nfr{NFR17}\\ \hline
Tìm kiếm & Elasticsearch 8.15 & Thực phẩm, công thức (chung và riêng) & \nfr{NFR09}, \nfr{NFR18}\\ \hline
Kho đối tượng & MinIO (API S3) & Ảnh, tệp xuất & \nfr{NFR22}\\ \hline
Dịch vụ ngoài & SMTP relay & Thư OTP & \br{BR02}, \nfr{NFR09}\\ \hline
Biên, triển khai & Caddy 2, Docker Compose, GitLab CI & TLS, cân bằng tải, CI/CD & \nfr{NFR13}, \nfr{NFR14}, \nfr{NFR20}\\ \hline
Vận hành & Prometheus, Grafana, Loki, OpenTelemetry, pgBackRest & Quan sát, sao lưu & \nfr{NFR09}, \nfr{NFR12}, \nfr{NFR13}\\ \hline
\end{longtable}
}

Như vậy, mỗi ràng buộc \nfr{NFR21}--\nfr{NFR23} đều có thành phần đáp ứng trong bảng; mục~\ref{sec:kien-truc} dùng bộ công nghệ này để dựng kiến trúc tổng thể.

% ---------------------------------------------------------------------
%  3.3. Thiết kế kiến trúc
% ---------------------------------------------------------------------
\section{Thiết kế kiến trúc}\label{sec:kien-truc}

Mục này chuyển các yếu tố dẫn dắt ở mục~\ref{ssec:yeu-to-dan-dat} và bộ công nghệ ở mục~\ref{sec:cong-nghe} thành cấu trúc phần mềm, trình bày theo các mức của mô hình C4~\cite{c4model} và kết thúc bằng các quyết định kiến trúc kèm đánh đổi.

\subsection{Lựa chọn kiến trúc phần mềm}\label{ssec:lua-chon-kien-truc}

Bốn đặc điểm chi phối lựa chọn: quy mô nhỏ (theo giả định thiết kế ở mục~\ref{ssec:uoc-luong-tai}, đỉnh khoảng 50 yêu cầu mỗi giây, dữ liệu dưới 2 GB mỗi năm); thao tác phải nguyên tử trên nhiều nghiệp vụ, như nhập mặt hàng đã mua vào tủ (\uc{UC11}) hay xác nhận đã nấu kèm trừ kho FEFO (\uc{UC19}, \br{BR15}, \nfr{NFR13}); nhóm chỉ có bốn người; và mã phải phân tầng, dễ bảo trì để lớn lên khi người dùng vượt giả định (\nfr{NFR20}). Bảng~\ref{tab:c3-so-sanh-kien-truc} so sánh bốn phương án: nguyên khối phân tầng truyền thống (layered monolith), khối nguyên khối chia module (modular monolith), vi dịch vụ (microservices) và nền tảng dịch vụ backend kèm hàm phi máy chủ (Backend-as-a-Service, serverless) như Firebase, Supabase.

{\footnotesize\setlength{\tabcolsep}{4pt}
\begin{longtable}{|L{2.5cm}|L{\dimexpr(\textwidth-10\tabcolsep-6\arrayrulewidth-2.5cm)/4\relax}|L{\dimexpr(\textwidth-10\tabcolsep-6\arrayrulewidth-2.5cm)/4\relax}|L{\dimexpr(\textwidth-10\tabcolsep-6\arrayrulewidth-2.5cm)/4\relax}|L{\dimexpr(\textwidth-10\tabcolsep-6\arrayrulewidth-2.5cm)/4\relax}|}
\caption{So sánh các phương án kiến trúc phần mềm}\label{tab:c3-so-sanh-kien-truc}\\
\hline
\hd{Tiêu chí} & \hd{Nguyên khối phân tầng} & \hd{Nguyên khối chia module} & \hd{Vi dịch vụ} & \hd{BaaS, serverless}\\ \hline
\endfirsthead
\multicolumn{5}{l}{\footnotesize\itshape Bảng \thetable\ (tiếp theo)}\\ \hline
\hd{Tiêu chí} & \hd{Nguyên khối phân tầng} & \hd{Nguyên khối chia module} & \hd{Vi dịch vụ} & \hd{BaaS, serverless}\\ \hline
\endhead
Độ phức tạp vận hành & Thấp & Thấp: một image, hai chế độ chạy & Cao: broker, trace phân tán & Thấp, phụ thuộc nhà cung cấp\\ \hline
Giao dịch \uc{UC11}, \uc{UC19} & ACID cục bộ & ACID cục bộ, chung giao dịch & Saga, giao dịch bù & Dồn vào hàm SQL, hàm đám mây\\ \hline
Chi phí hạ tầng & Một máy chủ & Một máy chủ & Nhiều nút, broker & Rẻ lúc nhỏ, khó dự báo\\ \hline
Tốc độ phát triển (4 người) & Nhanh, chậm dần khi mã rối & Nhanh, dễ phân công & Chậm vì hạ tầng, hợp đồng & Nhanh với CRUD, chậm với FEFO, gợi ý\\ \hline
Khả năng mở rộng & Nhân bản cả khối & Cả khối; tách \texttt{worker}, tách module & Độc lập từng dịch vụ & Tự động, trong giới hạn nền tảng\\ \hline
Ranh giới, bảo trì & Dễ xói mòn & Rõ nhờ gói \texttt{api}, schema, linter & Rất rõ & Logic rải ở client và hàm\\ \hline
Kiểm thử & Dễ & Dễ, theo module & Khó, cần kiểm thử hợp đồng & Khó mô phỏng cục bộ\\ \hline
Hợp quy mô giả định & Hợp, khó tiến hóa & Hợp, có lộ trình tách & Thừa & Hợp tải, khó áp \nfr{NFR11}\\ \hline
\hs{Kết luận} & Loại & \textbf{Chọn} & Loại & Loại\\ \hline
\end{longtable}
}

Nhóm chọn khối nguyên khối chia module, bên trong mỗi module dùng kiến trúc lục giác. Mỗi module là một ngữ cảnh giới hạn (bounded context)~\cite{evans} có mô hình miền và schema riêng, chỉ lộ gói \texttt{api} làm cổng công khai; bên trong, cổng và bộ chuyển đổi (ports and adapters)~\cite{hexagonal} cô lập lõi nghiệp vụ khỏi HTTP, CSDL và hàng đợi (mục~\ref{sec:thiet-ke-chi-tiet}). Lựa chọn này theo nguyên tắc “nguyên khối trước”~\cite{fowlermonolith}: tách dịch vụ khi ranh giới nghiệp vụ còn biến động dễ khóa chặt ranh giới sai. Lý do quyết định là tính nguyên tử: vi dịch vụ buộc \uc{UC11}, \uc{UC19} dùng saga và giao dịch bù~\cite{richardson}, ở đây chỉ cần một giao dịch PostgreSQL. BaaS bị loại vì logic FEFO, gợi ý món phải viết thành hàm SQL hoặc hàm đám mây, và cơ chế đồng bộ có sẵn áp “ghi sau thắng” cho mọi trường, trái với \nfr{NFR11}.

Cái giá phải trả là cả khối được nhân bản cùng nhau, các module tranh chấp một CSDL và ranh giới dễ xói mòn nếu chỉ là quy ước. Nhóm giảm thiểu bằng kiểm tra phụ thuộc tự động trong CI (mục~\ref{ssec:tong-quan-module}), tách tác vụ nền sang tiến trình \texttt{worker} và cho mọi tác dụng phụ đi qua sự kiện miền. Nhờ đó lộ trình tách dịch vụ đã sẵn: ứng viên đầu tiên là \texttt{notification}, vốn chỉ nhận sự kiện; khi tách chỉ cần thay bảng tác vụ River bằng message broker (mục~\ref{sec:trien-khai}).

\subsection{Kiến trúc tổng thể theo mô hình C4}\label{ssec:c4-tong-the}

Mô hình C4~\cite{c4model} có bốn mức: ngữ cảnh, vùng chứa (container), thành phần (component) và mã nguồn (mục~\ref{sec:thiet-ke-chi-tiet}). Hình~\ref{fig:c4-ngu-canh} là sơ đồ ngữ cảnh.

\diagram[\textwidth]{c3-01-c4-ngu-canh}{Sơ đồ ngữ cảnh hệ thống theo mô hình C4 (mức 1)}{fig:c4-ngu-canh}

Hệ thống gồm ứng dụng di động và máy chủ, phục vụ người dùng trên Android, iOS (\nfr{NFR19}) và quản trị viên qua khu vực quản trị trong cùng ứng dụng. FCM~\cite{fcm} là API đẩy chung, gửi thẳng tới Android và chuyển tiếp qua APNs~\cite{apns} tới iOS. Bộ lập lịch của chương 2 nay được hiện thực bên trong hệ thống bằng River. Hình~\ref{fig:c4-container} và Bảng~\ref{tab:c3-container} trình bày các vùng chứa; bốn kho dữ liệu phía máy chủ được gom vào ranh giới “Lưu trữ” vì chúng cùng nằm trong mạng Docker nội bộ \texttt{data}, còn các hệ thống bên ngoài xếp thành cột bên phải.

\diagram[\textwidth]{c3-02-c4-container}{Sơ đồ vùng chứa theo mô hình C4 (mức 2)}{fig:c4-container}

{\footnotesize\setlength{\tabcolsep}{4pt}
\begin{longtable}{|L{2.3cm}|L{3.0cm}|L{\dimexpr\textwidth-8\tabcolsep-5\arrayrulewidth-2.3cm-3.0cm-4.0cm\relax}|L{4.0cm}|}
\caption{Các vùng chứa của hệ thống}\label{tab:c3-container}\\
\hline
\hd{Vùng chứa} & \hd{Công nghệ} & \hd{Trách nhiệm} & \hd{Giao tiếp chính}\\ \hline
\endfirsthead
\multicolumn{4}{l}{\footnotesize\itshape Bảng \thetable\ (tiếp theo)}\\ \hline
\hd{Vùng chứa} & \hd{Công nghệ} & \hd{Trách nhiệm} & \hd{Giao tiếp chính}\\ \hline
\endhead
Ứng dụng di động & Flutter 3, Riverpod, dio & Giao diện, hàng đợi lệnh ngoại tuyến & HTTPS, WSS tới Caddy; PUT ảnh lên URL ký sẵn\\ \hline
CSDL cục bộ & SQLite qua drift & Bản sao dữ liệu không gian, lệnh chờ gửi (\nfr{NFR23}) & Trong tiến trình ứng dụng\\ \hline
Reverse proxy & Caddy 2~\cite{caddy} & TLS $\ge$ 1.2, HTTP/2, chia tải tới các bản API; phục vụ \path{/.well-known} cho liên kết mời & HTTP tới API; \texttt{media.} tới MinIO\\ \hline
API (×2 bản) & Go 1.23, chi v5, pgx, sqlc & REST \path{/it4788/v1}, WebSocket, phân quyền, nghiệp vụ, chèn tác vụ & PostgreSQL, Redis, ES, MinIO\\ \hline
Worker & Go 1.23, River~\cite{river} & Tác vụ nền và định kỳ: nhắc, đẩy, thư, chỉ mục, ảnh, xuất tệp & PostgreSQL, Redis, FCM, SMTP, ES, MinIO\\ \hline
CSDL chính & PostgreSQL 16~\cite{postgresql} & Nguồn sự thật, nhật ký thay đổi, bảng River & Chỉ trong mạng nội bộ\\ \hline
Bộ nhớ đệm & Redis 7~\cite{redis} & Cache, giới hạn tần suất, Pub/Sub & Dữ liệu tạm, mất không ảnh hưởng tính đúng\\ \hline
Máy tìm kiếm & Elasticsearch 8.15~\cite{elasticsearch} & Chỉ mục \texttt{foods}, \texttt{recipes} (bản chung và bản riêng của không gian) & Dựng lại được từ PostgreSQL; truy vấn luôn lọc theo không gian\\ \hline
Lưu trữ đối tượng & MinIO~\cite{minio} & Bucket \texttt{media}, \texttt{exports} & S3 API; bên ngoài chỉ qua URL ký sẵn\\ \hline
\end{longtable}
}

Luồng đồng bộ đi qua Caddy tới một trong hai bản API không giữ trạng thái phiên~\cite{twelvefactor}, nên được chia tải tùy ý và cập nhật cuốn chiếu (\nfr{NFR13}). Luồng nền bắt đầu khi API chèn tác vụ vào \texttt{river\_job} trong giao dịch nghiệp vụ, nên độ trễ của FCM, SMTP không cộng vào thời gian phản hồi (\nfr{NFR09}). Luồng thời gian thực dùng Redis Pub/Sub; ảnh được tải thẳng lên MinIO qua URL ký sẵn.

\subsection{Giao tiếp giữa các thành phần}\label{ssec:giao-tiep}

Giao tiếp đồng bộ gồm REST/JSON qua HTTPS~\cite{rfc9110} giữa ứng dụng và API theo hợp đồng OpenAPI (mục~\ref{sec:api}), và lời gọi hàm trong tiến trình giữa các module qua interface của gói \texttt{<module>/api}. \texttt{TxManager} đặt \texttt{pgx.Tx} vào \texttt{context}, nên khi \texttt{shopping} gọi \texttt{pantry} để nhập tủ (\uc{UC11}) hay \texttt{mealplan} gọi \texttt{pantry} để trừ kho (\uc{UC19}), mọi câu lệnh nằm trong một giao dịch.

Giao tiếp bất đồng bộ dựa trên sự kiện miền: module phát chèn một tác vụ River trong chính giao dịch nghiệp vụ, theo mẫu hộp thư đi trong giao dịch (transactional outbox)~\cite{richardson}. Mẫu này tránh lỗi ghi kép (dual write): phát tới broker sau commit có thể mất sự kiện, phát trước commit có thể sinh sự kiện “ma” khi rollback; tác vụ River là một dòng PostgreSQL nên tồn tại khi và chỉ khi dữ liệu được commit. Kênh thời gian thực chỉ báo có thay đổi: ở cuối giao dịch, cổng \texttt{ChangefeedAppender} (hiện thực bởi \texttt{SyncEngine}) tăng \texttt{change\_seq} và ghi \texttt{sync.space\_changes}; sau commit, API phát \texttt{\{seq, entityTypes\}} lên kênh \texttt{space:\{spaceId\}}, mọi bản API giữ WebSocket~\cite{rfc6455} của thành viên chuyển gợi ý xuống, thiết bị gọi \texttt{GET /changes?since=} tới một bản API bất kỳ để kéo phần chênh lệch (\nfr{NFR10}); khi ứng dụng chạy nền, iOS ngắt socket nên hệ thống dùng FCM. Hình~\ref{fig:giao-tiep-su-kien} thể hiện hai nhánh xử lý sau một lệnh ghi.

\diagram[\textwidth]{c3-03-giao-tiep-su-kien}{Luồng sự kiện miền: hộp thư đi trong giao dịch và phát tán thời gian thực}{fig:giao-tiep-su-kien}

Các bước 1–3 cùng một giao dịch: câu lệnh đầu tiên chèn khóa lũy đẳng, câu lệnh cuối cùng tăng \texttt{change\_seq} rồi \texttt{COMMIT}, nên dữ liệu, nhật ký thay đổi và tác vụ luôn cùng tồn tại hoặc cùng biến mất. Bước 4a chỉ chạy sau \texttt{COMMIT}, vì phát trước commit có thể khiến thiết bị kéo về một \texttt{seq} chưa nhìn thấy được. Nhánh “a” đi qua Redis Pub/Sub, vốn giao tối đa một lần; gợi ý bị mất chỉ làm thiết bị cập nhật chậm hơn vì thiết bị luôn kéo theo \texttt{seq} khi kết nối lại. Nhánh “b” do worker xử lý; ngoài việc gửi đẩy và thư, \texttt{NotifyWorker} phát \texttt{user:\{id\}} sau khi tạo thông báo để hộp thư trong ứng dụng của đúng người nhận cập nhật ngay, và cùng kênh này được dùng để đóng socket khi người dùng bị xóa khỏi nhóm hoặc bị khóa. Bảng~\ref{tab:c3-su-kien-mien} liệt kê các sự kiện miền và tác vụ nền tương ứng.

{\footnotesize\setlength{\tabcolsep}{4pt}
\begin{longtable}{|L{3.9cm}|L{1.7cm}|L{\dimexpr\textwidth-10\tabcolsep-6\arrayrulewidth-3.9cm-1.7cm-1.8cm-2.9cm\relax}|L{1.8cm}|L{2.9cm}|}
\caption{Các sự kiện miền, tác vụ nền và bảo đảm giao nhận}\label{tab:c3-su-kien-mien}\\
\hline
\hd{Sự kiện, tác vụ} & \hd{Module phát} & \hd{Bên xử lý (tác vụ River)} & \hd{Hàng đợi, kênh} & \hd{Bảo đảm}\\ \hline
\endfirsthead
\multicolumn{5}{l}{\footnotesize\itshape Bảng \thetable\ (tiếp theo)}\\ \hline
\hd{Sự kiện, tác vụ} & \hd{Module phát} & \hd{Bên xử lý (tác vụ River)} & \hd{Hàng đợi, kênh} & \hd{Bảo đảm}\\ \hline
\endhead
\texttt{user.registered}, \texttt{password.\allowbreak reset\_requested}; thư báo khóa tài khoản & identity, admin & notification: \texttt{SendEmail} (thư OTP, thư tài khoản) & \texttt{notify} & Outbox, ít nhất một lần\\ \hline
\texttt{member.joined}, \texttt{member.left}, \texttt{member.removed}, \texttt{leader.transferred} & household & notification: hoạt động nhóm; sync & \texttt{notify}; nhật ký & Outbox; nhật ký nguyên tử\\ \hline
Tác vụ \texttt{AssignmentDigest} (S10) & shopping & Chính shopping: sau 5 s đếm $k$ mặt hàng đang giao rồi phát \texttt{item.assigned} ($k$, mốc đợt) qua \texttt{PushNotification} & \texttt{notify} & Duy nhất theo người được giao và danh sách\\ \hline
\texttt{item.assigned}, \texttt{item.unassigned} & shopping & notification: đẩy cho người được giao, đọc \texttt{assignment\_\allowbreak notify\_on}; số mặt hàng lấy từ payload & \texttt{notify} & \texttt{dedupe\_key} gồm mốc đợt\\ \hline
\texttt{items.purchased} & shopping & sync; report: tổng hợp ngày & Nhật ký; \texttt{default} & Nguyên tử; ít nhất một lần\\ \hline
\texttt{items.stocked}, \texttt{lot.changed}, \texttt{meal.cooked} & shopping, pantry, mealplan & sync; xóa \texttt{suggest:\{sid\}} khi lô đổi & Nhật ký; \texttt{space:\allowbreak\{id\}} & Nguyên tử; gợi ý tối đa một lần\\ \hline
\texttt{food.changed}, \texttt{recipe.changed} & catalog, recipe & search: đánh chỉ mục; \texttt{recipe.changed} còn xóa \texttt{suggest:\{sid\}} & \texttt{index} & Ít nhất một lần, ghi đè theo phiên bản\\ \hline
\texttt{food.merged} (\uc{UC23}) & catalog & \texttt{RemapFood} cho từng module pantry, shopping, recipe; search cập nhật tài liệu & \texttt{default}, \texttt{index} & Lũy đẳng: đổi \texttt{food\_id} cũ sang đích\\ \hline
\texttt{config.changed} & admin & Mọi bản API và worker: xóa \texttt{cfg:all}; khóa \texttt{suggest.*} đổi thì xóa \texttt{suggest:*} & Pub/Sub \texttt{cfg} & Tối đa một lần, bù bằng TTL 5 phút\\ \hline
Tác vụ \texttt{CopySpaceData} (G02), \texttt{PurgeSpace} (giải tán nhóm) & household & Sao chép dữ liệu cá nhân sang nhóm mới; xóa mềm dữ liệu chung, xóa tài liệu ES & \texttt{default} & Lũy đẳng theo không gian\\ \hline
\texttt{media.uploaded}, \texttt{export.requested} & media, report & MediaWorker, ExportWorker & \texttt{media}, \texttt{report} & Outbox, ít nhất một lần\\ \hline
\texttt{export.completed} & report & notification: báo tệp xuất đã sẵn sàng & \texttt{notify} & Ít nhất một lần\\ \hline
Định kỳ mỗi phút: \texttt{ReminderDispatcher} & Bộ lập lịch & Chèn \texttt{UserReminder} cho từng người có lô trong cửa sổ nhắc (\uc{UC14}) & \texttt{default}, \texttt{notify} & Duy nhất theo ngày; theo người và ngày\\ \hline
\end{longtable}
}

River giao tác vụ theo ngữ nghĩa ít nhất một lần (at-least-once): worker chết sau khi gửi thư nhưng trước khi đánh dấu hoàn tất thì tác vụ chạy lại. Vì vậy mọi bên nhận phải lũy đẳng: thông báo có khóa \texttt{dedupe\_key} duy nhất (\br{BR13}, \nfr{NFR12}); chỉ mục ghi đè tài liệu theo \texttt{id} và \texttt{version}; \texttt{RemapFood} chỉ đổi những dòng còn trỏ tới thực phẩm đã gộp; ảnh thu nhỏ, tệp xuất có khóa đối tượng xác định trước; tác vụ \texttt{PushNotification} mang khóa gộp (collapse key). Thứ tự tác vụ không được bảo đảm, nên bên xử lý đọc trạng thái hiện tại từ CSDL của chính module sở hữu dữ liệu thay vì tin nội dung sự kiện cũ. \texttt{AssignmentDigest} là ví dụ: tác vụ thuộc \texttt{shopping} đếm lại các mặt hàng đang giao ngay lúc chạy rồi mới đóng số $k$ vào payload, nhờ đó \texttt{notification} soạn được nội dung “Bạn được giao mua $k$ mặt hàng” mà không phải đọc bảng của \texttt{shopping}.

\subsection{Thiết kế tổng quan các module}\label{ssec:tong-quan-module}

Backend gồm mười ba module trong \path{internal/modules}, gói hạ tầng \texttt{platform} và nhân \texttt{kernel}. Mỗi module sở hữu một schema PostgreSQL trùng tên gói, riêng \texttt{search} không có bảng. Bảng~\ref{tab:c3-module} nêu trách nhiệm, use case và phụ thuộc biên dịch trực tiếp của từng gói, là căn cứ cấu hình luật kiểm tra kiến trúc; tên cổng ghi trong ngoặc là interface công bố trong gói \texttt{<module>/api} của module được dùng.

{\footnotesize\setlength{\tabcolsep}{4pt}
\begin{longtable}{|L{2.4cm}|L{\dimexpr\textwidth-8\tabcolsep-5\arrayrulewidth-2.4cm-1.9cm-4.5cm\relax}|L{1.9cm}|L{4.5cm}|}
\caption{Các module của backend}\label{tab:c3-module}\\
\hline
\hd{Gói, schema} & \hd{Trách nhiệm} & \hd{Use case} & \hd{Phụ thuộc biên dịch}\\ \hline
\endfirsthead
\multicolumn{4}{l}{\footnotesize\itshape Bảng \thetable\ (tiếp theo)}\\ \hline
\hd{Gói, schema} & \hd{Trách nhiệm} & \hd{Use case} & \hd{Phụ thuộc biên dịch}\\ \hline
\endhead
\texttt{identity} & Đăng ký, OTP, phiên, hồ sơ, thiết lập, thiết bị; định nghĩa \texttt{ActivationHook} & \uc{UC01}--\uc{UC04} & media\\ \hline
\texttt{household} & Không gian, thành viên, mã mời, chuyển quyền, giải tán; định nghĩa \texttt{MemberLeaveHook} & \uc{UC05}--\uc{UC07} & identity (hiện thực \texttt{identity/api.\allowbreak ActivationHook}), sync, media\\ \hline
\texttt{catalog} & Thực phẩm, loại, đơn vị, hạn dùng gợi ý, duyệt và gộp thực phẩm riêng & \uc{UC23} & household, media\\ \hline
\texttt{shopping} & Danh sách, mặt hàng, phân công, đã mua; công bố \texttt{ListComposer}; tác vụ \texttt{AssignmentDigest} & \uc{UC08}, \uc{UC09}, \uc{UC11}, \uc{UC10} & household (hiện thực \texttt{household/api.\allowbreak MemberLeaveHook}, \br{BR07}), catalog, pantry (\texttt{StockIn}), sync\\ \hline
\texttt{pantry} & Lô, biến động kho, FEFO; công bố \texttt{StockQuery}, \texttt{StockIn}, \texttt{StockConsumer} & \uc{UC12}, \uc{UC11}, \uc{UC13} & household, catalog, sync\\ \hline
\texttt{recipe} & Công thức của không gian và công thức mẫu, yêu thích; công bố \texttt{RecipeQuery}, định nghĩa \texttt{RecipeDeleteHook} & \uc{UC15}, \uc{UC16}, \uc{UC10}, \uc{UC23} & household, catalog, pantry (\texttt{StockQuery}: \uc{UC16}, \uc{UC10}), shopping (\texttt{ListComposer}), sync, media\\ \hline
\texttt{mealplan} & Lịch bữa ăn, gợi ý món, xác nhận đã nấu; điều phối sinh danh sách & \uc{UC17}--\uc{UC19}, điều phối \uc{UC10} & household, catalog, pantry (\texttt{StockQuery}, \texttt{StockConsumer}), recipe (\texttt{RecipeQuery}; hiện thực \texttt{recipe/api.\allowbreak RecipeDeleteHook}, \uc{UC15}), shopping (\texttt{ListComposer}), sync\\ \hline
\texttt{search} & Chỉ mục, truy vấn ES lọc theo không gian, dự phòng trigram & \uc{UC20} & catalog, recipe, pantry (\texttt{StockQuery.\allowbreak SearchLots} cho Q01)\\ \hline
\texttt{report} & Mô hình đọc: báo cáo qua view, xuất tệp; sở hữu bảng \texttt{report.exports} & \uc{UC21}, \uc{UC22} & Không; \texttt{SELECT} qua view \texttt{report.v\_*} trên shopping, pantry, mealplan, catalog\\ \hline
\texttt{notification} & Hộp thư, nhắc hạn dùng, đẩy, thư & \uc{UC14} & identity, household, pantry; số mặt hàng phân công nhận qua payload, không đọc bảng shopping\\ \hline
\texttt{sync} & \texttt{SyncEngine}: nhật ký thay đổi (hiện thực \texttt{ChangefeedAppender}), khóa lũy đẳng; hub WebSocket & \nfr{NFR10}, \nfr{NFR11} & Không\\ \hline
\texttt{admin} & Tài khoản, cấu hình, nhật ký quản trị, cổng quản trị danh mục & \uc{UC24}, \uc{UC25} & identity, catalog, recipe\\ \hline
\texttt{media} & URL ký sẵn, xác nhận tải lên, xử lý ảnh & Dùng chung & Không\\ \hline
\texttt{platform} & db, \texttt{TxManager}, httpx, auth, events, idgen, i18n, observability\ldots & — & Không\\ \hline
\texttt{kernel} & \texttt{Quantity} (\br{BR16}), \texttt{Money}, \texttt{DateRange}, kiểu ID & — & Chỉ thư viện chuẩn\\ \hline
\end{longtable}
}

Hình~\ref{fig:phu-thuoc-module} biểu diễn các phụ thuộc bằng sơ đồ gói UML~\cite{uml251}, gom module thành bốn tầng. Sơ đồ dùng bốn loại quan hệ: nét liền «use» là phụ thuộc biên dịch qua gói \texttt{<module>/api}; nét đứt tam giác rỗng «hook» là hiện thực giao diện móc; nét đứt «flow» là luồng thông tin (information flow) của UML 2.5 dành cho sự kiện miền đi qua River, không phải phụ thuộc; nét đứt «read» là quyền đọc qua view của mô hình đọc \texttt{report}. Mũi tên giữa hai khung nghĩa là một số module khung trên dùng một số module khung dưới.

Phụ thuộc biên dịch không có chu trình và luôn hướng xuống. Trong tầng B, mỗi mũi tên ghi tên cổng được dùng: \texttt{mealplan} dùng \texttt{StockQuery}, \texttt{StockConsumer} (\uc{UC18}, \uc{UC19}), \texttt{RecipeQuery} và \texttt{ListComposer} (\uc{UC10}, nên API S15 do handler của \texttt{mealplan} xử lý); \texttt{recipe} dùng \texttt{ListComposer} để đưa phần thiếu vào danh sách và \texttt{StockQuery} để đối chiếu nguyên liệu với tủ lạnh (\uc{UC16}, \uc{UC10}); \texttt{shopping} dùng \texttt{StockIn} (\uc{UC11}). Từ tầng A, \texttt{search} dùng \texttt{StockQuery.SearchLots} để tìm lô trong không gian cho Q01, vì dữ liệu lô ở PostgreSQL chứ không ở Elasticsearch. Khi module được dùng cần kích hoạt hành vi ở module dùng nó, nhóm áp dụng nguyên lý đảo phụ thuộc: module được dùng định nghĩa giao diện móc (hook) trong gói \texttt{api} của mình, module kia hiện thực và đăng ký tại \texttt{cmd/api/main.go}. Ba móc như vậy là \texttt{identity/api.ActivationHook} (\texttt{OnActivated}, hiện thực bởi \texttt{household} để tạo không gian cá nhân khi A02 kích hoạt tài khoản), \texttt{household/api.MemberLeaveHook} (\texttt{OnMemberLeave}, hiện thực bởi \texttt{shopping} để hủy phân công của người rời nhóm, \br{BR07}) và \texttt{recipe/api.RecipeDeleteHook} (\texttt{OnRecipeDelete}, hiện thực bởi \texttt{mealplan} để xử lý món đã lên lịch, \uc{UC15}); cả ba chạy trong cùng giao dịch với thao tác gốc. Tương tự, \uc{UC12} do ứng dụng ghép từ P07 và M05 thay vì để \texttt{pantry} gọi \texttt{mealplan}.

\diagram[\textwidth]{c3-04-phu-thuoc-module}{Sơ đồ gói thể hiện phụ thuộc giữa các module}{fig:phu-thuoc-module}

Hình~\ref{fig:phu-thuoc-module} cũng cho thấy tầng B, C không nhập gói của tầng A: thông tin đi lên tầng A chỉ bằng «flow», tức sự kiện miền có kiểu định nghĩa trong \texttt{platform/events}, được chèn thành tác vụ River trong giao dịch nghiệp vụ. Trong tầng A cũng vậy: \texttt{report} báo cho \texttt{notification} bằng sự kiện \texttt{export.completed} chứ không gọi hàm. Sự kiện \texttt{food.merged} của \texttt{catalog} đi tới \texttt{pantry}, \texttt{shopping}, \texttt{recipe} theo cùng cơ chế nên không tạo phụ thuộc ngược từ tầng C lên tầng B. \texttt{sync} được gọi qua \texttt{ChangefeedAppender} nhưng không nhập ngược module nghiệp vụ. Ngoại lệ có chủ đích duy nhất về dữ liệu là \texttt{report}: là mô hình đọc (read model), module này không nhập gói nào của module khác nhưng được cấp quyền \texttt{SELECT} trên các view \texttt{report.v\_*} đặt trên bảng của \texttt{shopping}, \texttt{pantry}, \texttt{mealplan} và \texttt{catalog}; bảng duy nhất nó sở hữu là \texttt{report.exports}. Các view này là hợp đồng đọc có phiên bản, nên module nguồn đổi cấu trúc bảng mà giữ nguyên view thì báo cáo không bị ảnh hưởng.

Ranh giới được cưỡng chế ở ba lớp. Về ngôn ngữ, Go cấm chu trình nhập giữa các gói và cấm mã ngoài kho nhập gói \texttt{internal}. Về công cụ, luật \texttt{depguard} của \texttt{golangci-lint} ở bước verify trong CI chỉ cho phép nhập gói \texttt{api} của module khác theo đúng Bảng~\ref{tab:c3-module} và cấm tầng B, C nhập tầng A. Về dữ liệu, mỗi module chỉ truy vấn schema của mình; khóa ngoại liên schema chỉ trỏ tới \texttt{identity.users}, \texttt{household.spaces} và \texttt{catalog} (mục~\ref{sec:csdl}), nên khi tách dịch vụ chỉ cần chuyển chúng thành kiểm tra ở tầng ứng dụng.

\subsection{Thiết kế thành phần (C4 mức 3)}\label{ssec:thiet-ke-thanh-phan}

Vùng chứa API có quá nhiều thành phần để đọc được trên một trang, nên được trình bày thành hai hình: phía vào (nhận và kiểm soát yêu cầu) và phía ra (dịch vụ ứng dụng cùng các adapter tới kho dữ liệu, dịch vụ ngoài). Ở cả hai hình, handler của mọi module được gom thành một thành phần, dịch vụ ứng dụng của các module cũng vậy, vì cấu trúc bên trong chúng giống nhau và được trình bày ở mức lớp tại mục~\ref{sec:thiet-ke-chi-tiet}. Hình~\ref{fig:c4-thanh-phan-api-vao} là phía vào.

\diagram[\textwidth]{c3-05a-c4-thanh-phan-api-vao}{Sơ đồ thành phần của vùng chứa API theo mô hình C4 (mức 3): phía vào}{fig:c4-thanh-phan-api-vao}

Yêu cầu đi qua router chi và chuỗi middleware \texttt{net/http}: mã yêu cầu, nhật ký, trace~\cite{opentelemetry}, giới hạn thân 1~MB, phiên bản ứng dụng, giới hạn tần suất và lũy đẳng (chi tiết ở Hình~\ref{fig:pipeline-yeu-cau}). Bộ đếm giới hạn tần suất nằm trong Redis dưới khóa \texttt{rl:*} để mọi bản API dùng chung một hạn mức (\nfr{NFR17}). Lớp lũy đẳng của \texttt{SyncEngine} làm việc với bảng \texttt{sync.idempotency\_keys}: câu lệnh đầu tiên của giao dịch ghi chèn khóa ở trạng thái \texttt{in\_progress}, nên hai bản sao đồng thời của một yêu cầu va nhau ngay tại khóa chính; các API công khai như đăng ký, đăng nhập, làm mới phiên không dùng cơ chế này vì đã có giới hạn tần suất và OTP. Thành phần xác thực và phân quyền kiểm JWT, đối chiếu \texttt{token\_version} trong Redis (\texttt{auth:ver}) để thu hồi tức thời, rồi kiểm tra người gọi là thành viên không gian trong đường dẫn với đúng vai trò (\texttt{member:*}, \nfr{NFR16}, \br{BR23}). Hub WebSocket cũng gọi thành phần này để xác thực kết nối \texttt{/ws} trước khi nạp danh sách không gian của người dùng và đăng ký kênh \texttt{space:*} tương ứng cùng kênh riêng \texttt{user:*}; mỗi bản API chỉ đăng ký kênh của những người đang nối với chính nó. Handler sinh từ OpenAPI chỉ kiểm tra đầu vào, ánh xạ DTO và lỗi, rồi gọi dịch vụ ứng dụng, được phóng to ở Hình~\ref{fig:c4-thanh-phan-api-ra}.

\diagram[\textwidth]{c3-05b-c4-thanh-phan-api-ra}{Sơ đồ thành phần của vùng chứa API theo mô hình C4 (mức 3): phía ra}{fig:c4-thanh-phan-api-ra}

Dịch vụ ứng dụng là ranh giới giao dịch: nó mở giao dịch qua \texttt{TxManager} rồi gọi các adapter ra. Repository sinh bởi sqlc luôn lọc \texttt{space\_id} và khóa dòng tường minh khi cần; event publisher chỉ chèn tác vụ River trong giao dịch; \texttt{ChangefeedAppender} là lời gọi cuối cùng trước \texttt{COMMIT}, tăng \texttt{change\_seq}, ghi \texttt{space\_changes} và đăng ký hàm phát \texttt{space:\{id\}} chạy sau commit. Bộ đệm cấu hình và gợi ý giữ \texttt{cfg:all} trong bộ nhớ tiến trình với thời hạn sống 5 phút, đồng thời đăng ký kênh \texttt{cfg} để xóa ngay khi quản trị viên đổi tham số; kết quả gợi ý món được lưu ở \texttt{suggest:\{sid\}} và bị xóa khi lô, công thức hoặc trọng số gợi ý thay đổi. Search client luôn thêm điều kiện lọc theo không gian vào truy vấn Elasticsearch và tự chuyển sang trigram trên PostgreSQL khi Elasticsearch lỗi; S3 presigner chỉ ký URL và kiểm tra đối tượng bằng \texttt{HEAD}, không truyền nội dung ảnh qua API. Hình~\ref{fig:c4-thanh-phan-worker} phóng to vùng chứa Worker, chạy cùng tệp thực thi và dùng lại dịch vụ ứng dụng của các module.

\diagram[\textwidth]{c3-06-c4-thanh-phan-worker}{Sơ đồ thành phần của vùng chứa Worker theo mô hình C4 (mức 3)}{fig:c4-thanh-phan-worker}

River client bầu một bản làm leader qua bảng \texttt{river\_leader}; bộ lập lịch chỉ chạy trên leader nên tăng số bản worker không nhân đôi tác vụ định kỳ. Hàng đợi \texttt{notify} có ưu tiên cao nhất để \nfr{NFR12} không bị lô đánh chỉ mục lớn làm chậm. Luồng \uc{UC14} gồm ba tác vụ nối tiếp. Mỗi phút, \texttt{ReminderDispatcher} kiểm tra giờ Việt Nam đã tới \texttt{reminder.runAt} và hôm nay chưa chạy (tác vụ duy nhất theo ngày), rồi chèn một \texttt{UserReminder} cho mỗi người dùng có lô nằm trong cửa sổ nhắc, duy nhất theo người và ngày. \texttt{UserReminderWorker} bỏ qua người đã tắt nhắc, lọc lô trên mọi không gian của người đó theo ngưỡng $N$ riêng (\br{BR12}), nêu tên tối đa \texttt{reminder.maxLotsPerNotice} lô (mặc định 5), rồi chèn thông báo với \texttt{dedupe\_key} theo người và ngày; chỉ khi bản ghi thực sự mới được tạo thì tác vụ \texttt{PushNotification} mới được chèn trong cùng giao dịch, nên mỗi người nhận tối đa một thông báo nhắc mỗi ngày (\br{BR13}). \texttt{NotifyWorker} xử lý \texttt{PushNotification} và \texttt{SendEmail}: gửi FCM tới từng thiết bị, ghi \texttt{deliveries} và bỏ qua thiết bị đã gửi thành công, xóa thiết bị có mã không còn hợp lệ, thử lại tối đa ba lần với khoảng chờ tăng lũy thừa khi lỗi tạm thời. Cũng trên hàng đợi này, \texttt{AssignmentDigest} gộp các phân công cho cùng một người trong cùng danh sách trong 5 giây. Nhóm tác vụ dữ liệu gồm \texttt{CopySpaceData} khi tạo nhóm kèm dữ liệu cá nhân, \texttt{PurgeSpace} khi giải tán nhóm và \texttt{RemapFood} khi quản trị viên gộp thực phẩm. \texttt{MaintenanceWorker} chạy hằng đêm: làm mới materialized view \texttt{report.daily\_space\_stats} bằng \texttt{REFRESH MATERIALIZED VIEW CONCURRENTLY} để báo cáo không bị chặn, và dọn OTP, khóa lũy đẳng, phiên, tombstone hết hạn.

\subsection{Ghi nhận các quyết định kiến trúc}\label{ssec:adr}

Mỗi quyết định có ảnh hưởng rộng được ghi thành một bản ghi quyết định kiến trúc (Architecture Decision Record, ADR) gồm bối cảnh, quyết định và hệ quả, lấy đánh đổi về thuộc tính chất lượng làm trung tâm~\cite{bass}. Bảng~\ref{tab:c3-adr} tóm tắt mười quyết định; bản đầy đủ lưu tại \path{docs/adr/}.

{\footnotesize\setlength{\tabcolsep}{4pt}
\begin{longtable}{|C{1.5cm}|L{3.2cm}|L{\dimexpr(\textwidth-8\tabcolsep-5\arrayrulewidth-1.5cm-3.2cm)/2\relax}|L{\dimexpr(\textwidth-8\tabcolsep-5\arrayrulewidth-1.5cm-3.2cm)/2\relax}|}
\caption{Các quyết định kiến trúc chính}\label{tab:c3-adr}\\
\hline
\hd{Mã} & \hd{Quyết định} & \hd{Bối cảnh, lý do} & \hd{Hệ quả, đánh đổi}\\ \hline
\endfirsthead
\multicolumn{4}{l}{\footnotesize\itshape Bảng \thetable\ (tiếp theo)}\\ \hline
\hd{Mã} & \hd{Quyết định} & \hd{Bối cảnh, lý do} & \hd{Hệ quả, đánh đổi}\\ \hline
\endhead
\mbox{ADR-01} & Nguyên khối chia module, lục giác trong module & Quy mô nhỏ; \uc{UC11}, \uc{UC19} nguyên tử; nhóm 4 người & Một giao dịch, vận hành đơn giản; nhân bản cả khối, cần kỷ luật ranh giới\\ \hline
\mbox{ADR-02} & Go 1.23, một tệp thực thi hai chế độ & Tệp tĩnh, image nhỏ, goroutine hợp với WebSocket (mục~\ref{sec:cong-nghe}) & Phải tự ghép chi, sqlc, River; xử lý lỗi tường minh\\ \hline
\mbox{ADR-03} & PostgreSQL là nguồn sự thật, mỗi module một schema & Dữ liệu quan hệ, ràng buộc, ACID, khóa dòng cho FEFO & Một nơi sao lưu; khóa ngoại liên schema bị giới hạn\\ \hline
\mbox{ADR-04} & River làm hàng đợi và outbox thay message broker & Cần chèn tác vụ cùng giao dịch, định kỳ, duy nhất, thử lại & Không ghi kép, không thêm hệ thống; hàng đợi dùng tài nguyên PostgreSQL\\ \hline
\mbox{ADR-05} & ES cho thực phẩm, công thức (bản chung và bản riêng của không gian); PostgreSQL cho lô, mặt hàng & Tìm không dấu, chịu lỗi gõ, gợi ý $\le$ 500 ms (\nfr{NFR09}) cần cả cho thực phẩm riêng, công thức của nhóm (\uc{UC11}, \uc{UC20}); lô, mặt hàng ít, đổi liên tục, phải thấy ngay & Mọi truy vấn ES bắt buộc lọc \texttt{scope = global} hoặc \texttt{owner\_space\_id = sid} (công thức: mẫu đã công bố hoặc \texttt{space\_id = sid}) ngay trong search client (\br{BR23}); xóa tài liệu khi giải tán nhóm, xóa tài khoản (\nfr{NFR18}); cần đối soát, dự phòng trigram\\ \hline
\mbox{ADR-06} & WebSocket chỉ gửi gợi ý, client kéo delta & Nhiều thiết bị, mạng chập chờn, socket bị ngắt & Một đường đọc cho trực tuyến và ngoại tuyến; thêm một vòng yêu cầu, vẫn $\le$ 3 s\\ \hline
\mbox{ADR-07} & UUIDv7 sinh ở client, \texttt{Idempotency-\allowbreak Key} cho lệnh ghi đã xác thực & Tạo bản ghi ngoại tuyến, phát lại an toàn~\cite{idempotencykey} & Không cần mã tạm, UUIDv7~\cite{rfc9562} tốt cho chỉ mục; khóa chèn đầu giao dịch, chỉ lưu kết quả thành công, giữ 7 ngày\\ \hline
\mbox{ADR-08} & Trạng thái hạn dùng suy diễn lúc truy vấn & Ngưỡng $N$ thuộc người xem (\br{BR12}) & Không cần tác vụ hằng đêm, không lệch giữa thành viên; lọc qua \texttt{expires\_on}\\ \hline
\mbox{ADR-09} & JWT 30 phút ký EdDSA, refresh token xoay vòng & API phi trạng thái; thu hồi khi đăng xuất (\br{BR03}, \nfr{NFR15}) & Không truy vấn CSDL mỗi yêu cầu; cần \texttt{token\_version}, xoay khóa\\ \hline
\mbox{ADR-10} & Docker Compose trên một VPS (GĐ1) & Tải giả định nhỏ, ngân sách, nhóm nhỏ & Đơn giản, rẻ; một điểm lỗi, bù bằng sao lưu WAL; GĐ2 khi vượt 50.000 người dùng (Hình~\ref{fig:mo-rong})\\ \hline
\end{longtable}
}

ADR-01 và ADR-04 cùng dựa trên tiền đề một CSDL nên phải được xem lại cùng lúc khi tách \texttt{notification}; ADR-06 đến ADR-08 cùng phục vụ đồng bộ ngoại tuyến. ADR-05 là quyết định có rủi ro riêng tư cao nhất: đưa thực phẩm riêng và công thức của nhóm vào Elasticsearch giúp một đường gợi ý duy nhất đạt \nfr{NFR09}, nhưng biến điều kiện lọc theo không gian thành một quy tắc an ninh, nên điều kiện này được dựng tập trung trong search client chứ không để từng handler tự thêm, và có kiểm thử tích hợp riêng chứng minh không gian khác không thấy tài liệu của nhau. Khi số liệu vận hành vượt giả định mà một ADR dựa vào, ADR đó được mở lại.

\InputIfFileExists{chapters/c3/tich-hop-di-dong.tex}{}{}
% =====================================================================
%  3.4. THIẾT KẾ CHI TIẾT (3.4.1–3.4.4)
%  Mục 3.4.5 "Các kỹ thuật nổi bật" nằm ở tệp ky-thuat.tex.
% =====================================================================
\section{Thiết kế chi tiết}\label{sec:thiet-ke-chi-tiet}

Mục này đi xuống mức mã nguồn, tương ứng mức 4 của mô hình C4~\cite{c4model}, gồm tổ chức mã và phân tầng (mục~\ref{ssec:to-chuc-ma-nguon}), biểu đồ lớp (mục~\ref{ssec:thiet-ke-lop}), chuỗi xử lý một yêu cầu (mục~\ref{ssec:luong-yeu-cau}) và máy trạng thái phía máy chủ (mục~\ref{ssec:may-trang-thai}). Mục~\ref{ssec:ky-thuat} trình bày các kỹ thuật cho những bài toán khó nhất: khóa dòng theo FEFO, transactional outbox, đồng bộ ngoại tuyến, đẩy thay đổi thời gian thực, nhắc hạn không trùng lặp, gợi ý món và tìm kiếm tiếng Việt; mục~\ref{ssec:hien-thuc-uc} khép lại bằng cách hiện thực hóa các use case trọng tâm trên nền các thành phần này.

% ---------------------------------------------------------------------
\subsection{Tổ chức mã nguồn và phân tầng trong module}\label{ssec:to-chuc-ma-nguon}

Toàn bộ backend nằm trong một kho mã và một module Go~\cite{golang}, với cây thư mục rút gọn ở Mã nguồn~\ref{code:c3-cay-thu-muc}. Trình biên dịch Go cấm mã bên ngoài nhập các gói dưới \texttt{internal}, nên ranh giới công khai duy nhất là hợp đồng API. Bên trong mỗi module, nhóm áp dụng kiến trúc lục giác, hay cổng và bộ chuyển đổi (ports and adapters)~\cite{hexagonal}.

\codecaption{Cây thư mục kho mã backend (rút gọn)}{code:c3-cay-thu-muc}
\begin{Verbatim}
dicho-backend/
|-- cmd/api/, cmd/worker/  # composition root của hai tiến trình
|-- api/openapi.yaml       # hợp đồng OpenAPI 3.1, nguồn sinh mã
|-- db/migrations/         # tệp goose (expand-contract)
|-- db/queries/<module>/   # câu SQL có tên, đầu vào của sqlc
|-- internal/platform/     # config, db (pgxpool, TxManager), httpx,
|                          # auth, events, clock, i18n, observability
|-- internal/kernel/       # Quantity, Unit, Money, DateRange, ID
|-- internal/modules/
|   |-- pantry/
|   |   |-- api/           # cổng vào: StockIn, StockQuery...
|   |   |-- domain/        # FoodLot, StockMovement, FEFOAllocator
|   |   |-- app/           # PantryService và các cổng ra
|   |   `-- adapter/http/, adapter/postgres/  # handler; mã sqlc
|   |-- shopping/ mealplan/ recipe/ household/ ...
|   `-- notification/      # thêm adapter/river: worker đẩy, thư
|-- deploy/compose/        # docker-compose.*.yml, Caddyfile
`-- .gitlab-ci.yml         # verify, test, build, scan, deploy
\end{Verbatim}

Tệp \path{api/openapi.yaml} là nguồn sự thật của hợp đồng API theo hướng thiết kế API trước (contract-first)~\cite{openapi}: \texttt{oapi-codegen} sinh interface \texttt{ServerInterface} cho máy chủ, ứng dụng Flutter sinh client từ cùng tệp (\nfr{NFR20}). Câu SQL trong \path{db/queries} được \texttt{sqlc}~\cite{sqlc} biên dịch thành hàm Go an toàn kiểu; CI so sánh mã sinh lại với mã đã commit để chặn lệch hợp đồng hoặc lược đồ. Hình~\ref{fig:phan-tang-module} thể hiện cấu trúc lục giác của một module, lấy \texttt{shopping} làm ví dụ cho hai loại cổng nối với module khác.

\diagram[\textwidth]{c3-07-phan-tang-module}{Phân tầng lục giác bên trong một module}{fig:phan-tang-module}

Lõi lục giác gồm tầng ứng dụng và tầng miền. Tầng miền chỉ nhập thư viện chuẩn và \texttt{kernel}. Tầng ứng dụng có một phương thức cho mỗi use case, là ranh giới giao dịch, nơi kiểm tra quyền theo vai trò trong nhóm, và tự khai báo các cổng ra mà nó cần (repository, \texttt{ChangefeedAppender}, \texttt{EventPublisher}) theo thông lệ “interface đặt ở phía dùng” của Go. Trên mép lục giác có hai loại cổng nối module với module. Cổng vào là gói \texttt{<module>/api} do chính module công bố, ví dụ \path{shopping/api.ListComposer}: module khác (ở đây là \texttt{mealplan}) chỉ thấy interface này, còn dịch vụ ứng dụng hiện thực nó, nên cổng vào không phải một adapter mà là một phần hợp đồng của lõi. Cổng ra tới module khác là interface do module bên kia công bố, ví dụ \texttt{shopping} gọi \path{pantry/api.StockIn} khi nhập tủ và \texttt{PantryService} của \texttt{pantry} hiện thực nó. Adapter vào (handler HTTP sinh từ OpenAPI, worker River) gọi tầng ứng dụng; adapter ra (repository dùng mã \texttt{sqlc}, \texttt{RiverPublisher} chèn tác vụ bằng \texttt{InsertTx}, \texttt{SyncEngine} của module \texttt{sync}) hiện thực cổng ra. Mọi mũi tên phụ thuộc hướng vào lõi, và một module chỉ được nhập gói \texttt{api} của module khác theo đồ thị không chu trình ở Hình~\ref{fig:phu-thuoc-module}; luật \texttt{depguard} của \texttt{golangci-lint} kiểm tra điều này trong CI. Nhờ vậy tầng miền được kiểm thử đơn vị không cần cơ sở dữ liệu (\nfr{NFR20}).

Có ba chỗ chiều gọi ngược với chiều phụ thuộc cho phép. Khi thành viên rời hoặc bị xóa khỏi nhóm, \texttt{household} cần \texttt{shopping} trả các mặt hàng đang giao cho người đó về Cần mua (\br{BR07}), trong khi \texttt{shopping} đã phụ thuộc \texttt{household}. Nhóm áp dụng nguyên lý đảo phụ thuộc: \texttt{household} công bố interface \path{household/api.MemberLeaveHook}, \texttt{ShoppingService} hiện thực và được đăng ký tại gốc lắp ghép, rồi \texttt{household} gọi hook trong chính giao dịch xóa thành viên. Hai trường hợp còn lại dùng cùng khuôn: \path{identity/api.ActivationHook} do \texttt{household} hiện thực để tạo không gian cá nhân ngay khi tài khoản được kích hoạt (\uc{UC02}, \br{BR04}), và \path{recipe/api.RecipeDeleteHook} do \texttt{mealplan} hiện thực để chuyển các món dự kiến thành món tự do khi công thức bị xóa (\uc{UC33}).

Phụ thuộc được lắp ghép thủ công qua hàm khởi tạo tại gốc lắp ghép (composition root), tức hàm \texttt{main} của hai tiến trình: tạo pool, client Redis và River, rồi repository, dịch vụ, handler, và nối cổng giữa module, ví dụ \texttt{ShoppingService} nhận \texttt{PantryService} dưới dạng \texttt{pantryapi.StockIn}. \texttt{TxManager.WithinTx(ctx, fn)} mở \texttt{pgx.Tx} và gắn vào context; repository lấy đối tượng truy vấn qua \texttt{db.Querier(ctx)}, dùng giao dịch nếu context đang mang. Lời gọi lồng nhau tham gia giao dịch sẵn có, nên khi nhập tủ (\uc{UC21}) việc tạo lô, đổi trạng thái mặt hàng và ghi nhật ký thay đổi qua cổng \texttt{ChangefeedAppender}, do \texttt{SyncEngine} của module \texttt{sync} hiện thực, nằm trong một giao dịch mà không module nào biết \texttt{pgx.Tx}; tác dụng phụ ngoài cơ sở dữ liệu chỉ chạy sau commit.

% ---------------------------------------------------------------------
\subsection{Thiết kế lớp (C4 mức 4)}\label{ssec:thiet-ke-lop}

Biểu đồ lớp ánh xạ \texttt{struct} của Go thành lớp và \texttt{interface} thành lớp «interface», thuộc tính ghi theo kiểu Go (\texttt{*T} là giá trị có thể rỗng), ký hiệu theo UML~2.5~\cite{uml251}. Mọi phương thức của dịch vụ, cổng và repository nhận \texttt{ctx context.Context} làm tham số đầu và trả thêm \texttt{error}; hai thành phần lặp lại này được lược bỏ trên hình, cùng với kiểu lệnh và kiểu kết quả phụ. Tên dịch vụ, phương thức và cổng trên hình là tên dùng thống nhất trong các sơ đồ tuần tự của chương. Hình~\ref{fig:lop-tu-lanh} trình bày module \texttt{pantry}, trong đó \texttt{FoodLot} là gốc tập hợp (aggregate root)~\cite{evans}.

\diagram[\textwidth]{c3-08-lop-tu-lanh}{Biểu đồ lớp module tủ lạnh (pantry)}{fig:lop-tu-lanh}

\texttt{PantryService} hiện thực ba cổng vào của gói \path{pantry/api}: \path{StockIn.AddLots} cho nhập tủ từ danh sách mua sắm (\uc{UC21}), \path{StockConsumer.ConsumeForMeal} cho trừ kho khi nấu (\uc{UC41}) và \texttt{StockQuery} gồm \texttt{Available} (tồn kho còn hạn đến một ngày, dùng cho gợi ý món, sinh danh sách và đối chiếu nguyên liệu) cùng \texttt{SearchLots} (tìm lô trong không gian). Các phương thức riêng \texttt{AddLot}, \texttt{UpdateLot}, \texttt{Consume}, \texttt{Discard} phục vụ \uc{UC25}, \uc{UC27}, \uc{UC28}, \uc{UC29}. Mọi thay đổi lượng tồn đi qua phương thức của \texttt{FoodLot}, nơi kiểm tra bất biến (\br{BR11}, \br{BR14}, \br{BR16}) và sinh một \texttt{StockMovement} chỉ ghi thêm. Riêng \texttt{Adjust} phục vụ sửa lượng còn lại ở \uc{UC27}: khác các trường mô tả như vị trí, ngày hết hạn, ghi chú vốn được hợp nhất theo trường và lệnh sau thắng, lượng còn lại không áp dụng quy tắc này, \texttt{If-Match} lệch thì trả 409/4002, giá trị vượt lượng ban đầu trả 400/2002, giá trị 0 chuyển lô sang \texttt{consumed}, và mỗi lần đổi sinh một biến động loại \texttt{adjust}. Thuộc tính \texttt{TotalCost} lưu thành tiền lúc nhập tủ thay cho đơn giá, nên giá trị lãng phí tính theo tỷ lệ lượng bỏ đi trên lượng ban đầu vẫn đúng khi đơn vị của lô khác đơn vị mua. \texttt{Quantity} là value object của \texttt{kernel} (số thập phân chính xác kèm đơn vị, quy đổi về đơn vị cơ sở theo \br{BR16}) nên chỉ xuất hiện dưới dạng kiểu thuộc tính. \texttt{FEFOAllocator} là hàm thuần nhận các lô đã khóa theo thứ tự \texttt{(expires\_on, id)} (\br{BR15}) và trả phần phân bổ cùng lượng còn thiếu, còn \texttt{PantryService} chỉ điều phối: khóa lô, phân bổ, lưu, ghi nhật ký thay đổi, phát \path{lot.changed}. Cơ chế khóa dòng được phân tích ở mục~\ref{sssec:kt-fefo}.

Hai module \texttt{shopping} và \texttt{mealplan} liên kết chặt qua \uc{UC22} nhưng mỗi module có đủ lớp để cần một hình riêng: Hình~\ref{fig:lop-mua-sam-ke-hoach} trình bày \texttt{shopping}, Hình~\ref{fig:lop-ke-hoach} trình bày \texttt{mealplan}; ở mỗi hình, cổng của module khác chỉ được vẽ ở mức interface. Trong module \texttt{shopping}, năm phương thức \texttt{AddItem}, \texttt{AssignItem}, \texttt{MarkPurchased}, \texttt{UndoPurchase} và \texttt{StockItems} của \texttt{ShoppingService} lần lượt phục vụ \uc{UC17}, \uc{UC19}, \uc{UC20} (hai lệnh đánh dấu và hoàn tác) và \uc{UC21}. Dịch vụ còn hiện thực hai cổng do module khác gọi: \path{shopping/api.ListComposer} cho \texttt{mealplan} (\texttt{PendingQuantities} trả lượng đang chờ mua theo thực phẩm, \texttt{Compose} tạo hoặc gộp danh sách) và \path{household/api.MemberLeaveHook} cho \texttt{household}.

\diagram[\textwidth]{c3-09-lop-mua-sam-ke-hoach}{Biểu đồ lớp module mua sắm (shopping)}{fig:lop-mua-sam-ke-hoach}

Như Hình~\ref{fig:lop-mua-sam-ke-hoach} cho thấy, quy tắc trạng thái mặt hàng nằm ở \texttt{ShoppingItem}, còn quy tắc trạng thái danh sách nằm ở gốc tập hợp \texttt{ShoppingList}: \texttt{DeriveStatus} suy ra \texttt{ListStatus} theo \br{BR08}, \texttt{CheckEditable} từ chối sửa danh sách Hoàn tất trừ các lệnh được phép sau khi mua (mục~\ref{ssec:may-trang-thai}). Quy tắc vai trò theo \br{BR10} nằm ở \texttt{AssignmentPolicy}, nên đổi chính sách phân quyền không chạm vào máy trạng thái. \texttt{ShoppingService} chỉ điều phối các lớp này trong một giao dịch và gọi \path{StockIn.AddLots} khi nhập tủ.

\diagram[\textwidth]{c3-09b-lop-ke-hoach}{Biểu đồ lớp module kế hoạch bữa ăn (mealplan)}{fig:lop-ke-hoach}

Trong Hình~\ref{fig:lop-ke-hoach}, \texttt{MealPlanService} có \texttt{PlanMeals} (\uc{UC37}), \texttt{UpdateMeal} (\uc{UC39}), \texttt{Cook} (\uc{UC41}), \texttt{Suggest} (\uc{UC40}) và \texttt{GenerateList} (\uc{UC22}), đồng thời hiện thực \path{recipe/api.RecipeDeleteHook}. Với \uc{UC22}, handler \texttt{MealHandler} nhận lệnh S15 và gọi \texttt{GenerateList}; dịch vụ lấy nguyên liệu qua \path{RecipeQuery.Ingredients}, tồn kho qua \path{StockQuery.Available}, lượng đang chờ mua qua \path{ListComposer.PendingQuantities}, giao phép tính cho \path{ShortageCalculator.Compute} rồi ghi danh sách qua \path{ListComposer.Compose}. Vì \texttt{ListComposer} do \texttt{shopping} công bố, \texttt{mealplan} phụ thuộc \texttt{shopping} mà không có chiều ngược lại. \texttt{Cook} gọi \path{StockConsumer.ConsumeForMeal} trong cùng giao dịch với \path{MealEntry.MarkCooked}, bảo đảm \uc{UC41} nguyên tử (\nfr{NFR11}). \path{SuggestionEngine.Rank} và \path{ShortageCalculator.Compute} là hai dịch vụ miền thuần ngang hàng: cả hai nhận dữ liệu đã nạp sẵn và không truy cập cơ sở dữ liệu. Dịch vụ miền nhận ngày hiện tại và tham số cấu hình làm đối số, ví dụ \texttt{Reschedule} nhận cửa sổ lập kế hoạch đọc từ khóa \path{mealplan.horizonDays} (mặc định 28 ngày, \br{BR18}), nên ví dụ kiểm chứng ở Chương~\ref{chap:dac-ta} như cà rốt thiếu 0,3~kg trở thành ca kiểm thử đơn vị. Bảng~\ref{tab:c3-lop-chinh} tóm tắt trách nhiệm các lớp chính của ba biểu đồ lớp.

{\small
\begin{longtable}{|L{4.3cm}|L{\dimexpr\textwidth-4\tabcolsep-3\arrayrulewidth-4.3cm\relax}|}
\caption{Trách nhiệm của các lớp chính}\label{tab:c3-lop-chinh}\\
\hline
\hd{Lớp} & \hd{Trách nhiệm chính}\\ \hline
\endfirsthead
\multicolumn{2}{l}{\footnotesize\itshape Bảng \thetable\ (tiếp theo)}\\ \hline
\hd{Lớp} & \hd{Trách nhiệm chính}\\ \hline
\endhead
\texttt{FoodLot} & Gốc tập hợp; giữ bất biến của lô, sinh biến động kho (kể cả \texttt{adjust}), tính trạng thái hạn dùng suy diễn.\\ \hline
\texttt{FEFOAllocator} & Dịch vụ miền; phân bổ lượng trừ theo ngày hết hạn tăng dần, bỏ qua lô quá hạn (\br{BR15}).\\ \hline
\texttt{PantryService} & Dịch vụ ứng dụng của \uc{UC25}, \uc{UC27}--\uc{UC29} (\texttt{AddLot}, \texttt{UpdateLot}, \texttt{Consume}, \texttt{Discard}); hiện thực cổng \texttt{StockIn}, \texttt{StockConsumer}, \texttt{StockQuery}.\\ \hline
\texttt{ShoppingList} & Gốc tập hợp; gộp mặt hàng trùng (\br{BR09}), suy ra trạng thái danh sách và chặn sửa danh sách Hoàn tất (\br{BR08}).\\ \hline
\texttt{ShoppingItem} & Thực thể; hiện thực máy trạng thái mặt hàng ở mục~\ref{ssec:may-trang-thai}.\\ \hline
\texttt{ShoppingService} & Dịch vụ ứng dụng của \uc{UC17}--\uc{UC21}; gọi \texttt{StockIn.AddLots} trong cùng giao dịch khi nhập tủ; hiện thực \texttt{ListComposer} và \texttt{MemberLeaveHook}.\\ \hline
\texttt{AssignmentPolicy} & Dịch vụ miền; quyền phân công và mua thay theo \br{BR10}.\\ \hline
\texttt{MealEntry} & Gốc tập hợp; ràng buộc \br{BR18} với cửa sổ \texttt{mealplan.horizonDays} đọc từ cấu hình, chuyển sang đã nấu đúng một lần.\\ \hline
\texttt{MealPlanService} & Dịch vụ ứng dụng của \uc{UC37}, \uc{UC39}--\uc{UC41}, \uc{UC22}; hiện thực \texttt{RecipeDeleteHook}.\\ \hline
\texttt{SuggestionEngine} & Dịch vụ miền; tính $R$, $E$, $F$ và điểm $S$ theo \br{BR19}, trọng số và ngưỡng đọc từ cấu hình.\\ \hline
\texttt{ShortageCalculator} & Dịch vụ miền; tính phần thiếu theo ngày nấu, mỗi lượng chỉ tính một lần (\br{BR20}).\\ \hline
\texttt{SyncEngine} & Thuộc module \texttt{sync}; hiện thực \texttt{ChangefeedAppender} (tăng \texttt{change\_seq}, ghi \texttt{space\_changes}) và quản lý khóa lũy đẳng.\\ \hline
\end{longtable}
}

% ---------------------------------------------------------------------
\subsection{Luồng xử lý một yêu cầu}\label{ssec:luong-yeu-cau}

Mọi yêu cầu tới tiến trình \texttt{api} đi qua cùng một chuỗi bước ở Hình~\ref{fig:pipeline-yeu-cau}, sắp theo nguyên tắc kiểm tra rẻ trước, kiểm tra đắt sau để từ chối sớm yêu cầu không hợp lệ trước khi chiếm kết nối cơ sở dữ liệu.

\diagram[0.92\textwidth]{c3-10-pipeline-yeu-cau}{Chuỗi xử lý một yêu cầu HTTP trong tiến trình api}{fig:pipeline-yeu-cau}

Mã \texttt{X-Request-Id} được gắn đầu tiên để mọi log và span cùng định danh; bước ghi log và trace bao phần còn lại nên độ trễ đo được gồm cả yêu cầu bị từ chối (\nfr{NFR09}). Giới hạn thân 1~MB chặn tiêu tốn tài nguyên (API4 của OWASP~\cite{owaspapi}) và trả HTTP 413 với mã 2002; kiểm tra phiên bản đứng trước xác thực để ứng dụng quá cũ luôn nhận mã 4260. Giới hạn tần suất chia làm hai bước vì hai loại khóa đếm khác nhau. Bước 6 chỉ áp cho tuyến công khai như đăng ký, đăng nhập, gửi lại OTP và đếm theo địa chỉ IP, do lúc này chưa biết người gọi (\nfr{NFR17}). Bước 7 so claim \texttt{ver} của JWT với \texttt{token\_version} lưu đệm trong Redis, nên token mất hiệu lực ngay khi đăng xuất, đổi mật khẩu hoặc bị khóa; với nhóm đường dẫn \path{/admin/*}, bước này còn đòi \texttt{system\_role = admin}, thiếu thì trả 403/3007. Khi đã biết người gọi, bước 7b đếm theo người dùng, gồm luật chung cho mọi API và luật riêng như nhập mã mời, gợi ý tên, tra mã vạch; vượt ngưỡng trả 429/4290 kèm \texttt{Retry-After}. Bước 8 xác minh người gọi là thành viên không gian trong đường dẫn và đặt vai trò trong nhóm vào context; người ngoài nhận 404/4001 thay vì 403 để không xác nhận không gian tồn tại (\br{BR23}, \nfr{NFR16}), và mọi truy vấn repository vẫn lọc thêm \texttt{space\_id = \$sid}. Chi tiết bảo mật nằm ở mục~\ref{sec:bao-mat}.

Bước 9 chỉ đọc \texttt{Idempotency-Key} và băm phương thức, đường dẫn, thân yêu cầu vào context; việc kiểm tra thật nằm ở câu lệnh đầu tiên của giao dịch ở bước 11, do \texttt{SyncEngine} thực hiện: \texttt{INSERT} vào \path{sync.idempotency_keys} một dòng \texttt{status = 'in\_progress'} với \texttt{ON CONFLICT DO NOTHING}. Nếu khóa đã tồn tại với băm khác, yêu cầu nhận 400/2002; khóa đã \texttt{completed} được phát lại đúng phản hồi đã lưu kèm \texttt{Idempotent-Replayed: true}; ngay trước \texttt{COMMIT}, dòng khóa được cập nhật thành \texttt{completed} cùng phản hồi. Đặt khóa vào chính giao dịch nghiệp vụ cho hai tính chất. Lỗi nghiệp vụ dẫn tới \texttt{ROLLBACK} nên khóa biến mất cùng dữ liệu, tức chỉ kết quả thành công được lưu và lần gửi lại sau khi người dùng sửa dữ liệu vẫn được xử lý. Hai bản sao đồng thời của cùng yêu cầu cũng không thể cùng chạy, vì PostgreSQL buộc lệnh \texttt{INSERT ... ON CONFLICT} thứ hai chờ giao dịch đang giữ khóa trùng kết thúc: bản đầu commit thì bản sau thấy \texttt{completed} và phát lại, bản đầu rollback thì bản sau tự chèn được và xử lý bình thường. Nếu phải chờ quá \texttt{lock\_timeout}, khóa được coi là đang ở \texttt{in\_progress} và yêu cầu nhận 409/4002 kèm \texttt{Retry-After}. Sau khi khóa lũy đẳng được giữ, dịch vụ kiểm tra vai trò trong nhóm cho thao tác chỉ dành cho trưởng nhóm (403/3007) và so \texttt{If-Match} với cột \texttt{version} cho lệnh không được hợp nhất (409/4002), rồi mới chạy logic miền, ghi qua repository, ghi nhật ký thay đổi và chèn tác vụ River bằng \texttt{InsertTx}. Các khóa dòng được lấy theo thứ tự toàn cục trình bày ở mục~\ref{sssec:kt-fefo}, trong đó dòng \texttt{change\_seq} của không gian luôn được cập nhật cuối cùng.

Lỗi được ánh xạ tại một nơi: tầng miền trả lỗi có kiểu, \texttt{httpx.WriteError} dùng \texttt{errors.As} chuyển thành mã HTTP và mã phản hồi, ví dụ không tìm thấy thành 404/4001, xung đột phiên bản thành 409/4002 (\msg{MSG35}), vượt lượng còn lại thành 422/2003 (\msg{MSG17}), vi phạm quy tắc thành 422/4004 kèm mã BR. Khi lỗi xảy ra trong giao dịch, \texttt{TxManager} rollback trước rồi mới chuyển lỗi tới bước 14, như mũi tên từ bước 11 trên hình. Vi phạm ràng buộc cơ sở dữ liệu được repository dịch theo SQLSTATE và tên ràng buộc thành lỗi miền như Bảng~\ref{tab:c3-anh-xa-loi}, nên cùng mã 23505 có thể thành 4003, 4004 hay 4002 tùy ý nghĩa nghiệp vụ. Lỗi không nhận diện được luôn thành 500/5000, chi tiết chỉ ghi vào log kèm \texttt{traceId}. Handler không tự ghi mã trạng thái, nên mọi API dùng chung bộ mã ở mục~\ref{sec:api} (\nfr{NFR21}), thông điệp lấy từ Phụ lục~\ref{pl:thong-diep} theo \texttt{Accept-Language}. Sau commit, bước 13 phát \texttt{PUBLISH space:\{sid\}} để các thiết bị kéo thay đổi, và xóa \texttt{suggest:\{sid\}} khi lô hoặc công thức đổi để gợi ý món không dùng tồn kho cũ.

Mỗi yêu cầu có hạn chót 5~giây trong context, còn kết nối PostgreSQL đặt \texttt{statement\_timeout} và \texttt{lock\_timeout} cùng 3~giây: hạn câu lệnh ngắn hơn để cơ sở dữ liệu tự hủy truy vấn chậm thay vì giữ kết nối sau khi client đã bỏ cuộc~\cite{nygard}. Theo định luật Little $L = \lambda W$, đỉnh 50 yêu cầu/giây với thời gian giữ kết nối 20~ms (giả định thiết kế) chỉ cần trung bình 1 kết nối; nhóm đặt tối đa 20 kết nối cho mỗi bản \texttt{api} và 10 cho \texttt{worker}, tổng 50, dưới giới hạn 100 mặc định của PostgreSQL và còn dư cho báo cáo, đợt tăng tải, sao lưu và giám sát.

Khi gặp khóa chết (SQLSTATE 40P01) hoặc chờ khóa quá \texttt{lock\_timeout} (55P03), \texttt{TxManager} chạy lại cả hàm giao dịch tối đa 3 lần, chờ tăng lũy thừa từ 20~ms kèm nhiễu ngẫu nhiên (full jitter) và dừng sớm nếu hạn chót của context không đủ cho lần chạy tiếp; hết lượt thì trả 503/5003 (\msg{MSG33}). Ngoại lệ duy nhất là 55P03 phát sinh ngay ở câu chèn khóa lũy đẳng: chạy lại chỉ chờ tiếp chính yêu cầu đang xử lý, nên lỗi này được trả thẳng 409/4002 như trên. Chạy lại an toàn vì mọi tác dụng phụ, kể cả tác vụ River và khóa lũy đẳng, nằm trong giao dịch hoặc chạy sau commit. Mọi giao dịch ghi chạy ở mức \texttt{READ COMMITTED} với khóa dòng tường minh, nên lỗi tuần tự hóa 40001 không phát sinh và \texttt{TxManager} không thử lại mã này. Ở PostgreSQL, 40001 chỉ xuất hiện khi giao dịch \texttt{REPEATABLE READ} hoặc \texttt{SERIALIZABLE} sửa một dòng đã bị giao dịch khác sửa sau thời điểm chụp dữ liệu, hoặc khi \texttt{SERIALIZABLE} phát hiện phụ thuộc đọc ghi nguy hiểm~\cite{postgresql}. Các thao tác chỉ đọc cần ảnh chụp nhất quán qua nhiều câu truy vấn, như bản xem trước của S15 hay báo cáo, chạy \texttt{REPEATABLE READ READ ONLY}, vốn không bao giờ gặp 40001. Nhóm không nâng mức cô lập cho giao dịch ghi vì mọi giao dịch ghi đều cập nhật \texttt{spaces.change\_seq}: ở \texttt{REPEATABLE READ}, hai lệnh ghi đồng thời trong cùng không gian sẽ thường xuyên va nhau tại đúng dòng này.

% ---------------------------------------------------------------------
\subsection{Máy trạng thái}\label{ssec:may-trang-thai}

Với trạng thái được lưu, nhóm đặt hai lớp bảo vệ: phương thức của thực thể từ chối chuyển trạng thái không hợp lệ, và câu lệnh cập nhật có điều kiện \texttt{WHERE status IN (...)}, kèm \texttt{AND version = \$v} khi lệnh mang \texttt{If-Match}, giữ đúng máy trạng thái khi hai yêu cầu đồng thời cùng đọc trạng thái cũ. Nhãn chuyển trạng thái có dạng \textit{API [điều kiện] / hành động}; dòng bên trong một trạng thái là lệnh không đổi trạng thái kèm phản hồi, ví dụ \textit{A04 / 423 (3004)} nghĩa là A04 trả HTTP 423 với mã 3004. Mã API theo mục~\ref{sec:api}. Hình~\ref{fig:trang-thai-tai-khoan} thể hiện vòng đời tài khoản.

\diagram[0.8\textwidth]{c3-11-trang-thai-tai-khoan}{Biểu đồ trạng thái tài khoản người dùng}{fig:trang-thai-tai-khoan}

Tài khoản chỉ sang \texttt{active} khi OTP hợp lệ; trong cùng giao dịch, \texttt{household} tạo không gian cá nhân và dòng trưởng nhóm qua \texttt{ActivationHook} (\br{BR04}). Ở \texttt{pending}, A04 với mật khẩu đúng trả 403/3009 (\msg{MSG34}) và máy chủ gửi lại OTP nếu còn trong giới hạn của \br{BR02}, để người dùng quên bước xác thực được dẫn thẳng sang màn hình nhập mã. Lần đăng nhập sai thứ năm đặt \texttt{locked\_until} sau 15 phút (\br{BR03}); trong thời gian đó A04 trả 423/3004 (\msg{MSG06}) kèm \texttt{Retry-After}. Việc mở khóa được đánh giá lười ở lần đăng nhập kế tiếp nên không cần tác vụ định kỳ; A08 với \texttt{resetToken} hợp lệ cũng đưa bộ đếm về 0, xóa \texttt{locked\_until}, thu hồi mọi phiên và tăng \texttt{token\_version}. Trạng thái \texttt{disabled} do quản trị viên đặt qua D02 (\uc{UC48}) kèm \texttt{lock\_reason} và không tự hết hạn: phiên bị thu hồi, \texttt{token\_version} tăng nên mã truy cập còn hạn mất hiệu lực ngay ở bước 7 của Hình~\ref{fig:pipeline-yeu-cau}, đăng nhập trả 403/3005 (\msg{MSG21}). Đổi mật khẩu bằng A11 không đổi trạng thái nhưng thu hồi các phiên khác, tăng \texttt{token\_version} và cấp cặp mã mới cho phiên hiện tại.

Xóa tài khoản bằng A16 (\nfr{NFR18}) đòi nhập lại mật khẩu, sai thì trả 400/2002, và phải thỏa \br{BR07}: người đang là trưởng nhóm của một nhóm còn thành viên khác nhận 422/4004 (\msg{MSG24}) và phải chuyển quyền trước. Khi hợp lệ, một giao dịch đặt \texttt{deleted\_at}, thu hồi mọi phiên, xóa thiết bị và gỡ tư cách thành viên nhóm gia đình như khi rời nhóm, nên mặt hàng đang giao cho người đó trở về Cần mua qua \texttt{MemberLeaveHook}. Từ đó A04 trả 401/3003 như với email không tồn tại, không tiết lộ tài khoản từng có. Sau 30 ngày, tác vụ bảo trì ẩn danh email và họ tên, xóa ảnh, dữ liệu không gian cá nhân và tài liệu Elasticsearch; dòng người dùng được giữ ở dạng ẩn danh để dữ liệu đóng góp trong nhóm hiển thị “Người dùng đã xóa” mà không gãy khóa ngoại.

Hình~\ref{fig:trang-thai-mat-hang-be} là phiên bản phía máy chủ của Hình~\ref{fig:trang-thai-mat-hang}, gắn mỗi chuyển trạng thái với lệnh API và điều kiện bảo vệ. Hai trạng thái Cần mua và Đã phân công được gom vào trạng thái hợp “Chưa mua” vì chúng chung ba lối ra: đánh dấu đã mua, bỏ qua và xóa mềm. Bảng~\ref{tab:c3-chuyen-trang-thai} bổ sung sự kiện miền hoặc tác vụ nền phát ra ở mỗi chuyển trạng thái.

{\small
\begin{longtable}{|L{3.0cm}|L{1.7cm}|L{\dimexpr\textwidth-8\tabcolsep-5\arrayrulewidth-3.0cm-1.7cm-3.8cm\relax}|L{3.8cm}|}
\caption{Chuyển trạng thái của mặt hàng mua sắm}\label{tab:c3-chuyen-trang-thai}\\
\hline
\hd{Chuyển trạng thái} & \hd{API} & \hd{Điều kiện} & \hd{Sự kiện, tác vụ nền}\\ \hline
\endfirsthead
\multicolumn{4}{l}{\footnotesize\itshape Bảng \thetable\ (tiếp theo)}\\ \hline
\hd{Chuyển trạng thái} & \hd{API} & \hd{Điều kiện} & \hd{Sự kiện, tác vụ nền}\\ \hline
\endhead
(mới) → Cần mua & S07, S06, S15, R08 & Danh sách chưa Hoàn tất; gộp trùng (\br{BR09}) & Không\\ \hline
Cần mua → Đã phân công & S10 & Trưởng nhóm; version khớp (\br{BR10}) & \texttt{AssignmentDigest} lùi 5~s, phát \texttt{item.assigned}\\ \hline
Đã phân công → Cần mua & S11; G08, G10 & Trưởng nhóm; hoặc người được giao rời, bị xóa khỏi nhóm (\br{BR07}, \texttt{MemberLeaveHook}) & \texttt{item.unassigned}\\ \hline
Chưa mua → Đã mua & S12 & Lệnh được máy chủ ghi nhận trước thắng; mua thay cần \texttt{onBehalf} (\br{BR10}) & \texttt{items.purchased}; \texttt{PushNotification} khi mua thay\\ \hline
Đã mua → Cần mua & S13 & Chưa nhập tủ; version khớp; danh sách Hoàn tất trở về Đang mua & Không\\ \hline
Đã mua → Đã nhập tủ & S14 & \texttt{lot\_id IS NULL}; có \texttt{food\_id}; hạn mặc định theo \br{BR11} & \texttt{items.stocked}, \texttt{lot.changed}\\ \hline
Chưa mua → Đã bỏ qua & S08 & \texttt{status = skipped}; version khớp & \texttt{item.unassigned} nếu đang giao\\ \hline
Chưa mua → kết thúc & S09 & Version khớp; xóa mềm, để lại tombstone cho đồng bộ & Không\\ \hline
Đã mua → kết thúc & -- & Danh sách Hoàn tất mà mặt hàng không được nhập tủ & Không\\ \hline
\end{longtable}
}

\diagram[\textwidth]{c3-12-trang-thai-mat-hang-be}{Biểu đồ trạng thái mặt hàng mua sắm ở phía máy chủ}{fig:trang-thai-mat-hang-be}

Điều kiện \textit{[leader]} lấy từ vai trò do bước 8 của Hình~\ref{fig:pipeline-yeu-cau} đặt vào context; \textit{[version khớp]} so header \texttt{If-Match} với cột \texttt{version}, lệch thì trả 409/4002. S12 xử lý cả lô mặt hàng bằng một câu \texttt{UPDATE ... WHERE status IN ('to\_buy', 'assigned') RETURNING}, hiện thực quy tắc “lệnh được máy chủ ghi nhận trước thắng”. Đây là cách hiện thực “giữ bản ghi sớm hơn” của \uc{UC20} (\nfr{NFR11}): đồng hồ thiết bị không đáng tin, nên \texttt{purchasedAt} do ứng dụng gửi chỉ dùng để hiển thị và báo cáo, không dùng để phân xử. Mặt hàng không được trả về đã có người mua trước và nhận kết quả \texttt{already\_purchased} kèm người mua trong phản hồi 200; nếu không mặt hàng nào của lệnh được ghi, như khi phát lại một mặt hàng duy nhất, phản hồi là 409/4002 kèm \texttt{current.purchasedBy}. S14 thêm điều kiện \texttt{lot\_id IS NULL} và tạo lô qua \texttt{StockIn.AddLots} trong cùng giao dịch, nên một mặt hàng không thể sinh hai lô; gửi lại nhận \texttt{already\_stocked} kèm lô cũ. Mặt hàng ghi tên tự do không có \texttt{food\_id} nên không bao giờ sang Đã nhập tủ: S14 trả 422/4004 cho mặt hàng đó, và khi danh sách Hoàn tất nó kết thúc ở Đã mua với chi tiêu đã được ghi nhận. Trạng thái danh sách (\br{BR08}) được ghi cùng giao dịch; danh sách Hoàn tất từ chối lệnh sửa nội dung bằng 422/4004, trừ ba lệnh sau khi mua: S08 sửa \texttt{actualQuantity}, \texttt{unitPrice} của mặt hàng đã mua, S13 hoàn tác (danh sách trở về Đang mua) và S14 nhập tủ, vì người dùng thường cất hàng sau khi danh sách đã hoàn tất. Phân công gộp qua tác vụ \texttt{AssignmentDigest} của \texttt{shopping}, duy nhất theo cặp người được giao và danh sách, chạy sau 5~giây, nên giao liền nhiều mặt hàng chỉ sinh một thông báo.

Với lô thực phẩm, cơ sở dữ liệu chỉ lưu \texttt{lifecycle} (\texttt{active}, \texttt{consumed}, \texttt{discarded}); ba trạng thái hạn dùng ở Hình~\ref{fig:trang-thai-thuc-pham} được suy diễn lúc truy vấn theo ngưỡng $N$ của người xem (\br{BR12}). Lý do là $N$ thuộc thiết lập từng người (\uc{UC08}) nên một cột lưu sẵn không thể đúng cho mọi thành viên; trạng thái tự đổi theo ngày mà không cần tác vụ cập nhật hằng đêm, bước tính lại trạng thái của \uc{UC31} trở thành phép tính khi đọc; và dữ liệu không bao giờ lệch giữa các thành viên hay khi tác vụ nền chạy trễ. Mã nguồn~\ref{code:c3-trang-thai-han-dung} là biểu thức tính, với \texttt{\$2} là tham số \texttt{\$today}, tức ngày hôm nay theo giờ Việt Nam do \path{platform/clock} cung cấp (không dùng \texttt{current\_date} của phiên chạy UTC), và \texttt{\$3} là \texttt{reminder\_days} của người xem hoặc giá trị mặc định \texttt{reminder.defaultDays}; \texttt{FoodLot.ExpiryState} hiện thực cùng quy tắc trong Go và được kiểm thử đối chiếu.

\codecaption{Biểu thức SQL tính trạng thái hạn dùng suy diễn}{code:c3-trang-thai-han-dung}
\begin{Verbatim}
CASE WHEN l.expires_on < $2 THEN 'expired'
     WHEN l.expires_on - $2 <= $3 THEN 'expiring'
     ELSE 'fresh' END AS expiry_state
\end{Verbatim}

Cuối cùng, mọi tác vụ nền đi qua vòng đời của River~\cite{river} ở Hình~\ref{fig:trang-thai-tac-vu}, lưu bền vững trong bảng \texttt{river\_job}. Ba trạng thái chưa chạy được gom thành trạng thái hợp “Chờ chạy” vì cùng có thể bị hủy thủ công.

\diagram[\textwidth]{c3-13-trang-thai-tac-vu}{Vòng đời tác vụ nền trong hàng đợi River}{fig:trang-thai-tac-vu}

Tác vụ chèn bằng \texttt{InsertTx} chỉ tồn tại khi dữ liệu nghiệp vụ đã commit. Tác vụ khai báo khóa duy nhất (\texttt{UniqueOpts}) như \texttt{UserReminder} theo người dùng và ngày hay \texttt{AssignmentDigest} theo người được giao và danh sách: lần chèn trùng bị bỏ qua, nên chống trùng được giải ngay ở hàng đợi. Worker nhận việc bằng \texttt{FOR UPDATE SKIP LOCKED} nên không nhận trùng; tiến trình chết giữa chừng thì bộ cứu hộ (rescuer) đưa tác vụ treo về \texttt{retryable}. Lỗi vĩnh viễn, như dữ liệu đích đã bị xóa, được worker báo bằng \texttt{river.JobCancel} để không thử lại vô ích; người vận hành có thể hủy tác vụ chưa chạy bằng \texttt{client.JobCancel} và chạy lại tác vụ \texttt{discarded} bằng \texttt{client.JobRetry} sau khi khắc phục nguyên nhân, ví dụ khóa dịch vụ FCM hết hạn. Tác vụ \texttt{PushNotification} trên hàng đợi \texttt{notify} có tối đa 4 lần chạy, tức ba lần thử lại theo \nfr{NFR12}, chờ 20~s, 60~s, 180~s cộng nhiễu; độ trễ gửi xấu nhất khoảng 4,5 phút. Con số này thuộc chỉ số “gửi tới FCM”; chỉ số “thông báo được tạo trong 15 phút sau giờ nhắc” không phụ thuộc vào nó, vì bản ghi thông báo được \texttt{UserReminder} tạo trước khi gửi đẩy. Tác vụ \texttt{discarded} được đếm bằng chỉ số Prometheus để cảnh báo, còn thông báo vẫn nằm trong hộp thư của người dùng. Cơ chế chống nhắc trùng được trình bày ở mục~\ref{sssec:kt-nhac-han}.

% ---------------------------------------------------------------------
%  3.4.5. Các kỹ thuật nổi bật trong thiết kế (nối tiếp mục 3.4)
% ---------------------------------------------------------------------
\subsection{Các kỹ thuật nổi bật trong thiết kế}\label{ssec:ky-thuat}

Phần này trình bày tám kỹ thuật cho những tình huống dễ làm hệ thống sai lệch: ghi đồng thời, mất mạng, tiến trình dừng giữa chừng và dịch vụ ngoài gặp sự cố. Mỗi kỹ thuật đi theo trình tự vấn đề, giải pháp và đánh đổi, dựa trên ba nguyên tắc chung: PostgreSQL là nguồn dữ liệu gốc duy nhất; Redis, Elasticsearch, kênh WebSocket và thông báo đẩy được phép mất dữ liệu tạm thời mà không làm mất dữ liệu nghiệp vụ; mọi thao tác có thể bị lặp lại đều phải lũy đẳng (idempotent).

% ---------------------------------------------------------------------
\subsubsection{Giao dịch nguyên tử và khóa dòng khi trừ kho theo FEFO}\label{sssec:kt-fefo}

\uc{UC41} và \uc{UC28} trừ nhiều lô trong một lần thao tác, trong khi nhiều thành viên dùng chung một tủ lạnh. Nếu dịch vụ đọc lượng còn lại rồi ghi giá trị mới, hai giao dịch có thể cùng đọc 0,5 kg và cùng trừ 0,4 kg, làm mất một lần cập nhật (lost update) và vi phạm \br{BR14}. Nhóm dùng một giao dịch duy nhất ở mức cô lập READ COMMITTED kết hợp khóa dòng tường minh \texttt{SELECT ... FOR UPDATE}~\cite{postgresql}, đáp ứng yêu cầu cập nhật nguyên tử của \nfr{NFR11} và \nfr{NFR13}. Hình~\ref{fig:seq-xac-nhan-nau} thể hiện trình tự xác nhận đã nấu theo đúng ranh giới module: mealplan điều phối, pantry sở hữu dữ liệu tồn kho.

\diagram[\textwidth]{c3-14-seq-xac-nhan-nau}{Biểu đồ tuần tự xác nhận đã nấu với giao dịch nguyên tử và khóa dòng}{fig:seq-xac-nhan-nau}

\texttt{MealHandler} chỉ giải mã thân yêu cầu \texttt{\{servings, lines[], skipStock\}}, trong đó \texttt{lines[]} là lượng người dùng đã xác nhận ở bước 5 của \uc{UC41}, có thể kèm lô chọn đích danh, rồi gọi \texttt{MealPlanService.Cook(ctx, CookCmd)}. Câu đầu tiên của giao dịch là chèn khóa lũy đẳng (bước 3, xem mục~\ref{sssec:kt-ngoai-tuyen}); món ăn bị khóa ngay sau đó (bước 4), nên yêu cầu lặp lại hoặc yêu cầu từ thiết bị khác thấy trạng thái \texttt{cooked} và nhận kết quả cũ mà không trừ kho lần hai (bước 6b của use case, bước 5--7 của hình). Khi người dùng chọn “Bỏ qua trừ kho” (bước 5a), khối \texttt{opt} bị bỏ qua và món chỉ đổi trạng thái. Ngược lại, \texttt{MealPlanService} lấy định lượng qua cổng \texttt{RecipeQuery.Ingredients}, nhân theo tỷ lệ khẩu phần, điều chỉnh theo \texttt{lines[]}, sắp theo \texttt{food\_id} rồi giao cho cổng \texttt{StockConsumer.ConsumeForMeal}. Module mealplan không đọc hay ghi bảng nào của pantry: \texttt{PantryService} hiện thực cổng này và chạy trong chính \texttt{pgx.Tx} mà \texttt{TxManager} truyền qua \texttt{context}. \texttt{FEFOAllocator} là hàm thuần của tầng miền pantry: quy đổi lượng về đơn vị cơ sở (\br{BR16}), rồi trừ dần từ lô hết hạn sớm nhất theo \br{BR15}, sau khi sắp lại thứ tự trong bộ nhớ vì \texttt{ORDER BY} kèm \texttt{FOR UPDATE} có thể trả thứ tự lệch khi cột sắp xếp vừa bị sửa đồng thời. Lô quá hạn không bao giờ lọt vào truy vấn ở bước 10, trừ lô người dùng chọn đích danh (bước 4b). Bảng xem trước ở bước 4 của use case được ứng dụng tính từ bản sao cục bộ nên có thể đã cũ; nếu lượng còn lại sau khi khóa nhỏ hơn lượng đã xác nhận, cả giao dịch bị hủy và trả 409 mã 4002 kèm tồn kho mới để ứng dụng báo \msg{MSG35} (bước 6a). Sau commit, API xóa bộ đệm gợi ý \texttt{suggest:\{sid\}} và phát gợi ý thời gian thực (bước 21); mỗi lô bị trừ còn có một dòng \texttt{stock\_movements} gắn với món ăn để phục vụ \uc{UC45}.

Mọi phép so sánh ngày nghiệp vụ dùng tham số \texttt{\$today}, là ngày theo múi giờ Asia/Ho\_Chi\_Minh do gói \texttt{platform/clock} cấp, thay cho hàm \texttt{current\_date}. Phiên CSDL chạy theo UTC, nên trong khoảng 00:00--07:00 giờ Việt Nam \texttt{current\_date} vẫn là ngày hôm trước và lô vừa hết hạn hôm qua sẽ bị coi là còn hạn, vi phạm \br{BR15}; tham số hóa ngày còn cho phép kiểm thử tất định bằng đồng hồ giả. Mã nguồn~\ref{code:c3-sql-fefo} trình bày các câu lệnh cốt lõi do sqlc sinh thành hàm Go.

\codecaption{Các câu lệnh SQL cốt lõi của giao dịch trừ kho theo FEFO}{code:c3-sql-fefo}
\begin{Verbatim}
-- $3 = $today: ngày Asia/Ho_Chi_Minh do platform/clock cấp
-- Mỗi nguyên liệu, theo food_id tăng dần: khóa các lô còn hạn (BR15)
SELECT id, remaining_qty, unit_factor, expires_on
  FROM pantry.food_lots
 WHERE space_id = $1 AND food_id = $2 AND lifecycle = 'active'
   AND deleted_at IS NULL AND remaining_qty > 0
   AND expires_on >= $3          -- lô quá hạn không trừ tự động
 ORDER BY expires_on, id FOR UPDATE;

-- Trừ từng lô theo kết quả phân bổ; về 0 thì chuyển consumed (BR14)
UPDATE pantry.food_lots
   SET remaining_qty = remaining_qty - $2, version = version + 1,
       lifecycle = CASE WHEN remaining_qty = $2
                        THEN 'consumed' ELSE lifecycle END
 WHERE id = $1 AND remaining_qty >= $2;

-- Nhật ký tiêu thụ gắn với món ăn (UC45); bảng chỉ ghi thêm
INSERT INTO pantry.stock_movements (id, lot_id, space_id, kind,
       quantity, unit_id, meal_entry_id, actor_id, occurred_at)
VALUES ($1, $2, $3, 'cook', $4, $5, $6, $7, now());

-- Cấp số thứ tự thay đổi: luôn là câu khóa cuối cùng trước COMMIT
UPDATE household.spaces SET change_seq = change_seq + 1
 WHERE id = $1 RETURNING change_seq;
\end{Verbatim}

READ COMMITTED đủ an toàn vì khi \texttt{FOR UPDATE} phải chờ, PostgreSQL đánh giá lại điều kiện \texttt{WHERE} trên phiên bản mới nhất của dòng, nên lô vừa bị trừ về 0 tự rơi khỏi kết quả; SERIALIZABLE chỉ phát hiện xung đột lúc commit và buộc chạy lại cả giao dịch. Rủi ro còn lại là bế tắc (deadlock) khi hai giao dịch khóa cùng tập dòng theo thứ tự ngược nhau. Bảng~\ref{tab:c3-thu-tu-khoa} quy định thứ tự khóa toàn cục mà mọi giao dịch ghi trong hệ thống phải tuân theo, không riêng giao dịch trừ kho.

{\small
\begin{longtable}{|C{1.0cm}|L{2.7cm}|L{6.0cm}|L{\dimexpr\textwidth-8\tabcolsep-5\arrayrulewidth-1.0cm-2.7cm-6.0cm\relax}|}
\caption{Thứ tự khóa toàn cục trong giao dịch ghi}\label{tab:c3-thu-tu-khoa}\\
\hline
\hd{Bậc} & \hd{Đối tượng khóa} & \hd{Cách lấy khóa, thứ tự trong bậc} & \hd{Thao tác áp dụng}\\ \hline
\endfirsthead
\multicolumn{4}{l}{\footnotesize\itshape Bảng \thetable\ (tiếp theo)}\\ \hline
\hd{Bậc} & \hd{Đối tượng khóa} & \hd{Cách lấy khóa, thứ tự trong bậc} & \hd{Thao tác áp dụng}\\ \hline
\endhead
1 & Khóa lũy đẳng & \texttt{INSERT INTO sync.idempotency\_keys ... ON CONFLICT DO NOTHING}, câu đầu tiên của giao dịch & Mọi lệnh ghi có \texttt{Idempotency-Key}\\ \hline
2 & Khóa tư vấn theo không gian & \texttt{pg\_advisory\_xact\_lock(}\linebreak\texttt{hashtext(space\_id::text))}, tự nhả khi giao dịch kết thúc & G05--G10: mã mời, xóa thành viên, chuyển quyền, rời nhóm; G07 tra mã mời trước để biết \texttt{space\_id}\\ \hline
3 & Dòng thành viên & \texttt{space\_members} theo \texttt{user\_id} tăng dần; \texttt{FOR SHARE} khi chỉ kiểm vai trò, \texttt{FOR UPDATE} khi đổi vai trò hoặc xóa & G08--G10, S10, S11\\ \hline
4 & Thực thể nghiệp vụ & Danh sách trước mặt hàng, công thức trước món ăn; cùng loại theo \texttt{id} tăng dần; \texttt{SELECT ... FOR UPDATE} hoặc \texttt{UPDATE ... WHERE id} & S04--S14, R04, R05, M02--M04, M06\\ \hline
5 & Lô thực phẩm & Theo \texttt{(food\_id, expires\_on, id)}, \texttt{FOR UPDATE} & P04--P06, M06; S14 chỉ chèn lô mới\\ \hline
6 & Bộ đếm thay đổi & \texttt{UPDATE household.spaces SET change\_seq = change\_seq + 1 ... RETURNING}, câu khóa cuối cùng trước \texttt{COMMIT} & Mọi giao dịch ghi thực thể đồng bộ\\ \hline
\end{longtable}
}

Mọi giao dịch chỉ xin khóa theo chiều tăng của cột Bậc, và trong cùng bậc theo thứ tự ghi ở cột thứ ba, nên không thể hình thành chu trình chờ: một chu trình đòi hỏi ít nhất một giao dịch đang giữ khóa đứng sau lại chờ khóa đứng trước, điều mà thứ tự toàn phần này loại trừ. Dòng \texttt{household.spaces} chỉ bị khóa ở bậc 6, sau mọi khóa khác; mọi giao dịch ghi của không gian đều cần dòng này, nhưng không giao dịch nào giữ nó trong lúc chờ khóa khác, nên nó không thể nằm trong chu trình. Thời gian giữ ngắn chỉ giảm tranh chấp chứ không phải lý do loại trừ bế tắc, và đó cũng là lý do thiết kế không khóa dòng không gian \texttt{FOR UPDATE} ở đầu giao dịch. Các thao tác đổi thành viên, vai trò và mã mời cần tuần tự hóa theo không gian suốt giao dịch nên dùng khóa tư vấn ở bậc 2; G07 tra mã mời không khóa để lấy \texttt{space\_id} rồi mới xin khóa này. Hai nguồn khóa ngầm cũng đã được xét: kiểm tra khóa ngoại lấy \texttt{FOR KEY SHARE} trên dòng được tham chiếu, không xung đột với \texttt{FOR NO KEY UPDATE} của câu cấp \texttt{change\_seq} vì cột này không thuộc khóa; chèn trùng chỉ mục duy nhất phải chờ giao dịch đang chèn cùng khóa, điều chỉ xảy ra ở bậc 1. Dù vậy, \texttt{TxManager} vẫn đặt \texttt{lock\_timeout} 3 giây và tự chạy lại tối đa ba lần khi gặp mã lỗi 40P01 (\texttt{deadlock\_detected}) hoặc 55P03 (\texttt{lock\_not\_available}, khi vượt \texttt{lock\_timeout}) để phòng lỗi lập trình làm sai thứ tự; ở READ COMMITTED không phát sinh 40001 nên mã này không nằm trong danh sách thử lại. Chạy lại an toàn vì giao dịch bị hủy không để lại dấu vết, kể cả khóa lũy đẳng; hết ba lần, API trả 503 mã 5003 (\msg{MSG33}). Đánh đổi là các thao tác ghi trong cùng một không gian bị tuần tự hóa trong thời gian rất ngắn, không đáng kể với nhóm tối đa 10 người (\br{BR05}).

% ---------------------------------------------------------------------
\subsubsection{Hộp thư đi giao dịch (transactional outbox) với River}\label{sssec:kt-outbox}

Khi một thành viên mua thay người được giao (bước 2a của \uc{UC20}), hệ thống phải vừa cập nhật CSDL vừa gửi thông báo đẩy. Gửi sau commit thì tiến trình có thể dừng giữa chừng làm mất thông báo, gửi trước commit thì giao dịch có thể bị hủy sau khi đã báo: bài toán ghi kép (dual write)~\cite{kleppmann}. Theo mẫu hộp thư đi giao dịch~\cite{richardson}, giao dịch nghiệp vụ chỉ ghi “ý định gửi” vào CSDL, một tiến trình khác thực hiện việc gửi; với River~\cite{river}, bảng \texttt{river\_job} là hộp thư đi và worker là bộ chuyển tiếp, như Hình~\ref{fig:seq-outbox-mua-hang}.

\diagram[\textwidth]{c3-15-seq-outbox-mua-hang}{Biểu đồ tuần tự đánh dấu đã mua với hộp thư đi giao dịch}{fig:seq-outbox-mua-hang}

Sau khóa lũy đẳng, \texttt{ShoppingService.MarkPurchased} khóa danh sách trước rồi lần lượt từng mặt hàng theo \texttt{id}, đúng bậc 4 của Bảng~\ref{tab:c3-thu-tu-khoa}. Mỗi mặt hàng được chuyển bằng một câu \texttt{UPDATE} có điều kiện \texttt{status IN ('to\_buy', 'assigned')} kèm \texttt{RETURNING} (bước 3): khi hai thiết bị cùng đánh dấu, câu của giao dịch đến sau phải chờ, đánh giá lại điều kiện trên dòng mới và không cập nhật được dòng nào. Đây là cách hiện thực quy tắc “giữ bản ghi sớm hơn” ở bước 6a của \uc{UC20}: lệnh được máy chủ ghi nhận trước thắng, vì đồng hồ thiết bị không đáng tin, còn \texttt{purchasedAt} do ứng dụng gửi chỉ dùng để hiển thị và báo cáo. Với yêu cầu chỉ có một mặt hàng, trường hợp của thao tác chạm ô đánh dấu và của phát lại ngoại tuyến, lệnh thua nhận 409 mã 4002 kèm \texttt{current.purchasedBy} để ứng dụng báo \msg{MSG35} (bước 4--5); với yêu cầu nhiều mặt hàng, mặt hàng đó nhận \texttt{outcome = already\_purchased} kèm \texttt{current} trong phản hồi 200 và các mặt hàng khác vẫn được ghi. Mặt hàng đang giao cho người khác mà thiếu cờ \texttt{onBehalf} lại là lỗi dữ liệu vào: cả yêu cầu bị hủy với 422 mã 4004 (\br{BR10}, \msg{MSG12}) để ứng dụng hỏi xác nhận mua thay (bước 6--7), vì mọi lỗi nhóm 2xxx và 4004 đều theo nguyên tắc tất cả hoặc không. Trạng thái danh sách theo \br{BR08} được tính lại trong cùng giao dịch (bước 9).

\texttt{InsertTx} ở bước 8 dùng chính \texttt{pgx.Tx} của giao dịch nghiệp vụ, nên tác vụ \texttt{PushNotification} tồn tại khi và chỉ khi mặt hàng đã sang \texttt{purchased}; \texttt{PUBLISH} nằm ngoài giao dịch vì chỉ là gợi ý thời gian thực (mục~\ref{sssec:kt-thoi-gian-thuc}). \texttt{NotifyWorker} lấy tác vụ bằng \texttt{SKIP LOCKED} nên nhiều goroutine cùng quét hàng đợi \texttt{notify} mà không lấy trùng; nó đọc \texttt{assignment\_notify\_on} của người nhận, nếu người này tắt loại thông báo thì tác vụ hoàn tất mà không gửi, rồi đọc danh sách thiết bị (bước 12). Đánh đổi của mẫu là giao nhận \emph{ít nhất một lần} (at-least-once): worker có thể gửi xong rồi dừng trước khi đánh dấu \texttt{completed}, và River chạy lại tác vụ theo vòng đời ở Hình~\ref{fig:trang-thai-tac-vu}. Vì vậy \texttt{dedupe\_key} duy nhất biến lần chèn thông báo thứ hai thành thao tác rỗng (bước 13), bảng \texttt{deliveries} chỉ cho gửi lại thiết bị chưa \texttt{sent} (bước 14--15), và mã thông báo được đặt vào khóa gộp (\texttt{collapse\_key} và \texttt{tag} trên Android, \texttt{apns-collapse-id} trên iOS) để bản gửi trùng thay thế bản cũ trên khay thông báo. Đổi lại, hệ thống không phải vận hành thêm message broker.

Nhập tủ (\uc{UC21}) dùng cùng khuôn giao dịch nhưng đi qua cổng \texttt{StockIn} của module pantry, như Hình~\ref{fig:seq-nhap-tu}.

\diagram[0.85\textwidth]{c3-15b-seq-nhap-tu}{Biểu đồ tuần tự chuyển mặt hàng đã mua vào tủ lạnh qua cổng StockIn}{fig:seq-nhap-tu}

\texttt{ShoppingService.StockItems} không ghi bảng \texttt{food\_lots}: với mỗi mặt hàng đã khóa, nó gọi \texttt{StockIn.AddLots} (bước 7) với \texttt{source\_item\_id}, ngày hết hạn người dùng xác nhận hoặc mặc định bằng ngày mua cộng hạn dùng gợi ý tại vị trí lưu trữ (\br{BR11}), và \texttt{total\_cost} bằng lượng thực tế nhân đơn giá. Lưu tổng tiền thay cho đơn giá giúp giá trị lãng phí ở \uc{UC46} vẫn đúng khi người dùng đổi đơn vị lúc nhập tủ. \texttt{PantryService} chèn lô và biến động \texttt{stock\_in} (bước 8), rồi mặt hàng chuyển \texttt{stocked} kèm \texttt{lot\_id} (bước 10); mọi việc nằm trong một giao dịch nên lỗi giữa chừng hoàn tác cả việc tạo lô lẫn đổi trạng thái. Mặt hàng đã \texttt{stocked} không phải lỗi mà là kết quả \texttt{already\_stocked} kèm \texttt{lotId} (bước 6), nên yêu cầu lặp lại hay thiết bị thứ hai chỉ nhận lại lô đã có (bước 5b của \uc{UC21}); mặt hàng tên tự do chưa gắn thực phẩm là lỗi dữ liệu vào, hủy cả yêu cầu với 422 mã 4004 để ứng dụng yêu cầu chọn thực phẩm (bước 2a). \br{BR08} chặn sửa nội dung danh sách Hoàn tất (422, mã 4004) với ba ngoại lệ có chủ đích: S08 sửa lượng thực tế và đơn giá của mặt hàng đã mua, S13 hoàn tác (danh sách trở về Đang mua) và S14 nhập tủ, vì người dùng thường nhập giá và cất đồ sau khi đã mua xong cả danh sách.

% ---------------------------------------------------------------------
\subsubsection{Đồng bộ ngoại tuyến và giải quyết xung đột}\label{sssec:kt-ngoai-tuyen}

\nfr{NFR11} yêu cầu bốn nhóm thao tác dùng được khi mất mạng mà không mất hay ghi trùng dữ liệu. Ứng dụng theo hướng ưu tiên ngoại tuyến (offline-first): mỗi thao tác ghi cập nhật bản sao SQLite (drift) và chèn vào hàng đợi \texttt{pending\_ops} trong cùng một giao dịch cục bộ, kèm \texttt{Idempotency-Key}, phiên bản gốc, ảnh chụp các trường trước khi sửa và thời điểm thao tác, đồng thời hiển thị \msg{MSG20}. Khi có mạng, bộ đồng bộ của ứng dụng phát lại hàng đợi bằng đúng các API REST thông thường, nên quy tắc nghiệp vụ và phân quyền chỉ cài đặt một nơi (Hình~\ref{fig:seq-dong-bo-ngoai-tuyen}); phía máy chủ không có thành phần đồng bộ riêng ngoài nhật ký thay đổi và bảng khóa lũy đẳng.

\diagram[\textwidth]{c3-16-seq-dong-bo-ngoai-tuyen}{Biểu đồ tuần tự đồng bộ ngoại tuyến với phát lại lũy đẳng}{fig:seq-dong-bo-ngoai-tuyen}

Câu đầu tiên của mọi giao dịch ghi có \texttt{Idempotency-Key} chèn \texttt{(user\_id, key)} cùng băm SHA-256 của yêu cầu vào \path{sync.idempotency_keys} với \texttt{status = 'in\_progress'} và \texttt{ON CONFLICT DO NOTHING}~\cite{idempotencykey} (bước 5). Nếu không chèn được dòng nào thì khóa đã tồn tại: cùng băm và đã \texttt{completed}, máy chủ phát lại phản hồi đã lưu kèm header \texttt{Idempotent-Replayed: true} (bước 6); khác băm nghĩa là ứng dụng dùng lại khóa cho một yêu cầu khác và bị từ chối với 400 mã 2002 (bước 7). Ngay trước \texttt{COMMIT}, dòng khóa được cập nhật thành \texttt{completed} cùng mã trạng thái và thân phản hồi; khi có lỗi, \texttt{ROLLBACK} xóa luôn dòng khóa, nên bảng chỉ lưu kết quả thành công và một lệnh bị từ chối có thể gửi lại sau khi sửa. Hai bản phát lại đồng thời tự xếp hàng ở chỉ mục khóa chính: bản sau chờ bản đầu kết thúc rồi thấy dòng \texttt{completed}, hoặc chèn được nếu bản đầu đã rollback. Vì dòng \texttt{in\_progress} chưa commit là vô hình với giao dịch khác, bản sau không đọc được trạng thái này mà chỉ có thể chờ quá \texttt{lock\_timeout}; khi đó máy chủ trả 409 mã 4002 để ứng dụng thử lại sau. Các API công khai A01--A05, A07, A08 không dùng khóa lũy đẳng vì đã được bảo vệ bằng giới hạn tần suất và OTP. Xung đột nghiệp vụ được xử lý theo Bảng~\ref{tab:c3-xung-dot}.

{\small
\begin{longtable}{|L{2.4cm}|L{4.3cm}|L{3.3cm}|L{\dimexpr\textwidth-8\tabcolsep-5\arrayrulewidth-2.4cm-4.3cm-3.3cm\relax}|}
\caption{Quy tắc giải quyết xung đột khi đồng bộ}\label{tab:c3-xung-dot}\\
\hline
\hd{Loại thao tác} & \hd{Quy tắc giải quyết} & \hd{Phản hồi máy chủ} & \hd{Ứng dụng xử lý}\\ \hline
\endfirsthead
\multicolumn{4}{l}{\footnotesize\itshape Bảng \thetable\ (tiếp theo)}\\ \hline
\hd{Loại thao tác} & \hd{Quy tắc giải quyết} & \hd{Phản hồi máy chủ} & \hd{Ứng dụng xử lý}\\ \hline
\endhead
Tạo danh sách, mặt hàng, lô & Id UUIDv7 do thiết bị sinh; mặt hàng trùng thực phẩm được gộp theo \br{BR09} & 201, hoặc phản hồi đã lưu khi phát lại & Thay bản cục bộ bằng bản máy chủ, báo \msg{MSG11} nếu đã gộp\\ \hline
Đánh dấu đã mua & Lệnh được máy chủ ghi nhận trước thắng, không dựa vào đồng hồ thiết bị; \texttt{purchasedAt} của thiết bị chỉ để hiển thị và báo cáo & 409, mã 4002 (\msg{MSG35}) kèm \texttt{current.}\-\texttt{purchasedBy}; trong yêu cầu nhiều mặt hàng: 200, \texttt{already\_}\-\texttt{purchased} & Áp trạng thái máy chủ, bỏ lệnh cục bộ, báo người đã mua\\ \hline
Sửa thuộc tính (kể cả vị trí, ngày hết hạn, ghi chú của lô) & Hợp nhất theo trường; trường bị sửa ở cả hai phía thì lệnh đến sau thắng (LWW) & 200 kèm \texttt{overwritten}\-\texttt{Fields} & So với ảnh chụp gốc, báo trường đã ghi đè lên thay đổi của người khác\\ \hline
Sửa lượng còn lại của lô (P04) & Không áp LWW: \texttt{If-Match} phải khớp; chênh lệch được ghi thành biến động \texttt{adjust}, lượng bằng 0 thì lô chuyển \texttt{consumed} & 409, mã 4002 (\msg{MSG35}) khi lệch phiên bản; 400, mã 2002 khi vượt lượng ban đầu & Tải lô mới nhất, mở lại ô nhập để người dùng nhập theo tồn kho mới\\ \hline
Sửa mặt hàng đã xóa, đã nhập tủ hoặc thuộc danh sách Hoàn tất & Xóa mềm và trạng thái \texttt{stocked} là trạng thái cuối, luôn thắng; danh sách Hoàn tất chỉ nhận S08 (lượng thực tế, đơn giá), S13 và S14 (\br{BR08}) & 404 mã 4001, hoặc 422 mã 4004 (\msg{MSG32}) & Áp tombstone hoặc trạng thái máy chủ, báo người dùng\\ \hline
Lệnh có quy tắc theo ngày (\br{BR08}, \br{BR11}, \br{BR18}) & Kiểm theo ngày của \texttt{X-Client-Time} nếu thời điểm này trong vòng 7 ngày và không vượt giờ máy chủ; ngoài khoảng đó dùng ngày máy chủ & 422, mã 4004 khi vi phạm & Giữ bản nháp, mở lại biểu mẫu cho người dùng sửa\\ \hline
\end{longtable}
}

Với sửa thuộc tính, \texttt{PATCH} chỉ chứa các trường đã sửa; khi \texttt{If-Match} không khớp, máy chủ vẫn ghi và trả giá trị trước khi ghi của các trường đó, còn ứng dụng so với ảnh chụp gốc để chỉ báo những trường thực sự bị người khác sửa. Lượng còn lại của lô là ngoại lệ có chủ đích: nó bị \uc{UC28}, \uc{UC29}, \uc{UC41} thay đổi liên tục, nên một giá trị tuyệt đối tính từ bản cũ sẽ xóa mất lần tiêu thụ của người khác; P04 vì vậy đòi phiên bản khớp và ghi chênh lệch vào \texttt{stock\_movements} để báo cáo tiêu thụ luôn khớp nhật ký. Lệnh phát lại mang header \texttt{X-Client-Time} là thời điểm thao tác trên thiết bị, để danh sách tạo ngoại tuyến tối qua cho ngày hôm qua vẫn hợp lệ khi đồng bộ sáng nay; giới hạn 7 ngày và điều kiện không vượt giờ máy chủ ngăn việc chỉnh đồng hồ thiết bị để lách quy tắc (bước 10--11). Khi hàng đợi rỗng, ứng dụng kéo \texttt{GET .../changes?since=<seq>} lặp lại đến khi \texttt{hasMore} bằng \texttt{false}, tối đa 500 phần tử mỗi lần, với phản hồi như Mã nguồn~\ref{code:c3-json-changes}; mỗi phần tử mang trạng thái thực thể thay vì thao tác nên áp lại nhiều lần vẫn cho cùng kết quả.

\codecaption{Phản hồi mẫu của API kéo thay đổi (ảnh chụp thực thể đã rút gọn)}{code:c3-json-changes}
\begin{Verbatim}
{ "code": 1000, "message": "Thành công",
  "data": { "nextSince": 1845, "hasMore": false, "changes": [
    { "seq": 1843, "entityType": "shopping_item", "op": "upsert",
      "entityId": "01929f4e-6c1a-7d3e-9b2f-4a5c8e7d1f20", "version": 4,
      "entity": { "status": "purchased", "actualQuantity": 1.5 } },
    { "seq": 1844, "entityType": "food_lot", "op": "upsert",
      "entityId": "01929f4e-7d02-7a11-8c3e-0b6d2f9a4e51", "version": 1,
      "entity": { "remainingQty": 1.5, "expiresOn": "2026-10-28" } },
    { "seq": 1845, "entityType": "shopping_item", "op": "delete",
      "entityId": "01929f4e-80b4-7f62-a1d7-5e3c9b0f2a88" } ] },
  "meta": { "serverTime": "2026-10-25T03:15:03Z" } }
\end{Verbatim}

Số \texttt{seq} là bộ đếm riêng từng không gian, không phải sequence toàn cục. Sequence được cấp lúc chèn nhưng dòng chỉ hiện ra lúc commit: nếu giao dịch T1 nhận số 101, T2 nhận 102 và commit trước, thiết bị kéo đúng lúc đó lưu \texttt{since = 102} và bỏ sót vĩnh viễn thay đổi 101. Bộ đếm \texttt{change\_seq} tăng bằng \texttt{UPDATE ... RETURNING} ở cuối giao dịch (bậc 6 của Bảng~\ref{tab:c3-thu-tu-khoa}) nên T2 phải chờ T1 commit mới được cấp số, và thứ tự số trùng thứ tự commit; bộ đếm phải riêng từng không gian để không tuần tự hóa mọi giao dịch của hệ thống. Con trỏ cũ hơn 90 ngày nhận 410 mã 4100 và ứng dụng tải lại toàn bộ. Khóa chính UUIDv7~\cite{rfc9562} cho phép tạo thực thể khi ngoại tuyến với id không đổi sau đồng bộ, và nhờ tiền tố thời gian, bản ghi mới được chèn vào cuối chỉ mục B-tree; máy chủ chỉ nhận id đúng phiên bản 7 và trả mã 4003 khi trùng. Cái giá là ứng dụng phức tạp hơn và LWW có thể ghi đè một chỉnh sửa có chủ ý, được chấp nhận vì người dùng luôn được báo.

% ---------------------------------------------------------------------
\subsubsection{Đồng bộ thời gian thực đa thiết bị trên Android và iOS}\label{sssec:kt-thoi-gian-thuc}

\nfr{NFR10} yêu cầu thay đổi hiển thị trên thiết bị khác trong 3 giây. Hỏi vòng (polling) tốn pin; iOS đóng kết nối vài giây sau khi ứng dụng vào nền, Android áp dụng chế độ Doze; thiết bị ghi và thiết bị nhận lại thường gắn với hai bản API khác nhau. Thiết kế dùng WebSocket~\cite{rfc6455} khi ứng dụng ở tiền cảnh, FCM khi chạy nền và Redis Pub/Sub làm cầu nối giữa các bản API. Hình~\ref{fig:seq-thoi-gian-thuc} minh họa trưởng nhóm phân công trên thiết bị A, thiết bị B đang mở ứng dụng, còn thiết bị C của người được giao đang chạy nền.

\diagram[\textwidth]{c3-17-seq-thoi-gian-thuc}{Biểu đồ tuần tự đồng bộ thời gian thực qua WebSocket, Redis Pub/Sub và FCM}{fig:seq-thoi-gian-thuc}

Khi mở socket, bản API xác thực JWT, nạp các không gian mà người dùng là thành viên rồi đăng ký các kênh \texttt{space:\{sid\}} tương ứng cùng kênh riêng \texttt{user:\{uid\}} (bước 1--3). Giao dịch phân công không tạo thông báo trực tiếp mà chèn tác vụ \texttt{AssignmentDigest} của module shopping (bước 6), duy nhất theo cặp (người được giao, danh sách) với lịch chạy lùi 5 giây, nên các phân công liên tiếp cho cùng người trong cùng danh sách chỉ sinh một tác vụ. Khi chạy, tác vụ này đếm $k$ mặt hàng đang giao trong bảng của chính shopping rồi chèn \texttt{PushNotification} mang sự kiện \texttt{item.assigned} gồm $k$ và mốc đợt giao trong cùng giao dịch (bước 14); nhờ vậy module notification không phải đọc bảng của shopping. \texttt{NotifyWorker} tôn trọng thiết lập \texttt{assignment\_notify\_on} của người nhận và dùng \texttt{dedupe\_key} chứa mốc đợt, nên mỗi đợt chỉ sinh một thông báo “Bạn được giao mua $k$ mặt hàng” dù tác vụ bị chạy lại (bước 15--16).

WebSocket chỉ mang \emph{gợi ý thay đổi}, dữ liệu luôn được kéo qua \texttt{/changes}. Nhờ vậy thời gian thực và bắt kịp sau ngoại tuyến dùng chung một đường dữ liệu, quyền được kiểm tra lại ở mỗi lần kéo, và việc Redis Pub/Sub có thể mất thông điệp trở nên vô hại vì lần kéo sau sẽ bù theo \texttt{seq}. Máy chủ gửi \texttt{ping} mỗi 25 giây và đóng kết nối nếu không có \texttt{pong} sau 10 giây; ứng dụng nối lại với độ trễ tăng theo hàm mũ có nhiễu ngẫu nhiên (1–30 giây). Sau mỗi lần làm mới access token, ứng dụng gửi token mới trên socket, quá hạn thì máy chủ đóng với mã 4401; khi người dùng đăng xuất hoặc bị xóa khỏi nhóm, sự kiện trên kênh \texttt{user:\{userId\}} khiến bản API đóng socket hoặc hủy theo dõi không gian, giữ \br{BR23} cả trên kênh thời gian thực. Thiết kế không dựa vào thông báo nền im lặng vì iOS có thể trì hoãn hoặc bỏ qua chúng~\cite{apns}; chỉ sự kiện cần biết ngay (phân công, nhắc hạn dùng, hoạt động nhóm) được gửi thành thông báo hiển thị qua FCM, chuyển tiếp qua APNs với iOS. Bảng~\ref{tab:c3-ngan-sach-do-tre} phân bổ ngân sách 3 giây cho từng chặng; các con số là mục tiêu thiết kế, không phải số đo.

{\small
\begin{longtable}{|L{4.4cm}|L{2.9cm}|C{1.6cm}|L{\dimexpr\textwidth-8\tabcolsep-5\arrayrulewidth-4.4cm-2.9cm-1.6cm\relax}|}
\caption{Phân bổ ngân sách độ trễ đồng bộ thời gian thực (mục tiêu thiết kế)}\label{tab:c3-ngan-sach-do-tre}\\
\hline
\hd{Chặng} & \hd{Thành phần} & \hd{Mục tiêu} & \hd{Biện pháp giữ ngân sách}\\ \hline
\endfirsthead
\multicolumn{4}{l}{\footnotesize\itshape Bảng \thetable\ (tiếp theo)}\\ \hline
\hd{Chặng} & \hd{Thành phần} & \hd{Mục tiêu} & \hd{Biện pháp giữ ngân sách}\\ \hline
\endhead
Thiết bị A gửi yêu cầu ghi & Mạng 4G, Caddy & 400 ms & HTTP/2 giữ kết nối, không bắt tay TLS lại\\ \hline
Xử lý và commit giao dịch & API, PostgreSQL & 250 ms & Truy vấn theo chỉ mục, khóa ngắn\\ \hline
Phát tán tới bản API giữ socket & Redis Pub/Sub & 50 ms & Cùng mạng Docker nội bộ\\ \hline
Gửi gợi ý xuống thiết bị B & WebSocket & 300 ms & Socket mở sẵn, khung dữ liệu nhỏ\\ \hline
Gom gợi ý và kéo \texttt{/changes} & Thiết bị B, API & 800 ms & Gom 200 ms; truy vấn theo khóa \texttt{(space\_id, seq)}\\ \hline
Áp vào SQLite, vẽ lại giao diện & drift, Riverpod & 300 ms & Áp theo lô trong một giao dịch SQLite\\ \hline
\textbf{Tổng} & & \textbf{2,1 s} & Dự trữ 0,9 s cho dao động mạng di động\\ \hline
\end{longtable}
}

% ---------------------------------------------------------------------
\subsubsection{Nhắc hạn dùng bảo đảm không gửi trùng}\label{sssec:kt-nhac-han}

\uc{UC31} và \br{BR13} yêu cầu mỗi người nhận tối đa một thông báo tổng hợp mỗi ngày, gộp các lô sắp hoặc vừa hết hạn trong mọi không gian mà người đó là thành viên; \nfr{NFR12} yêu cầu lỗi tạm thời được thử lại ba lần và ít nhất 99\% thông báo được tạo trong 15 phút kể từ giờ chạy. Hình~\ref{fig:seq-nhac-nho} ở Chương~\ref{chap:dac-ta} là góc nhìn nghiệp vụ; Hình~\ref{fig:seq-nhac-han-dung-be} là góc nhìn bên trong backend, gồm ba tác vụ River nối tiếp nhau.

\diagram[\textwidth]{c3-18-seq-nhac-han-dung-be}{Biểu đồ tuần tự nhắc hạn dùng phía backend với ba lớp chống trùng}{fig:seq-nhac-han-dung-be}

\texttt{ReminderDispatcher} là tác vụ định kỳ do bản worker giữ vai trò leader chèn mỗi phút một lần. Nó đọc \texttt{reminder.runAt} từ bộ đệm \texttt{cfg:all} (sống 5 phút và bị xóa ngay qua kênh Pub/Sub \texttt{cfg} khi \uc{UC53} đổi cấu hình), rồi chỉ làm việc khi giờ Việt Nam đã đến hoặc vượt \texttt{runAt} và hôm nay chưa chạy (bước 2--5). Điều kiện “đã đến hoặc vượt” thay cho “đúng 08:00” giúp tự chạy bù khi worker ngừng đúng giờ hẹn (bước 1a của \uc{UC31}), còn chu kỳ một phút giới hạn độ trễ khởi động ở mức một phút, để dành gần trọn ngân sách 15 phút cho việc tạo thông báo. Dấu “hôm nay đã chạy” là một tác vụ duy nhất theo ngày được chèn bằng \texttt{InsertTx} ở đầu giao dịch; cùng giao dịch đó, một truy vấn tìm người dùng có lô \texttt{active} với $-1 \le \texttt{expires\_on} - \texttt{\$today} \le 7$ ở bất kỳ không gian nào họ là thành viên, rồi chèn hàng loạt tác vụ \texttt{UserReminder} duy nhất theo cặp (người dùng, ngày) (bước 6). Nếu tiến trình chết giữa chừng, giao dịch rollback làm mất cả dấu ngày, và lượt chạy phút sau làm lại từ đầu.

Mỗi tác vụ \texttt{UserReminder} đọc thiết lập của người dùng và tính các lô theo ngưỡng $N$ riêng (\texttt{reminder\_days}, hoặc \texttt{reminder.defaultDays} khi để trống) trên mọi không gian của người đó, với $-1 \le d \le N$ để bao gồm cả lô vừa quá hạn hôm qua theo \br{BR13} (bước 7). Người tắt nhắc hạn dùng (bước 4a) hoặc không còn lô nào trong ngưỡng (bước 5a) được kết thúc tác vụ mà không tạo thông báo (bước 8). Ngược lại, thông báo nêu tên tối đa \texttt{reminder.maxLotsPerNotice} lô gần hạn nhất (tham số cấu hình của \uc{UC53}, mặc định 5, miền 1--10) để nội dung vừa màn hình khóa, phần còn lại ghi “và $k$ lô khác”; câu \texttt{INSERT} dùng \texttt{dedupe\_key = userId:ngày} với \texttt{ON CONFLICT DO NOTHING RETURNING} (bước 9), và chỉ khi thực sự chèn dòng mới thì \texttt{PushNotification} mới được chèn vào hàng đợi \texttt{notify} trong cùng giao dịch (bước 10). Chống trùng vì vậy gồm ba lớp: dấu ngày của bộ điều phối, tác vụ \texttt{UserReminder} duy nhất theo người và ngày, và \texttt{dedupe\_key} duy nhất trong \texttt{notifications}; hai lớp đầu tiết kiệm công việc, lớp cuối bảo đảm tính đúng kể cả khi quản trị viên chạy lại tác vụ bằng tay. Bước 2 của \uc{UC31} không cần ghi gì vì trạng thái hạn dùng được suy diễn lúc truy vấn theo ngưỡng của người xem, và thông báo chỉ chứa lô thuộc không gian người nhận là thành viên nên không vi phạm \br{BR23}.

\texttt{NotifyWorker} gửi tới từng thiết bị chưa có bản giao \texttt{sent}, với khóa gộp là mã thông báo (bước 12--13). Gửi thành công ghi \texttt{sent}; FCM báo mã không hợp lệ (\texttt{UNREGISTERED}) thì thiết bị bị xóa khỏi tài khoản (bước 7a); lỗi tạm thời được ghi \texttt{failed} và tác vụ trả lỗi ở cuối vòng lặp để River thử lại tối đa ba lần với độ trễ tăng theo hàm mũ, mỗi lần chỉ gửi tiếp các thiết bị chưa \texttt{sent} (bước 7b); thông báo vẫn nằm trong hộp thư trong ứng dụng dù đẩy thất bại. Tiêu chí của \nfr{NFR12} được tách thành hai chỉ số mức dịch vụ đo trực tiếp từ dữ liệu: chỉ số \emph{tạo} là khoảng cách từ \texttt{runAt} tới \texttt{notifications.created\_at}, với mục tiêu 99\% trong 15 phút; chỉ số \emph{gửi} là khoảng cách từ \texttt{created\_at} tới \texttt{deliveries.sent\_at}, theo dõi riêng vì phụ thuộc FCM nằm ngoài quyền kiểm soát của hệ thống. Với giả định thiết kế 5.000 tài khoản, mỗi ngày có tối đa 5.000 thông báo nhắc, một con số nhỏ so với năng lực chèn hàng loạt của PostgreSQL và hạn mức của FCM.

% ---------------------------------------------------------------------
\subsubsection{Gợi ý món và sinh danh sách mua sắm}\label{sssec:kt-goi-y}

\uc{UC40} phải xếp hạng 500 công thức theo 200 lô trong 2 giây (\nfr{NFR09}). Pha lọc dùng ba truy vấn SQL: tồn kho còn hạn gộp theo thực phẩm và nhóm đại lượng; công thức ứng viên (mẫu đã công bố và công thức của không gian, lọc theo bữa \texttt{mealSlot}) có ít nhất một nguyên liệu trong tồn kho; nguyên liệu của ứng viên cùng danh sách yêu thích. Khi người dùng mở từ màn hình thực phẩm sắp hết hạn (\uc{UC30}), tham số \texttt{expiringOnly} chỉ giữ ứng viên có dùng lô sắp hết hạn. Pha chấm điểm theo \br{BR19} chạy trong bộ nhớ như Mã nguồn~\ref{code:c3-goi-y}, với trọng số, $\theta$ và $K$ đọc từ cấu hình của \uc{UC53} (\nfr{NFR20}).

\codecaption{Mã giả Go của bước chấm điểm gợi ý món}{code:c3-goi-y}
\begin{Verbatim}
// Suggest chấm điểm theo BR19; trọng số, θ, K đọc từ system_configs.
func Suggest(inv map[FoodKey]Stock, cands []Recipe, fav map[ID]bool,
    n, servings int, cfg Config) (top, near []Scored) {
    best, alt := NewMinHeap(cfg.MaxResults), NewMinHeap(3)
    for _, c := range cands {
        k := float64(servings) / float64(c.Servings)
        ok, req, e := 0, 0, 0.0
        for _, ing := range c.Ingredients {
            st, has := inv[ing.Key()] // cùng thực phẩm, cùng đại lượng
            if ing.Required {
                req++
                if has && st.BaseQty >= ing.BaseQty*k { ok++ }
            }
            if has && st.DaysLeft >= 0 && st.DaysLeft <= n {
                e += float64(n+1-st.DaysLeft) / float64(n+1)
            }
        }
        r := float64(ok) / float64(max(req, 1))
        s := cfg.WR*r + cfg.WE*min(1, e) + cfg.WF*b2f(fav[c.ID])
        alt.Push(Scored{c, r, s}, r) // dự phòng MSG31: R cao nhất
        if r >= cfg.Theta { best.Push(Scored{c, r, s}, s) }
    }
    return best.SortedDesc(), alt.SortedDesc()
}
\end{Verbatim}

\texttt{DaysLeft} là số ngày còn lại của lô gần hết hạn nhất như trong công thức~\eqref{eq:diem-han-dung}; nguyên liệu khác nhóm đại lượng với tồn kho được coi là không đáp ứng (\br{BR16}). Độ phức tạp là $O(C \cdot I + C \log K)$ với $C$ công thức, mỗi công thức $I$ nguyên liệu, tức vài nghìn phép tra bảng băm; thời gian bị chi phối bởi ba truy vấn có chỉ mục. Đống thứ hai giữ ba món có $R$ cao nhất để trả kèm \msg{MSG31}. Kết quả được lưu đệm 10 phút trong Redis hash \texttt{suggest:\{sid\}}, trường theo người dùng, bữa và số khẩu phần vì $N$, $F$ thuộc từng người. Bộ đệm bị xóa bởi móc sau commit của mọi giao dịch phát sinh \texttt{lot.changed} hoặc \texttt{recipe.changed}, bởi sự kiện \texttt{config.changed} khi khóa \texttt{suggest.*} đổi, và khi người dùng đổi ngưỡng $N$ qua A13 (xóa cho mọi không gian của người đó); cấu hình chính nó nằm trong \texttt{cfg:all}, sống 5 phút và bị xóa ngay qua kênh \texttt{cfg}, nên trọng số mới có hiệu lực ở lần gợi ý kế tiếp.

Sinh danh sách (\uc{UC22}, \br{BR20}) là bài toán ghép nhu cầu với nguồn cung có khoảng hiệu lực: nhu cầu là cặp (ngày nấu, thực phẩm) đã quy về đơn vị cơ sở; lô còn hạn dùng được đến \texttt{expires\_on}, mặt hàng chờ mua dùng được từ \texttt{planned\_date}. Thuật toán duyệt nhu cầu theo ngày nấu tăng dần, lấy từ nguồn hợp lệ có hạn kết thúc sớm nhất trước (lô theo FEFO, hàng chờ mua sau cùng), phần còn lại là lượng thiếu. Chiến lược tham lam “kết thúc sớm nhất trước” này cho tổng lượng thiếu nhỏ nhất, không tính một lượng hai lần và tái tạo đúng kết quả 0,3 kg trong ví dụ ở Chương~\ref{chap:dac-ta}; với hàng đợi ưu tiên theo từng thực phẩm, độ phức tạp là $O((D + L + P)\log(L + P))$ cho $D$ nhu cầu, $L$ lô và $P$ mặt hàng chờ mua.

% ---------------------------------------------------------------------
\subsubsection{Tìm kiếm tiếng Việt với Elasticsearch}\label{sssec:kt-tim-kiem}

Người dùng thường gõ không dấu (“thit ba chi”), dùng từ địa phương (heo và lợn), gõ sai và chờ gợi ý trong 500 ms (\nfr{NFR09}), điều truy vấn \texttt{LIKE} không đáp ứng được. Elasticsearch~\cite{elasticsearch} vì vậy phục vụ \uc{UC42}, \uc{UC43} cho hai loại dữ liệu: danh mục thực phẩm chung cùng thực phẩm riêng của các không gian, và công thức mẫu cùng công thức của các không gian. Lô trong tủ và mặt hàng chỉ vài trăm bản ghi mỗi không gian nên được tìm trực tiếp bằng PostgreSQL qua cổng \texttt{StockQuery.SearchLots} của pantry. Mã nguồn~\ref{code:c3-es-mapping} là thiết lập của chỉ mục \texttt{foods\_v1}.

\codecaption{Bộ phân tích và ánh xạ của chỉ mục thực phẩm}{code:c3-es-mapping}
\begin{Verbatim}
PUT /foods_v1
{ "settings": { "analysis": {
    "filter": {
      "vi_syn": { "type": "synonym_graph", "updateable": true,
                  "synonyms": ["heo, lợn", "ngô, bắp", "dứa, thơm"] },
      "pfx": { "type": "edge_ngram", "min_gram": 2, "max_gram": 15 } },
    "analyzer": {
      "vi_base": { "tokenizer": "standard",
        "filter": ["lowercase", "asciifolding"] },
      "vi_fold": { "tokenizer": "standard",
        "filter": ["lowercase", "asciifolding", "vi_syn"] },
      "vi_pfx": { "tokenizer": "standard",
        "filter": ["lowercase", "asciifolding", "pfx"] } } } },
  "mappings": { "properties": {
    "name": { "type": "text", "analyzer": "vi_base",
      "search_analyzer": "vi_fold", "fields": {
        "prefix": { "type": "text", "analyzer": "vi_pfx",
                    "search_analyzer": "vi_base" },
        "exact": { "type": "text", "analyzer": "standard" } } },
    "aliases": { "type": "text", "analyzer": "vi_base",
                 "search_analyzer": "vi_fold" },
    "scope": { "type": "keyword" }, "status": { "type": "keyword" },
    "owner_space_id": { "type": "keyword" },
    "category_id": { "type": "keyword" } } } }
\end{Verbatim}

\texttt{asciifolding} bỏ dấu và chuyển “đ” thành “d”, nên “thịt” và “thit” cho cùng từ tố. \texttt{vi\_fold} thêm \texttt{synonym\_graph} và chỉ dùng lúc truy vấn, đúng phạm vi thiết kế của bộ lọc đồng nghĩa dạng đồ thị; nhờ cờ \texttt{updateable}, danh sách đồng nghĩa được nạp lại mà không đánh lại chỉ mục. Trường \texttt{name.prefix} với \texttt{edge\_ngram} 2–15 phục vụ gợi ý khi gõ; vì bỏ dấu làm “cá” và “cà” trùng nhau, trường \texttt{name.exact} giữ dấu để cộng điểm khi người dùng gõ có dấu. Vì một chỉ mục chứa dữ liệu của nhiều không gian, mọi truy vấn bắt buộc kèm bộ lọc phạm vi (\br{BR23}): với thực phẩm là \texttt{scope = global} $\lor$ \texttt{owner\_space\_id = sid}, với công thức là (\texttt{kind = template} $\land$ \texttt{status = published}) $\lor$ \texttt{space\_id = sid}. Bộ lọc được thêm trong một hàm dựng truy vấn duy nhất của adapter search, không để từng handler tự ghép, và có kiểm thử riêng. Hình~\ref{fig:seq-tim-kiem-chi-muc} thể hiện hai luồng của module search: đánh chỉ mục và truy vấn.

\diagram[\textwidth]{c3-37-seq-tim-kiem-chi-muc}{Biểu đồ tuần tự đánh chỉ mục và truy vấn tìm kiếm có dự phòng}{fig:seq-tim-kiem-chi-muc}

Ở luồng đánh chỉ mục, mọi API ghi thực phẩm hoặc công thức (D04, P02 khi tạo thực phẩm riêng, R02, R04...) chèn tác vụ \texttt{IndexFood} hoặc \texttt{IndexRecipe} vào hàng đợi \texttt{index} trong cùng giao dịch (bước 2), nên chỉ mục không bao giờ bỏ sót một lần ghi đã commit. \texttt{IndexWorker} không dùng dữ liệu trong payload mà đọc bản ghi hiện tại từ PostgreSQL rồi ghi bằng \texttt{\_bulk} với \texttt{version\_type = external} theo cột \texttt{version} (bước 3--4), nên tác vụ đến muộn không ghi đè tài liệu mới hơn. Bản ghi đã xóa, thực phẩm đã gộp vào bản chung, dữ liệu của nhóm đã giải tán hoặc của tài khoản đã xóa đều dẫn tới xóa tài liệu, giữ đúng yêu cầu xóa dữ liệu của \nfr{NFR18}; đối soát hằng tuần so cặp \texttt{id} và \texttt{version} giữa hai kho để nạp lại phần lệch (bước 5). Đổi ánh xạ bằng cách nạp \texttt{foods\_v2} rồi chuyển alias \texttt{foods} trong một lệnh nguyên tử. Ở luồng truy vấn, gợi ý tên Q02 bị giới hạn 20 yêu cầu mỗi giây cho mỗi người dùng bằng khóa \texttt{rl:suggest:\{uid\}}, vượt thì nhận 429 mã 4290 (\msg{MSG36}); truy vấn \texttt{multi\_match} trên \texttt{name} và \texttt{name.prefix} dùng \texttt{fuzziness: AUTO}, lọc \texttt{status = active} cùng bộ lọc phạm vi và đặt thời hạn 300 ms (bước 9). Khi Elasticsearch lỗi hoặc quá hạn, truy vấn chuyển sang so khớp trigram của PostgreSQL trên hàm bao \texttt{f\_unaccent} khai báo \texttt{IMMUTABLE}, vì hàm \texttt{unaccent} gốc chỉ là \texttt{STABLE} nên không dùng được trong chỉ mục (bước 11); chế độ suy giảm này mất đồng nghĩa nhưng vẫn trả kết quả với cùng bộ lọc phạm vi. Q01 còn tìm lô trong tủ qua \texttt{StockQuery.SearchLots} (bước 12) và trả kết quả theo nhóm tủ lạnh, công thức, danh mục; không có kết quả thì ứng dụng hiển thị \msg{MSG18}.

% ---------------------------------------------------------------------
\subsubsection{Tra cứu mã vạch có bộ đệm và ngắt mạch}\label{sssec:kt-ma-vach}

\uc{UC26} phụ thuộc Open Food Facts~\cite{openfoodfacts}, một dịch vụ cộng đồng có độ trễ khó đoán và giới hạn tần suất, trong khi use case yêu cầu phản hồi trong 5 giây. Hình~\ref{fig:seq-ma-vach} thể hiện chuỗi tra cứu có bộ đệm và bộ ngắt mạch (circuit breaker).

\diagram[0.85\textwidth]{c3-19-seq-ma-vach}{Biểu đồ tuần tự tra cứu mã vạch qua bộ đệm, danh mục và Open Food Facts}{fig:seq-ma-vach}

Ứng dụng kiểm tra chữ số kiểm tra và tra thực phẩm riêng của không gian trong bản sao cục bộ trước, nên \texttt{BarcodeService.Lookup} của API P08 chỉ tra danh mục với điều kiện \texttt{scope = 'global'} (bước 4), và khóa \texttt{barcode:\{code\}} dùng chung được cho mọi người mà không lộ dữ liệu nhóm (\br{BR23}). Ràng buộc duy nhất của mã vạch vì thế cũng tách theo phạm vi: duy nhất toàn cục với bản chung, duy nhất trong từng không gian với thực phẩm riêng. Bộ ngắt mạch được hỏi trạng thái trước khi rẽ nhánh (bước 5--6): thấy trong danh mục thì kết quả được lưu đệm 30 ngày và trả ngay; không thấy mà mạch đang mở thì API trả ngay 504 mã 5004 (\msg{MSG23}) thay vì chiếm giữ goroutine; chỉ khi mạch đóng hoặc nửa mở mới gọi dịch vụ ngoài với thời hạn 3 giây, giữ cả chuỗi trong giới hạn 5 giây. Kết quả “không có” được lưu 1 ngày (bộ đệm phủ định) và trả 404 mã 4001, còn lỗi quá thời gian không được lưu vì chỉ là tạm thời. Bộ ngắt mạch~\cite{nygard} mở sau 5 lỗi liên tiếp; sau 30 giây ở trạng thái mở, nó cho đúng một lời gọi thăm dò ở trạng thái nửa mở. Dữ liệu từ dịch vụ ngoài được coi là không tin cậy~\cite{owaspapi}: chỉ lấy trường cần thiết, cắt tên còn 100 ký tự và không tải ảnh theo URL trả về. Khi người dùng lưu lô ở \uc{UC25}, máy chủ tạo thực phẩm riêng của không gian ở trạng thái \texttt{pending\_review}, nguồn \texttt{open\_food\_facts}, kèm mã vạch trong cùng giao dịch. Khi quản trị viên duyệt (\uc{UC49}), thực phẩm này được nâng thành bản chung, hoặc được gộp vào bản chung nếu danh mục đã có thực phẩm trùng mã vạch; cả hai trường hợp đều phát sự kiện xóa khóa bộ đệm tương ứng.

% ---------------------------------------------------------------------
%  Hiện thực hóa use case (cuối mục Thiết kế chi tiết). Ma trận đầy đủ và
%  biểu đồ tuần tự từng use case nằm ở Phụ lục (backmatter/pl/e-hien-thuc-uc.tex).
% ---------------------------------------------------------------------
\subsection{Hiện thực hóa các use case trọng tâm}\label{ssec:hien-thuc-uc}

Các mục trước mô tả thiết kế theo chiều ngang: module, tầng và cơ chế dùng chung. Mục này đi theo chiều dọc, lần theo từng use case của Chương~\ref{chap:dac-ta} xuyên qua các tầng để kiểm chứng rằng mỗi bước, mỗi nhánh và mỗi quy tắc đều có phần tử thiết kế đảm nhận. Kỹ thuật sử dụng là hiện thực hóa use case (use case realization)~\cite{uml251}: use case được ánh xạ thành một cộng tác giữa các thành phần cụ thể và mô tả bằng biểu đồ tuần tự. Khác với biểu đồ hoạt động của Chương~\ref{chap:dac-ta} vốn coi hệ thống là hộp đen, các biểu đồ ở đây nhìn vào bên trong máy chủ, từ API qua dịch vụ ứng dụng và tầng miền tới PostgreSQL, Redis, River và dịch vụ bên ngoài.

Mọi use case theo cùng một khuôn. Use case được giao cho một dịch vụ ứng dụng (application service) của module sở hữu dữ liệu chính, gọi qua một API trong hợp đồng OpenAPI; riêng \uc{UC14} do tác vụ định kỳ của River khởi phát. Handler chỉ phân tích yêu cầu, xác thực và ánh xạ lỗi miền thành mã phản hồi. Dữ liệu vào được kiểm tra hợp lệ trước mọi truy vấn; mọi phép so sánh ngày nghiệp vụ dùng tham số \texttt{\$today} là ngày theo giờ Việt Nam do \texttt{platform/clock} cấp, không dùng \texttt{current\_date} của phiên CSDL vốn chạy theo UTC. Mỗi lệnh ghi chạy trong đúng một giao dịch do \texttt{TxManager} mở, với \texttt{statement\_timeout} và \texttt{lock\_timeout} 3 giây; khóa chết (40P01) hoặc quá thời gian chờ khóa được thử lại tối đa ba lần, sau đó trả 503 mã 5003. Với API có \texttt{Idempotency-Key}, câu lệnh đầu tiên của giao dịch là chèn khóa lũy đẳng ở trạng thái \texttt{in\_progress}; khóa đã hoàn tất thì phản hồi cũ được phát lại, khóa đang xử lý trả 409 mã 4002, cùng khóa khác nội dung trả 400 mã 2002, còn giao dịch lỗi bị hủy cùng khóa nên chỉ kết quả thành công được lưu (Hình~\ref{fig:seq-dong-bo-ngoai-tuyen}). Các khóa tiếp theo được lấy đúng thứ tự toàn cục ở mục~\ref{sssec:kt-fefo}: khóa tư vấn (advisory lock) theo không gian cho thao tác đổi thành viên, các dòng \texttt{space\_members} theo \texttt{user\_id}, rồi danh sách, món, mặt hàng, lô, và sau cùng bộ đếm \texttt{spaces.change\_seq}. Module khác tham gia qua cổng và dùng chung \texttt{pgx.Tx}. Mọi truy vấn đều lọc theo \texttt{space\_id}: người không phải thành viên của không gian trong đường dẫn, hoặc truy cập tài nguyên thuộc không gian khác, nhận 404 mã 4001 để không lộ sự tồn tại của tài nguyên (chống dò BOLA, \br{BR23}); thành viên thiếu vai trò cần thiết nhận 403 mã 3007 kèm \msg{MSG13}. Quy tắc nghiệp vụ được bảo vệ hai lớp: tầng miền kiểm tra trước để trả đúng mã quy tắc và thông điệp, còn ràng buộc CSDL (chỉ mục duy nhất từng phần, \texttt{CHECK}, khóa dòng) là hàng rào cuối khi có truy cập đồng thời. Tác dụng phụ ra ngoài giao dịch như đẩy thông báo, gửi thư, đánh chỉ mục hay xuất tệp được phát thành sự kiện miền, tức tác vụ River chèn trong cùng giao dịch (Hình~\ref{fig:seq-outbox-mua-hang}). Thực thể đồng bộ ghi \texttt{sync.space\_changes} ngay trước \texttt{COMMIT} và phát gợi ý thay đổi \texttt{space:\{sid\}} qua Redis sau commit (Hình~\ref{fig:seq-thoi-gian-thuc}); phép đọc trải nhiều bảng dùng giao dịch \texttt{REPEATABLE READ READ ONLY} để có ảnh chụp nhất quán. Xung đột phiên bản trả 409 mã 4002 kèm \msg{MSG35}, vượt giới hạn tần suất trả 429 mã 4290 kèm \msg{MSG36}. Khi một quy tắc buộc module thượng nguồn tác động lên dữ liệu của module hạ nguồn, module thượng nguồn định nghĩa giao diện móc (hook) và module hạ nguồn hiện thực nó lúc khởi động: \texttt{identity/api.ActivationHook} do household hiện thực để tạo không gian cá nhân, \texttt{household/api.MemberLeaveHook} do shopping hiện thực để trả mặt hàng đang giao về Cần mua (\br{BR07}), \texttt{recipe/api.RecipeDeleteHook} do mealplan hiện thực; nhờ vậy đồ thị phụ thuộc không có chu trình (mục~\ref{sec:kien-truc}).

Để giữ mạch trình bày của chương, phần chi tiết của việc hiện thực hóa được đặt ở Phụ lục~\ref{pl:hien-thuc-uc}. Phụ lục này gồm ma trận truy vết đủ 25 use case tới module, API, bảng dữ liệu, sự kiện và sơ đồ (Bảng~\ref{tab:c3-uc-thiet-ke}), cùng biểu đồ tuần tự phía máy chủ của các luồng chưa được minh họa ở mục~\ref{ssec:ky-thuat}: đăng ký và xác thực OTP (Hình~\ref{fig:seq-dang-ky-otp}), tham gia nhóm (Hình~\ref{fig:seq-tham-gia-nhom}), xóa thành viên, chuyển quyền và rời nhóm (Hình~\ref{fig:seq-roi-nhom-chuyen-quyen}), phân công người mua (Hình~\ref{fig:seq-phan-cong}), sinh danh sách từ kế hoạch (Hình~\ref{fig:seq-sinh-ds-ke-hoach}), thêm thực phẩm (Hình~\ref{fig:seq-them-thuc-pham}), gợi ý món (Hình~\ref{fig:seq-goi-y-mon}), lên lịch bữa ăn (Hình~\ref{fig:seq-lap-lich}), báo cáo và xuất tệp (Hình~\ref{fig:seq-bao-cao-xuat}), quản trị cấu hình, khóa tài khoản (Hình~\ref{fig:seq-quan-tri-cau-hinh}) và duyệt thực phẩm (Hình~\ref{fig:seq-duyet-thuc-pham}). Cùng với các hình ở mục~\ref{ssec:ky-thuat} và mục~\ref{sec:bao-mat}, mọi use case trọng tâm đều có biểu đồ hiện thực riêng.

Ma trận ở Phụ lục~\ref{pl:hien-thuc-uc} cho thấy ba dạng hiện thực: thao tác một thực thể trong một module (phần lớn gói tài khoản, công thức, quản trị); giao dịch trải nhiều module qua cổng (\uc{UC07}, \uc{UC11}, \uc{UC10} và \uc{UC19}), lý do chính để chọn một CSDL dùng chung; và phần xử lý bất đồng bộ qua River (\uc{UC05}, \uc{UC06}, \uc{UC08}, \uc{UC14}, \uc{UC22}, \uc{UC23}). Mười hai use case đặc tả đầy đủ ở Chương~\ref{chap:dac-ta} đều có biểu đồ tuần tự. Bốn use case \uc{UC09}, \uc{UC11}, \uc{UC14} và \uc{UC19} có biểu đồ riêng ở mục~\ref{ssec:ky-thuat}, còn \uc{UC13} dùng chung cơ chế FEFO của Hình~\ref{fig:seq-xac-nhan-nau}. Bảy use case còn lại là \uc{UC06}, \uc{UC08}, \uc{UC10}, \uc{UC11}, \uc{UC17}, \uc{UC18}, \uc{UC21}, cùng \uc{UC22}, được trình bày dưới đây. Ngoài nhóm đó, mục này hiện thực thêm những use case có luồng bảo mật hoặc đồng thời đáng chú ý: đăng ký, xác thực OTP và khôi phục mật khẩu (\uc{UC01} và \uc{UC03}); xóa thành viên, chuyển quyền và rời nhóm (\uc{UC07}); quản trị tài khoản, cấu hình và duyệt thực phẩm riêng (\uc{UC24}, \uc{UC23}, \uc{UC25}). Tìm kiếm và lọc (\uc{UC20}) có biểu đồ ở mục~\ref{sssec:kt-tim-kiem}.

% ---------------------------------------------------------------------
%  3.5. Thiết kế cơ sở dữ liệu
% ---------------------------------------------------------------------
\section{Thiết kế cơ sở dữ liệu}\label{sec:csdl}

Mục này cụ thể hóa tầng dữ liệu của kiến trúc ở mục~\ref{sec:kien-truc}, suy ra từ mô hình khái niệm (Hình~\ref{fig:mo-hinh-khai-niem}) và từ điển dữ liệu (Phụ lục~\ref{pl:tu-dien}). Quy tắc nghiệp vụ nào diễn đạt được bằng ràng buộc khai báo thì được cưỡng chế ngay trong cơ sở dữ liệu (CSDL), quy tắc phụ thuộc ngày hoặc cấu hình được kiểm tra ở ứng dụng trong cùng giao dịch.

\subsection{Chiến lược lưu trữ đa cơ sở dữ liệu}\label{ssec:luu-tru-da-csdl}

Hệ thống lưu trữ đa dạng (polyglot persistence)~\cite{kleppmann}, mỗi kho chọn theo cách truy cập dữ liệu (Bảng~\ref{tab:c3-luu-tru}). Để tránh ghi kép (dual write), PostgreSQL là nguồn sự thật (source of truth) duy nhất; kho khác chỉ chứa dữ liệu dẫn xuất hoặc tạm, cập nhật sau commit qua tác vụ River ghi trong cùng giao dịch (mục~\ref{ssec:ky-thuat}). Bộ đếm bảo mật lâu dài (\br{BR02}, \br{BR03}) cũng nằm trong PostgreSQL để mất Redis không làm mất trạng thái khóa tài khoản.

{\footnotesize\setlength{\tabcolsep}{4pt}\renewcommand{\arraystretch}{1.1}
\begin{longtable}{|L{3.5cm}|L{1.8cm}|L{3.8cm}|L{1.8cm}|L{\dimexpr\textwidth-10\tabcolsep-6\arrayrulewidth-3.5cm-1.8cm-3.8cm-1.8cm\relax}|}
\caption{Phân bổ dữ liệu cho các kho lưu trữ}\label{tab:c3-luu-tru}\\
\hline
\hd{Loại dữ liệu} & \hd{Nơi lưu} & \hd{Lý do} & \hd{Vai trò} & \hd{Phục hồi khi mất}\\ \hline
\endfirsthead
\multicolumn{5}{l}{\footnotesize\itshape Bảng \thetable\ (tiếp theo)}\\ \hline
\hd{Loại dữ liệu} & \hd{Nơi lưu} & \hd{Lý do} & \hd{Vai trò} & \hd{Phục hồi khi mất}\\ \hline
\endhead
Dữ liệu nghiệp vụ, nhật ký thay đổi, khóa lũy đẳng, hàng đợi River & PostgreSQL 16 & Giao dịch ACID nhiều module (\uc{UC11}, \uc{UC19}), ràng buộc, khóa dòng & Nguồn sự thật & pgBackRest: đầy đủ hằng ngày, WAL liên tục (\nfr{NFR13})\\ \hline
Giới hạn tần suất, bộ đệm, Pub/Sub & Redis 7 & Độ trễ thấp, TTL, thao tác nguyên tử & Tạm thời & Không sao lưu; tự nạp lại\\ \hline
Chỉ mục thực phẩm, công thức (chung và riêng của không gian) & Elastic\-search 8.15 & Tìm không dấu, lỗi gõ, gợi ý $\le$ 500 ms (\nfr{NFR09}) & Dẫn xuất & Đánh chỉ mục lại; tạm dùng trigram\\ \hline
Ảnh, tệp báo cáo xuất & MinIO & Tệp nhị phân, URL ký sẵn & Ảnh gốc: nguồn sự thật & Bucket sao lưu; biến thể và tệp xuất sinh lại\\ \hline
Bản sao không gian, hàng đợi ngoại tuyến & SQLite (thiết bị) & Dùng khi mất mạng (\nfr{NFR11}) & Bản sao & Tải lại qua API\\ \hline
\end{longtable}
}

\subsection{Mô hình dữ liệu quan hệ}\label{ssec:mo-hinh-quan-he}

Lược đồ PostgreSQL~\cite{postgresql} theo các quy ước chung. Khóa chính là UUIDv7~\cite{rfc9562}: thiết bị sinh được khi ngoại tuyến (\nfr{NFR11}) và tiền tố thời gian làm các lần chèn gần như tăng dần, ít phân mảnh chỉ mục hơn UUIDv4. Thời điểm dùng \texttt{timestamptz} (UTC); ngày nghiệp vụ dùng \texttt{date} theo giờ Asia/Ho\_Chi\_Minh, mọi câu SQL so sánh ngày nhận ngày hiện tại qua tham số \texttt{\$today} do \texttt{platform/clock} cấp thay vì \texttt{current\_date} của phiên chạy UTC, vì từ 00:00 đến 07:00 giờ Việt Nam hai giá trị này lệch nhau một ngày và làm sai \br{BR12}, \br{BR15}. Số lượng dùng \texttt{numeric(12,2)} (\br{BR14}), tiền dùng \texttt{bigint} (đồng), tập giá trị cố định dùng ENUM. Thực thể đồng bộ có \texttt{version} cho khóa lạc quan (lệnh sửa kèm \texttt{If-Match} thêm điều kiện \texttt{version = \$v}; không dòng nào được cập nhật thì trả mã 4002) và \texttt{deleted\_at} cho xóa mềm, tạo bia mộ (tombstone) giữ 90 ngày. Mỗi module sở hữu một schema cùng tên; khóa ngoại liên schema chỉ trỏ tới \texttt{identity.users}, \texttt{household.spaces}, \texttt{catalog.*}, còn tham chiếu giữa module ngang hàng (mặt hàng -- lô, bữa ăn -- công thức, biến động kho -- bữa ăn) là tham chiếu logic do giao dịch của module sở hữu giữ đúng, tránh chu trình khóa ngoại và giữ khả năng tách module. Riêng schema \texttt{report} là mô hình đọc (read model): ngoài bảng \texttt{exports} do report sở hữu, nó gồm các view \texttt{report.v\_*} chỉ \texttt{SELECT} trên shopping, pantry, mealplan, catalog và một materialized view; đây là ngoại lệ có chủ đích của quy tắc “mỗi module chỉ đọc schema của mình”, chấp nhận được vì report không bao giờ ghi vào schema khác. Toàn lược đồ có 29 bảng (kể cả materialized view) trong 12 schema. Hình~\ref{fig:erd-tong-quan} là ERD tổng quan của chín schema chứa dữ liệu nghiệp vụ; ba schema hỗ trợ notification, media, admin chỉ tham chiếu \texttt{users} nên được vẽ riêng ở Hình~\ref{fig:erd-thong-bao-he-thong}.

\diagram[\textwidth]{c3-20-erd-tong-quan}{Sơ đồ thực thể -- liên kết tổng quan theo schema}{fig:erd-tong-quan}

Lược đồ có hai trục. Trục sở hữu đi từ \texttt{spaces}: điều kiện \texttt{space\_id} trên mọi truy vấn sinh bằng sqlc~\cite{sqlc} vừa lọc dữ liệu vừa là ranh giới phân quyền (\br{BR23}); \texttt{foods} cũng có khóa ngoại tùy chọn \texttt{owner\_space\_id} tới \texttt{spaces} cho thực phẩm riêng của không gian (\uc{UC11}), còn công thức mẫu có \texttt{space\_id} rỗng. Trục danh mục hội tụ về \texttt{foods}, được mặt hàng, lô và nguyên liệu cùng tham chiếu. Ba nét đứt là ba tham chiếu logic giữa module ngang hàng: mặt hàng -- lô sinh ra khi nhập tủ (\uc{UC11}), bữa ăn -- công thức (\uc{UC17}) và biến động \texttt{cook} -- bữa ăn đã nấu (\uc{UC19}); nét đứt từ \texttt{spaces} tới \texttt{daily\_space\_stats} chỉ nguồn của materialized view. Schema \texttt{river} do River tự quản lý nên không vẽ. Trong các ERD chi tiết, chấm tròn đặc là NOT NULL, «GEN» là cột sinh, nét đứt là tham chiếu logic không có khóa ngoại, phần dưới đường chấm của một bảng ghi ràng buộc đáng chú ý, \texttt{created\_at} và \texttt{updated\_at} được lược bớt; các cột người thực hiện như \texttt{created\_by} được đánh dấu «FK» nhưng không vẽ đường nối để sơ đồ dễ đọc. Hình~\ref{fig:erd-tai-khoan-nhom} trình bày schema identity và household.

\diagram[\textwidth]{c3-21-erd-tai-khoan-nhom}{ERD schema identity và household}{fig:erd-tai-khoan-nhom}

Ba ghi chú dưới household là ba quy tắc “tối đa một” được cưỡng chế bằng chỉ mục duy nhất từng phần: mỗi không gian một trưởng nhóm (\br{BR05}), mỗi người một nhóm gia đình và một không gian cá nhân (\br{BR04}), mỗi nhóm một mã mời còn hiệu lực (\br{BR06}). Để chỉ mục về \texttt{space\_kind} không bị sai lệch, \texttt{space\_members} sao chép loại không gian và khóa ngoại ghép (\texttt{space\_id}, \texttt{space\_kind}) trỏ tới ràng buộc duy nhất (\texttt{id}, \texttt{kind}) của \texttt{spaces}, nên giá trị sao chép không thể khác bản gốc. \texttt{spaces.change\_seq} là bộ đếm thay đổi riêng từng không gian cho đồng bộ (Hình~\ref{fig:seq-dong-bo-ngoai-tuyen}). Trong identity, email tối đa 100 ký tự theo Phụ lục~\ref{pl:tu-dien}; \texttt{last\_login\_at} và \texttt{lock\_reason} phục vụ màn hình quản lý tài khoản của \uc{UC24}, trong đó \texttt{locked} là khóa tạm do đăng nhập sai (\br{BR03}, kèm \texttt{locked\_until}) còn \texttt{disabled} là khóa của quản trị viên. Khi xác minh OTP đặt lại mật khẩu (\uc{UC03}), máy chủ ghi băm của mã đặt lại 256 bit vào \texttt{otp\_challenges.reset\_token\_hash} với hạn 10 phút; A08 tiêu thụ mã bằng một câu \texttt{UPDATE ... SET reset\_token\_hash = NULL ... RETURNING} có điều kiện còn hạn, nên mã chỉ dùng được một lần kể cả khi hai yêu cầu đến cùng lúc. Hình~\ref{fig:erd-mua-sam-tu-lanh} thể hiện schema shopping và pantry.

\diagram[\textwidth]{c3-22-erd-mua-sam-tu-lanh}{ERD schema shopping và pantry}{fig:erd-mua-sam-tu-lanh}

Liên kết “tạo ra” giữa mặt hàng và lô là quan hệ không -- hoặc -- một ở cả hai đầu, được ghi hai chiều (\texttt{lot\_id}, \texttt{source\_item\_id}) trong giao dịch \uc{UC11}; mặt hàng tên tự do không có \texttt{food\_id} nên không sinh lô. Lô không sao chép đơn giá mà lưu \texttt{total\_cost}, bằng lượng thực mua nhân đơn giá tại thời điểm nhập tủ; giá trị lãng phí của một lần loại bỏ được tính bằng $\texttt{total\_cost} \times \text{lượng bỏ} / \texttt{initial\_qty}$. Cách tính theo tỷ lệ này đúng cả khi S14 đổi đơn vị của lô so với mặt hàng (ví dụ mua theo bó nhưng cất theo gam), điều mà công thức “đơn giá $\times$ lượng bỏ” không bảo đảm; lô nhập tay không có giá thì \texttt{total\_cost} rỗng và chỉ được tính theo số lượng (\br{BR21}). Cột sinh \texttt{base\_remaining} cộng tồn kho theo đơn vị cơ sở (\br{BR16}). Bảng \texttt{stock\_movements} chỉ ghi thêm: khi P04 sửa lượng còn của lô (\uc{UC13}), lượng còn không được hợp nhất theo kiểu lệnh sau thắng mà phải khớp \texttt{If-Match} (lệch trả 409 mã 4002), và chênh lệch được ghi thành một biến động \texttt{adjust} có dấu; nhờ vậy tổng biến động của một lô luôn khớp \texttt{initial\_qty} trừ \texttt{remaining\_qty}. Hình~\ref{fig:erd-cong-thuc-ke-hoach} gom schema catalog, recipe và mealplan.

\diagram[\textwidth]{c3-23-erd-cong-thuc-ke-hoach}{ERD schema catalog, recipe và mealplan}{fig:erd-cong-thuc-ke-hoach}

\texttt{foods} chứa cả thực phẩm chung lẫn thực phẩm riêng của không gian (\uc{UC11}), nên tính duy nhất của tên trong cùng loại được xét theo từng phạm vi (dấu «UQ*»), vì hai gia đình có thể cùng tự tạo một thực phẩm trước khi quản trị viên chuẩn hóa. Hai vòng tự tham chiếu thể hiện hai thao tác của quản trị và người dùng. \texttt{foods.merged\_into} ghi kết quả duyệt thực phẩm riêng (\uc{UC23}): duyệt thì bản ghi thành \texttt{global}; gộp thì bản ghi chuyển \texttt{inactive} và trỏ tới bản chuẩn, sự kiện \texttt{food.merged} sinh các tác vụ River lũy đẳng đổi \texttt{food\_id} trong pantry, shopping và recipe. Thực phẩm riêng có tên trùng một bản chung trong cùng loại chỉ được gộp, vì duyệt sẽ vi phạm chỉ mục duy nhất của phạm vi chung. \texttt{recipes.origin\_recipe\_id} ghi nguồn khi R09 sao chép công thức mẫu thành bản của không gian (\uc{UC15}, \uc{UC23}). Công thức thuộc không gian (\texttt{kind = space}) chứ không thuộc cá nhân; Chương~\ref{chap:dac-ta} gọi đó là bộ sưu tập cá nhân vì xét từ không gian cá nhân. Công thức mẫu dùng đủ ba trạng thái \texttt{draft}, \texttt{published}, \texttt{hidden} của \uc{UC23}, còn công thức không gian luôn \texttt{published}; R01 và M05 chỉ đọc công thức \texttt{published}. Với \texttt{units}, cờ \texttt{is\_base} cùng chỉ mục duy nhất từng phần bảo đảm mỗi nhóm đại lượng có đúng một đơn vị cơ sở hệ số 1 (\br{BR16}). Đổi nhóm đại lượng của một đơn vị đã được tham chiếu sẽ làm sai mọi lượng đang lưu, nên bị trigger \texttt{BEFORE UPDATE OF dimension} của \texttt{catalog.units} từ chối (422, mã 4004); kiểm tra đặt trong CSDL vì phải nhìn các bảng tham chiếu ở schema khác, điều module catalog không được phép truy vấn. Hình~\ref{fig:erd-thong-bao-he-thong} gồm các schema hỗ trợ.

\diagram[\textwidth]{c3-24-erd-thong-bao-he-thong}{ERD các schema notification, sync, media, admin và report}{fig:erd-thong-bao-he-thong}

\texttt{notifications.dedupe\_key} là chốt chống trùng cuối cùng của thông báo: \texttt{userId:ngày} với nhắc hạn dùng (\br{BR13}, mỗi người tối đa một thông báo mỗi ngày gộp mọi không gian) và khóa kèm mốc đợt giao với thông báo phân công. \texttt{deliveries.device\_id} là tham chiếu logic vì thiết bị có mã FCM không còn hợp lệ bị xóa khỏi \texttt{devices} nhưng lịch sử gửi vẫn giữ. \texttt{idempotency\_keys.status} phân biệt khóa đang xử lý (\texttt{in\_progress}) với khóa đã có phản hồi (\texttt{completed}); dòng được chèn bằng câu đầu tiên của giao dịch ghi và chỉ ghi phản hồi ngay trước commit, nên lỗi làm giao dịch rollback thì khóa cũng biến mất và chỉ kết quả thành công được phát lại. \texttt{media.objects} chỉ còn theo dõi ảnh, với \texttt{purpose} (ảnh đại diện, không gian, công thức, thực phẩm) và \texttt{target\_id} để Y04 kiểm tra quyền trên đối tượng đích. Tệp báo cáo xuất có bảng riêng \texttt{report.exports} (\uc{UC22}) lưu loại báo cáo, khoảng ngày, định dạng PDF hoặc CSV, trạng thái \texttt{pending}, \texttt{ready}, \texttt{failed}, khóa đối tượng, lỗi và hạn 7 ngày; B05 lọc theo cả \texttt{space\_id} và \texttt{requested\_by} sau khi kiểm tra thành viên. \texttt{daily\_space\_stats} được vẽ nét đứt vì là materialized view, không có khóa ngoại. Hai cột người thực hiện của admin (\texttt{updated\_by}, \texttt{actor\_id}) trỏ tới \texttt{users}; \texttt{updated\_by} rỗng với giá trị mặc định do di trú nạp.

Đối chiếu với mô hình khái niệm, phần lớn khái niệm ánh xạ một -- một thành bảng, trừ năm điểm. (1) “Nhóm” thành \texttt{household.spaces} có cột \texttt{kind}: không gian cá nhân là nhóm một thành viên nên mọi dữ liệu nghiệp vụ chỉ có một kiểu chủ sở hữu \texttt{space\_id} (\br{BR04}). (2) \texttt{space\_members} sao chép \texttt{space\_kind} để cưỡng chế \br{BR04} bằng ràng buộc duy nhất từng phần. (3) “Nhật ký sử dụng” thành sổ cái \texttt{pantry.stock\_movements} có thêm \texttt{stock\_in}, \texttt{cook} (\uc{UC19}) và \texttt{adjust} (\uc{UC13}) để đối soát tồn kho. (4) Thuộc tính đa trị “hạn dùng gợi ý” thành bảng \texttt{food\_shelf\_lives}. (5) “Trạng thái” lô tách đôi: \texttt{lifecycle} được lưu, còn trạng thái hạn dùng được suy diễn theo ngưỡng $N$ của người xem (\br{BR12}).

\subsection{Đặc tả bảng và tham số cấu hình}\label{ssec:dac-ta-bang}

Đặc tả chi tiết từng cột, kiểu dữ liệu và ràng buộc của các bảng trọng tâm, cùng bảng tóm tắt các bảng còn lại, được đặt ở Phụ lục~\ref{pl:dac-ta-bang} để giữ mạch trình bày; riêng \texttt{pantry.food\_lots} được minh họa bằng khối DDL ở mục~\ref{ssec:chi-muc}. Mục này trình bày phần có ý nghĩa thiết kế trực tiếp là cách các tham số nghiệp vụ được lưu và đọc.

Bảng \texttt{admin.system\_configs} hiện thực các tham số vận hành của \uc{UC25}. Bảng~\ref{tab:c3-tham-so} ánh xạ từng tham số của Bảng~\ref{tab:tham-so-he-thong} sang khóa cấu hình, miền được kiểm tra khi lưu, thành phần đọc và thời điểm có hiệu lực. \texttt{PUT /admin/configs} ghi nhiều khóa trong một giao dịch: kiểm tra miền và tổng trọng số bằng 1 trước khi ghi (sai thì trả \msg{MSG27}, không áp dụng một phần), ghi một dòng \texttt{audit\_logs} chứa giá trị cũ và mới, rồi phát đúng một sự kiện \texttt{config.changed}. Mọi bản API và worker đọc cấu hình qua \texttt{cfg:all} có TTL 5 phút và xóa bản đệm ngay khi nhận thông điệp trên kênh \texttt{cfg}, nên TTL chỉ là lưới an toàn khi một thông điệp Pub/Sub bị mất. Thời điểm hiệu lực theo một nguyên tắc chung: giá trị được chụp vào dữ liệu lúc tạo bản ghi (hạn của OTP, của mã mời) không hồi tố, còn giá trị dùng để tính lúc truy vấn ($N$, $\theta$, trọng số) có hiệu lực từ yêu cầu kế tiếp.

{\footnotesize\setlength{\tabcolsep}{4pt}\renewcommand{\arraystretch}{1.1}
\begin{longtable}{|L{2.2cm}|L{3.4cm}|L{1.95cm}|L{3.9cm}|L{\dimexpr\textwidth-10\tabcolsep-6\arrayrulewidth-2.2cm-3.4cm-1.95cm-3.9cm\relax}|}
\caption{Tham số cấu hình: khóa, miền, thành phần đọc và thời điểm hiệu lực}\label{tab:c3-tham-so}\\
\hline
\hd{Tham số} & \hd{Khóa} & \hd{Miền (mặc định)} & \hd{Thành phần đọc} & \hd{Thời điểm hiệu lực}\\ \hline
\endfirsthead
\multicolumn{5}{l}{\footnotesize\itshape Bảng \thetable\ (tiếp theo)}\\ \hline
\hd{Tham số} & \hd{Khóa} & \hd{Miền (mặc định)} & \hd{Thành phần đọc} & \hd{Thời điểm hiệu lực}\\ \hline
\endhead
Ngưỡng nhắc mặc định $N$ & \path{reminder.defaultDays} & 1--7 ngày (3) & Biểu thức \texttt{expiry\_state} khi \texttt{reminder\_days} rỗng (P01, P07, Q01); \texttt{UserReminder}; A12 & Yêu cầu và lượt nhắc kế tiếp\\ \hline
Khung giờ chạy nhắc nhở & \path{reminder.runAt} & 05:00--21:00 (08:00) & \texttt{ReminderDispatcher} kiểm tra mỗi phút; A12 & Lần kiểm tra kế tiếp; ngày đã chạy thì từ hôm sau\\ \hline
Số thành viên tối đa & \path{group.maxMembers} & 2--30 (10) & G05, G07 (đếm trong khóa tư vấn của không gian) & Lần mời, tham gia kế tiếp; không loại thành viên hiện có\\ \hline
Thời hạn mã mời & \path{invite.ttlHours} & 1--168 giờ (48) & G05 ghi \texttt{expires\_at} & Mã tạo sau thay đổi\\ \hline
Thời hạn OTP & \path{otp.ttlMinutes} & 2--15 phút (5) & A01, A03, A07 ghi \texttt{expires\_at} & Mã sinh sau thay đổi\\ \hline
Ngưỡng đáp ứng $\theta$ & \path{suggest.theta} & 0,3--1,0 (0,6) & M05 (\path{SuggestionEngine.Rank}) & Ngay; \texttt{config.changed} xóa \texttt{suggest:*}\\ \hline
Trọng số $R$, $E$, $F$ & \path{suggest.weights} & Mỗi số 0--1, tổng 1 (0,6; 0,3; 0,1) & M05 & Như trên\\ \hline
Số món gợi ý tối đa & \path{suggest.maxResults} & 3--30 (10) & M05 & Như trên\\ \hline
Khoảng lập kế hoạch tối đa & \path{mealplan.horizonDays} & 7--60 ngày (28) & M02, M03, S15 so với \texttt{\$today} (\br{BR18}, \msg{MSG28}) & Yêu cầu kế tiếp; món đã lập không bị xóa\\ \hline
Số lô tối đa mỗi thông báo tổng hợp & \path{reminder.maxLotsPerNotice} & 1--10 lô (5) & \texttt{UserReminder} & Lượt nhắc kế tiếp\\ \hline
\end{longtable}
}

Mười khóa trùng mười dòng của Bảng~\ref{tab:tham-so-he-thong}. Cột \texttt{value\_type}, \texttt{min\_value}, \texttt{max\_value} của \texttt{system\_configs} lưu luôn miền hợp lệ, nên trình kiểm tra của \texttt{PUT /admin/configs} là mã chung cho mọi khóa, chỉ riêng ràng buộc tổng trọng số bằng 1 cần luật riêng. Khi nội dung thông báo nhắc vượt \texttt{reminder.maxLotsPerNotice} lô, phần còn lại được gộp thành “và $k$ lô khác”; miền tối đa 10 lô cũng giữ tải tin dưới giới hạn 4 KB của FCM.

\subsection{Chỉ mục, ràng buộc toàn vẹn và tối ưu truy vấn}\label{ssec:chi-muc}

Chỉ mục được thiết kế từ truy vấn thực tế: B-tree phức hợp theo thứ tự lọc rồi sắp xếp, chỉ mục từng phần (partial index) và ràng buộc duy nhất từng phần cho quy tắc “tối đa một” (Bảng~\ref{tab:c3-chi-muc}, không liệt kê B-tree thông thường). Chỉ mục trigram dùng hàm bọc \texttt{f\_unaccent} khai báo \texttt{IMMUTABLE} vì \texttt{unaccent} gốc không dùng được. Mã nguồn~\ref{code:c3-ddl-food-lots} là định nghĩa bảng \texttt{pantry.food\_lots}, nơi tập trung nhiều quy tắc nhất.

\codecaption{Định nghĩa bảng \texttt{pantry.food\_lots} và các chỉ mục từng phần}{code:c3-ddl-food-lots}
\nopagebreak[4]
\begin{Verbatim}
CREATE TABLE pantry.food_lots (
  id               uuid PRIMARY KEY,                       -- UUIDv7
  space_id         uuid NOT NULL REFERENCES household.spaces (id),
  food_id          uuid NOT NULL REFERENCES catalog.foods (id),
  initial_qty      numeric(12,2) NOT NULL CHECK (initial_qty > 0),
  remaining_qty    numeric(12,2) NOT NULL,
  unit_id          uuid NOT NULL REFERENCES catalog.units (id),
  unit_factor      numeric(12,4) NOT NULL CHECK (unit_factor > 0),
  base_remaining   numeric GENERATED ALWAYS AS
                     (remaining_qty * unit_factor) STORED,
  storage_location catalog.storage_location NOT NULL,
  purchased_on     date NOT NULL,
  expires_on       date NOT NULL,
  lifecycle        pantry.lot_lifecycle NOT NULL DEFAULT 'active',
  source_item_id   uuid,              -- tham chiếu logic sang shopping
  total_cost       bigint CHECK (total_cost >= 0),      -- BR21, S14
  note             varchar(200),
  created_by       uuid NOT NULL REFERENCES identity.users (id),
  version          int NOT NULL DEFAULT 1,
  created_at       timestamptz NOT NULL DEFAULT now(),
  updated_at       timestamptz NOT NULL DEFAULT now(),
  deleted_at       timestamptz,
  CONSTRAINT ck_lot_remaining CHECK                            -- BR14
    (remaining_qty BETWEEN 0 AND initial_qty),
  CONSTRAINT ck_lot_dates CHECK (expires_on >= purchased_on),  -- BR11
  CONSTRAINT ck_lot_lifecycle CHECK                            -- BR14
    ((lifecycle = 'active') = (remaining_qty > 0)));
CREATE INDEX ix_lots_fefo ON pantry.food_lots       -- FEFO (BR15)
  (space_id, food_id, expires_on, id) WHERE lifecycle = 'active';
CREATE INDEX ix_lots_expiry ON pantry.food_lots (expires_on)  -- UC14
  WHERE lifecycle = 'active';
\end{Verbatim}

\texttt{ck\_lot\_lifecycle} diễn đạt bất biến của \br{BR14}: lô \texttt{active} khi và chỉ khi còn lượng, nên P04 đặt lượng còn bằng 0 phải đồng thời chuyển lô sang \texttt{consumed}. \texttt{ix\_lots\_fefo} có thứ tự cột trùng điều kiện lọc và \texttt{ORDER BY expires\_on, id}, nên câu \texttt{SELECT ... WHERE expires\_on >= \$today ... FOR UPDATE} của FEFO duyệt chỉ mục đúng thứ tự khóa dòng mà không sắp xếp (Hình~\ref{fig:seq-xac-nhan-nau}). \texttt{ix\_lots\_expiry} phục vụ bộ điều phối nhắc hạn dùng: lọc lô có \texttt{expires\_on} $\le$ \texttt{\$today} + 7 (ngưỡng $N$ lớn nhất) rồi nối \texttt{space\_members} để ra danh sách người cần nhắc.

{\footnotesize\setlength{\tabcolsep}{4pt}\renewcommand{\arraystretch}{1.1}
\begin{longtable}{|L{2.75cm}|L{7.3cm}|L{2.0cm}|L{\dimexpr\textwidth-8\tabcolsep-5\arrayrulewidth-2.75cm-7.3cm-2.0cm\relax}|}
\caption{Chỉ mục và ràng buộc chính}\label{tab:c3-chi-muc}\\
\hline
\hd{Bảng} & \hd{Chỉ mục, ràng buộc} & \hd{Loại} & \hd{Phục vụ}\\ \hline
\endfirsthead
\multicolumn{4}{l}{\footnotesize\itshape Bảng \thetable\ (tiếp theo)}\\ \hline
\hd{Bảng} & \hd{Chỉ mục, ràng buộc} & \hd{Loại} & \hd{Phục vụ}\\ \hline
\endhead
\path{users} & \texttt{uq\_users\_email} \texttt{(email)}; CK \texttt{char\_length(email) <= 100} & UQ; CK & \br{BR01}\\ \hline
\path{space_members} & \texttt{uq\_space\_leader} \texttt{(space\_id) WHERE role = 'leader'} & UQ từng phần & \br{BR05}\\ \hline
\path{space_members} & \texttt{uq\_member\_family} \texttt{(user\_id) WHERE space\_kind = 'family'}, \texttt{uq\_member\_personal} tương tự với \texttt{'personal'}; FK \texttt{(space\_id, space\_kind)} tới \texttt{spaces (id, kind)} & UQ từng phần; FK ghép & \br{BR04}\\ \hline
\path{invite_codes} & UQ \texttt{(space\_id) WHERE status = 'active'}; UQ \texttt{(code)}; CK \texttt{code \textasciitilde{} '\textasciicircum[A-HJ-NP-Z2-9]\{8\}\$'} & UQ từng phần; CK & \br{BR06}\\ \hline
\path{foods} & UQ \texttt{(category\_id, lower(name)) WHERE scope = 'global'} và \texttt{(owner\_space\_id, category\_id, lower(name)) WHERE scope = 'space'}; CK \texttt{(scope = 'global') = (owner\_space\_id IS NULL)}; CK \texttt{merged\_into IS NULL OR status = 'inactive'}; \texttt{GIN (f\_unaccent(name) gin\_trgm\_ops)} & UQ từng phần; CK; GIN & \br{BR22}, \uc{UC23}; tìm dự phòng\\ \hline
\path{units} & UQ \texttt{(dimension) WHERE is\_base}; CK \texttt{NOT is\_base OR factor\_to\_base = 1}; trigger chặn đổi \texttt{dimension} khi đang được tham chiếu & UQ từng phần; CK & \br{BR16}\\ \hline
\path{shopping_items} & \texttt{(list\_id, position)}; \texttt{(assignee\_id) WHERE status = 'assigned'}; \texttt{(space\_id, food\_id) WHERE status IN ('to\_buy', 'assigned', 'purchased')} & B-tree từng phần & \br{BR10}; hàng đang chờ mua của \br{BR20}\\ \hline
\path{food_lots} & \texttt{ix\_lots\_fefo}, \texttt{ix\_lots\_expiry}; ba CK (Mã nguồn~\ref{code:c3-ddl-food-lots}) & Từng phần; CK & \br{BR15}, \uc{UC14}, \br{BR14}\\ \hline
\path{stock_movements} & \texttt{(space\_id, occurred\_at)}; CK \texttt{reason}, \texttt{meal\_entry\_id} theo \texttt{kind}; CK \texttt{quantity <> 0 AND (kind = 'adjust' OR quantity > 0)} & B-tree; CK & \uc{UC21} và \uc{UC13}\\ \hline
\path{recipes} & CK \texttt{(kind = 'template') = (space\_id IS NULL)}; CK \texttt{kind = 'template' OR status = 'published'}; \texttt{(space\_id) WHERE deleted\_at IS NULL} & CK; B-tree từng phần & \br{BR17}, \uc{UC23}\\ \hline
\path{notifications} & UQ \texttt{(dedupe\_key)}; \texttt{(user\_id, created\_at DESC)} & UQ; B-tree & \br{BR13}, \uc{UC14}\\ \hline
\path{idempotency_keys} & PK \texttt{(user\_id, key)}; CK \texttt{status = 'in\_progress' OR response\_body IS NOT NULL} & PK; CK & \nfr{NFR11}\\ \hline
\path{space_changes} & PK \texttt{(space\_id, seq)}; \texttt{BRIN (changed\_at)} & B-tree; BRIN & Y01 đọc theo con trỏ; G03\\ \hline
\path{exports} & \texttt{(space\_id, requested\_by, created\_at DESC)} & B-tree & \uc{UC22}\\ \hline
\path{daily_space_stats} & UQ \texttt{(space\_id, stat\_date, category\_id, dimension)} & UQ & \texttt{REFRESH CONCURRENTLY}\\ \hline
Bảng tham chiếu danh mục & Mọi FK tới \texttt{food\_categories}, \texttt{units}, \texttt{foods} là \texttt{ON DELETE RESTRICT} & FK & \br{BR22}\\ \hline
\end{longtable}
}

Các ràng buộc khai báo tạo phòng thủ nhiều lớp (defense in depth) cho \br{BR04}--\br{BR06}, \br{BR11}, \br{BR14}. Mọi thao tác đổi thành viên, vai trò hoặc mã mời của một không gian trước hết lấy khóa tư vấn \path{pg_advisory_xact_lock(hashtext(space_id::text))}, nên các yêu cầu cùng không gian được tuần tự hóa và phép đếm thành viên so với \texttt{group.maxMembers} của G07 luôn chính xác, điều không ràng buộc khai báo nào diễn đạt được. Khóa tư vấn không che được hai yêu cầu của cùng một người tham gia hai nhóm khác nhau cùng lúc, vì chúng giữ hai khóa khác nhau; khi đó \texttt{uq\_member\_family} chặn yêu cầu commit sau. Chỉ mục một trưởng nhóm không thể trì hoãn (deferrable) nên chuyển quyền \uc{UC07} phải hạ vai trò người cũ trước rồi mới nâng người mới trong cùng giao dịch. Ứng dụng vẫn kiểm tra trước để trả đúng thông điệp (Phụ lục~\ref{pl:thong-diep}); khi hai yêu cầu tranh nhau, lỗi của PostgreSQL được dịch theo SQLSTATE và \emph{tên ràng buộc} chứ không theo SQLSTATE đơn thuần, vì cùng mã 23505 có thể là dữ liệu đã tồn tại (\texttt{uq\_users\_email}, trùng khóa chính do thiết bị sinh: 4003), vi phạm quy tắc (\texttt{uq\_member\_family}: 4004) hay xung đột đồng thời (\texttt{uq\_space\_leader}: 4002). Ánh xạ đầy đủ, gồm 23514 của \texttt{ck\_lot\_dates}, \texttt{ck\_lot\_remaining} và 23503 khi xóa mục danh mục đang được tham chiếu (D05, D06 trả 4004 kèm \br{BR22}), nằm ở Bảng~\ref{tab:c3-anh-xa-loi}; đó là lý do mọi ràng buộc đều được đặt tên tường minh.

\subsection{Thiết kế dữ liệu Redis, Elasticsearch và MinIO}\label{ssec:du-lieu-phu}

Kho phụ không được thành nguồn sự thật thứ hai; khóa Redis~\cite{redis} mang tiền tố môi trường (ví dụ \texttt{prod:}) và có thời hạn sống (Bảng~\ref{tab:c3-khoa-redis}). Bộ đệm luôn bị xóa bởi móc sau commit của giao dịch làm thay đổi dữ liệu nguồn, còn TTL chỉ giới hạn thời gian sai lệch nếu một lần xóa bị bỏ lỡ.

{\footnotesize\setlength{\tabcolsep}{4pt}\renewcommand{\arraystretch}{1.1}
\begin{longtable}{|L{5.2cm}|L{2.3cm}|L{1.7cm}|L{\dimexpr\textwidth-8\tabcolsep-5\arrayrulewidth-5.2cm-2.3cm-1.7cm\relax}|}
\caption{Các mẫu khóa Redis}\label{tab:c3-khoa-redis}\\
\hline
\hd{Mẫu khóa} & \hd{Kiểu Redis} & \hd{TTL} & \hd{Mục đích}\\ \hline
\endfirsthead
\multicolumn{4}{l}{\footnotesize\itshape Bảng \thetable\ (tiếp theo)}\\ \hline
\hd{Mẫu khóa} & \hd{Kiểu Redis} & \hd{TTL} & \hd{Mục đích}\\ \hline
\endhead
\texttt{rl:<luật>:<chủ thể>} & string (GCRA) & Theo luật & Giới hạn tần suất, ví dụ \path{rl:login:<ip>}, \path{rl:join:<userId>}; ngưỡng ở mục~\ref{sec:bao-mat}\\ \hline
\path{otp:cooldown:<email>:<purpose>}; \path{otp:hourly:<email>} & string, \texttt{SET NX}, \texttt{INCR} & 60 s; 1 giờ & Chờ 60 s; tối đa 5 lần mỗi giờ (\br{BR02})\\ \hline
\path{auth:ver:<userId>} & string & 30 phút & Bộ đệm \texttt{token\_version} (\br{BR03})\\ \hline
\path{member:<spaceId>:<userId>} & string & 60 s & Vai trò thành viên; xóa khi đổi thành viên, vai trò\\ \hline
\path{suggest:<spaceId>} & hash (trường theo người, bữa, khẩu phần) & 10 phút & Gợi ý món (\uc{UC18}); xóa khi \texttt{lot.changed}, \texttt{recipe.changed}, \texttt{config.changed} khóa \texttt{suggest.*}, khi A13 đổi $N$\\ \hline
\path{cfg:all} & hash & 5 phút & Cấu hình hệ thống; xóa ngay khi có thông điệp kênh \texttt{cfg}\\ \hline
\path{catalog:snapshot:<etag>} & string (JSON) & 24 giờ & Danh mục kèm ETag (Q03)\\ \hline
\path{space:<id>}, \path{user:<id>}, \path{cfg} & Kênh Pub/Sub & -- & Gợi ý thay đổi; sự kiện cá nhân, đóng socket; xóa bộ đệm cấu hình\\ \hline
\end{longtable}
}

Elasticsearch~\cite{elasticsearch} giữ \texttt{foods\_v1}, \texttt{recipes\_v1} sau bí danh, cập nhật bởi hàng đợi \texttt{index} (Hình~\ref{fig:seq-tim-kiem-chi-muc}); ánh xạ đầy đủ ở mục~\ref{ssec:ky-thuat}, các trường ở Bảng~\ref{tab:c3-truong-es}. Hai chỉ mục chứa cả danh mục chung lẫn thực phẩm và công thức riêng của các không gian, nên các trường phạm vi \texttt{scope}, \texttt{owner\_space\_id}, \texttt{kind}, \texttt{status}, \texttt{space\_id} là bắt buộc trong mọi tài liệu và mọi truy vấn. Tài liệu của một không gian bị xóa khi nhóm giải tán (tác vụ \texttt{PurgeSpace}) hoặc khi tài khoản sở hữu không gian cá nhân bị xóa hẳn (\nfr{NFR18}); thực phẩm bị gộp được cập nhật cùng tác vụ xử lý \texttt{food.merged}.

{\footnotesize\setlength{\tabcolsep}{4pt}\renewcommand{\arraystretch}{1.1}
\begin{longtable}{|L{1.45cm}|L{5.6cm}|L{4.2cm}|L{\dimexpr\textwidth-8\tabcolsep-5\arrayrulewidth-1.45cm-5.6cm-4.2cm\relax}|}
\caption{Các trường của chỉ mục Elasticsearch}\label{tab:c3-truong-es}\\
\hline
\hd{Chỉ mục} & \hd{Trường} & \hd{Kiểu, analyzer} & \hd{Mục đích}\\ \hline
\endfirsthead
\multicolumn{4}{l}{\footnotesize\itshape Bảng \thetable\ (tiếp theo)}\\ \hline
\hd{Chỉ mục} & \hd{Trường} & \hd{Kiểu, analyzer} & \hd{Mục đích}\\ \hline
\endhead
\texttt{foods} & \path{name}, \path{aliases} (kèm \path{prefix}, \path{exact}) & text: \texttt{vi\_base} khi ghi, \texttt{vi\_fold} khi tìm; edge\_ngram 2--15; giữ dấu & Tìm không dấu, đồng nghĩa, lỗi gõ; gợi ý tên\\ \hline
\texttt{foods} & \path{category_id}, \path{scope}, \path{owner_space_id}, \path{status}, \path{updated_at} & keyword; date & Lọc \texttt{scope = global} $\lor$ \texttt{owner\_space\_id = sid}, \texttt{status = active}; đối soát hằng tuần\\ \hline
\texttt{recipes} & \path{name}, \path{ingredient_names}; \path{kind}, \path{status}, \path{space_id}, \path{food_ids} & text như \texttt{foods}; keyword & Tìm theo tên, nguyên liệu; mẫu đã công bố hoặc của không gian (\br{BR23})\\ \hline
\end{longtable}
}

MinIO~\cite{minio} dùng bucket riêng tư, khóa đối tượng chứa UUID nên nội dung không đổi và ảnh được lưu đệm theo khóa (Bảng~\ref{tab:c3-bucket}).

{\footnotesize\setlength{\tabcolsep}{4pt}\renewcommand{\arraystretch}{1.1}
\begin{longtable}{|L{1.8cm}|L{5.6cm}|L{4.2cm}|L{\dimexpr\textwidth-8\tabcolsep-5\arrayrulewidth-1.8cm-5.6cm-4.2cm\relax}|}
\caption{Các bucket MinIO}\label{tab:c3-bucket}\\
\hline
\hd{Bucket} & \hd{Mẫu khóa đối tượng} & \hd{Quyền truy cập} & \hd{Vòng đời}\\ \hline
\endfirsthead
\multicolumn{4}{l}{\footnotesize\itshape Bảng \thetable\ (tiếp theo)}\\ \hline
\hd{Bucket} & \hd{Mẫu khóa đối tượng} & \hd{Quyền truy cập} & \hd{Vòng đời}\\ \hline
\endhead
\texttt{media} & \path{<loại>/<ownerId>/<uuid>.jpg} (loại theo \texttt{purpose}: avatars, spaces, recipes, foods); biến thể \path{thumb.webp} & URL ký sẵn: PUT 5 phút, $\le$ 5 MB; GET ngắn hạn & \texttt{pending} quá 24 giờ bị dọn\\ \hline
\texttt{exports} & \path{<spaceId>/<exportId>.pdf} hoặc \path{.csv}; \texttt{exportId} là \texttt{report.exports.id} & URL GET ký sẵn 15 phút, chỉ cho người yêu cầu (\uc{UC22}) & Xóa sau 7 ngày, trùng \texttt{expires\_at}\\ \hline
Sao lưu & Kho pgBackRest; bản sao bucket \texttt{media} & Tài khoản sao lưu & 7 bản (\nfr{NFR13})\\ \hline
\end{longtable}
}

\subsection{Di trú lược đồ và vòng đời dữ liệu}\label{ssec:di-tru}

Lược đồ được quản lý bằng goose: mỗi thay đổi là một tệp SQL có số phiên bản với phần \texttt{-- +goose Up} và \texttt{Down}, chạy bằng \texttt{goose up} (kèm di trú River) trước khi cuốn chiếu API (mục~\ref{sec:trien-khai}). Vì API cũ và mới cùng chạy khi cuốn chiếu và ứng dụng di động chậm cập nhật, thay đổi phá vỡ tương thích theo chiến lược mở rộng -- thu hẹp (expand--contract): thêm cấu trúc mới và ghi song song, chuyển mã sang đọc dạng mới, chỉ xóa dạng cũ khi phiên bản ứng dụng tối thiểu (bản cũ hơn nhận mã 4260) không còn dùng. Ví dụ, đổi tên một cột được thực hiện qua ba bản phát hành: thêm cột mới và ghi cả hai cột, chuyển mã đọc sang cột mới và nạp lại dữ liệu cũ bằng tác vụ theo lô, rồi mới bỏ cột cũ. Chỉ mục được tạo bằng \texttt{CREATE INDEX CONCURRENTLY}, ràng buộc thêm \texttt{NOT VALID} rồi \texttt{VALIDATE} để không khóa bảng lâu. Dữ liệu tạm có thời hạn lưu giữ (\nfr{NFR18}), dọn hằng đêm trên hàng đợi \texttt{maintenance} (Bảng~\ref{tab:c3-luu-giu}; mốc ngoài yêu cầu là giả định thiết kế).

{\footnotesize\setlength{\tabcolsep}{4pt}\renewcommand{\arraystretch}{1.1}
\begin{longtable}{|L{4.6cm}|L{3.0cm}|L{\dimexpr\textwidth-6\tabcolsep-4\arrayrulewidth-4.6cm-3.0cm\relax}|}
\caption{Thời hạn lưu giữ và tác vụ dọn dẹp}\label{tab:c3-luu-giu}\\
\hline
\hd{Dữ liệu} & \hd{Thời hạn lưu giữ} & \hd{Tác vụ dọn dẹp}\\ \hline
\endfirsthead
\multicolumn{3}{l}{\footnotesize\itshape Bảng \thetable\ (tiếp theo)}\\ \hline
\hd{Dữ liệu} & \hd{Thời hạn lưu giữ} & \hd{Tác vụ dọn dẹp}\\ \hline
\endhead
OTP; khóa lũy đẳng; phiên đã thu hồi & +1 ngày; 7 ngày; 30 ngày & Xóa theo thời điểm hết hạn; khóa \texttt{in\_progress} không tồn đọng vì rollback xóa cùng giao dịch\\ \hline
Nhật ký thay đổi; bản ghi xóa mềm & 90 ngày & Xóa; con trỏ cũ hơn: mã 4100\\ \hline
Thông báo; bản ghi gửi & 90 ngày; 30 ngày & Xóa hằng đêm\\ \hline
Mã mời hết hạn; nhật ký quản trị & Giữ lại; tối thiểu 12 tháng & Chuyển \texttt{expired}; không xóa\\ \hline
Tệp xuất và \texttt{report.exports}; ảnh tải dở; tác vụ River & 7 ngày; 24 giờ; 1--7 ngày & Vòng đời bucket và xóa bản ghi quá \texttt{expires\_at}; dọn \texttt{pending}; River tự dọn\\ \hline
Tài khoản đã xóa (A16) & 30 ngày & Ẩn danh email, họ tên; xóa ảnh, dữ liệu không gian cá nhân, tài liệu ES; giữ đóng góp trong nhóm\\ \hline
Nhóm giải tán (G10, \uc{UC07}) & Từ lúc ghi \texttt{dissolved\_at} & \texttt{PurgeSpace} xóa mềm dữ liệu chung, xóa tài liệu ES\\ \hline
\texttt{daily\_space\_stats} & Làm mới hằng đêm & \texttt{REFRESH MATERIALIZED VIEW CONCURRENTLY}\\ \hline
\end{longtable}
}

Xóa tài khoản không xóa cứng dòng \texttt{users}, vì dữ liệu nhóm tham chiếu tới người tạo, người mua và người nấu qua khóa ngoại. Khi A16 thành công, giao dịch ghi \texttt{deleted\_at}, thu hồi mọi phiên, tăng \texttt{token\_version} và xóa thiết bị, nên tài khoản không đăng nhập được ngay. Sau 30 ngày, tác vụ nền thay email bằng một giá trị duy nhất không định danh (giữ được \texttt{uq\_users\_email}), đặt họ tên thành “Người dùng đã xóa”, vô hiệu băm mật khẩu, rồi xóa ảnh, dữ liệu của không gian cá nhân và tài liệu Elasticsearch tương ứng; dữ liệu người đó đóng góp trong nhóm gia đình được giữ nguyên với tên hiển thị mới (\nfr{NFR18}).

% ---------------------------------------------------------------------
%  3.6. Thiết kế API Backend
% ---------------------------------------------------------------------
\section{Thiết kế API Backend}\label{sec:api}

Giao diện lập trình ứng dụng (API) là hợp đồng duy nhất giữa ứng dụng Flutter và backend: mọi nghiệp vụ ở Chương~\ref{chap:dac-ta}, kể cả thao tác được phát lại sau khi mất mạng, đều đi qua đây. Thiết kế phải thỏa bốn ràng buộc: REST/JSON với đường dẫn \texttt{/it4788/} và bộ mã phản hồi thống nhất (\nfr{NFR21}), kiểm tra không gian và vai trò ở mọi API (\nfr{NFR16}), gửi lại yêu cầu không gây ghi trùng (\nfr{NFR11}) và có tài liệu OpenAPI (\nfr{NFR20}).

% =====================================================================
\subsection{Quy ước thiết kế API}\label{ssec:quy-uoc-api}

API được thiết kế theo tài nguyên (resource-oriented), tên tài nguyên là danh từ số nhiều nối bằng dấu gạch ngang. Dữ liệu của một không gian (\br{BR04}) nằm dưới \path{/spaces/{sid}}: một middleware xác nhận tư cách thành viên cho mọi API bên dưới và repository luôn lọc theo \texttt{space\_id} lấy từ đường dẫn, chặn lỗi truy cập đối tượng trái phép (BOLA) đứng đầu danh sách rủi ro của OWASP~\cite{owaspapi}. Phương thức giữ ngữ nghĩa của RFC 9110~\cite{rfc9110}: \texttt{GET} an toàn, \texttt{PUT} thay thế toàn bộ, \texttt{PATCH} sửa một phần theo JSON Merge Patch, \texttt{DELETE} xóa mềm và để lại bản ghi bia mộ (tombstone) cho đồng bộ. Chuyển trạng thái có điều kiện như đánh dấu đã mua hay xác nhận đã nấu được mô hình hóa thành tài nguyên lệnh với \texttt{POST} (\path{.../purchase}, \path{.../cook}) thay vì \texttt{PATCH} một trường trạng thái, để mỗi điều kiện bảo vệ có một điểm vào và được kiểm thử riêng. Lệnh mà use case có bước xem trước để người dùng xác nhận (G07, S15) nhận cờ \texttt{dryRun}: lần gọi \texttt{true} chỉ đọc và trả kết quả dự kiến, lần gọi xác nhận kiểm tra lại toàn bộ trong giao dịch ghi.

API được phục vụ tại \path{https://api.dicho.example/it4788/v1} (tên miền ví dụ; local, staging, production chỉ khác tên máy chủ), chỉ qua HTTPS với TLS từ 1.2 (\nfr{NFR14}), thân yêu cầu là JSON UTF-8 tối đa 1~MB (vượt nhận 413, mã 2002); ảnh được tải thẳng lên kho đối tượng bằng URL ký sẵn. Phân đoạn \texttt{v1} tồn tại vì ứng dụng trên kho ứng dụng được cập nhật chậm, máy chủ phải phục vụ nhiều thế hệ ứng dụng cùng lúc. Thay đổi bổ sung như API mới hay trường tùy chọn được giữ trong \texttt{v1} với điều kiện ứng dụng bỏ qua trường lạ; chỉ thay đổi phá vỡ hợp đồng mới sinh \texttt{v2} (mục~\ref{ssec:openapi}). Bảng~\ref{tab:c3-header} liệt kê các header dùng chung.

{\scriptsize\renewcommand{\arraystretch}{1.05}
\begin{longtable}{|L{3.3cm}|C{1.25cm}|L{3.2cm}|L{\dimexpr\textwidth-8\tabcolsep-5\arrayrulewidth-3.3cm-1.25cm-3.2cm\relax}|}
\caption{Các header HTTP dùng chung}\label{tab:c3-header}\\
\hline
\hd{Header} & \hd{Bắt buộc} & \hd{Áp dụng} & \hd{Ý nghĩa}\\ \hline
\endfirsthead
\multicolumn{4}{l}{\footnotesize\itshape Bảng \thetable\ (tiếp theo)}\\ \hline
\hd{Header} & \hd{Bắt buộc} & \hd{Áp dụng} & \hd{Ý nghĩa}\\ \hline
\endhead
\multicolumn{4}{|l|}{\hs{Header của yêu cầu}}\\ \hline
\texttt{Authorization} & Có & Trừ A01--A05, A07, A08 & \texttt{Bearer <JWT>} sống 30 phút; sai trả 3001, hết hạn trả 3002\\ \hline
\texttt{Accept-Language} & Không & Mọi API & \texttt{vi} hoặc \texttt{en}, chọn ngôn ngữ của thông điệp\\ \hline
\texttt{X-App-Version}, \texttt{X-Platform} & Có & Mọi API & Phiên bản (\texttt{1.4.2}) và nền tảng (\texttt{android}, \texttt{ios}) của ứng dụng\\ \hline
\texttt{Idempotency-Key} & Có & \texttt{POST} ghi, trừ A01--A05, A07, A08 & UUID của một thao tác, giữ nguyên khi gửi lại; G07, S15 chỉ cần khi \texttt{dryRun = false}\\ \hline
\texttt{If-Match} & Có & \texttt{PATCH}, \texttt{PUT} trên thực thể có cột \texttt{version}; M06 & \texttt{"<version>"} của thực thể mà ứng dụng đã thấy; không dùng cho R06, A13, D09\\ \hline
\texttt{X-Client-Time} & Không & Lệnh tạo, sửa phát lại từ hàng đợi ngoại tuyến & Thời điểm thao tác trên thiết bị (RFC 3339); trong 7 ngày và không sau giờ máy chủ thì BR08, BR11, BR18 được xét theo ngày của giá trị này, ngoài khoảng đó trả 422, mã 4004\\ \hline
\texttt{X-Request-Id} & Không & Mọi API & Mã đối chiếu nhật ký, được trả lại trong phản hồi\\ \hline
\multicolumn{4}{|l|}{\hs{Header của phản hồi}}\\ \hline
\texttt{ETag} & -- & \texttt{GET} một thực thể; Q03 & \texttt{"<version>"}; Q03 dùng với \texttt{If-None-Match} để trả 304\\ \hline
\texttt{Idempotent-Replayed} & -- & Phản hồi phát lại & \texttt{true} khi phản hồi lấy từ bản đã lưu\\ \hline
\texttt{Retry-After} & -- & 423, 429, 503; 409 của khóa lũy đẳng đang xử lý & Số giây cần chờ trước khi gửi lại\\ \hline
\texttt{Deprecation}, \texttt{Sunset} & -- & Phiên bản sắp ngừng & Báo phiên bản lỗi thời và thời điểm ngừng\\ \hline
\texttt{Content-Language} & -- & Phản hồi có thông điệp & Ngôn ngữ thực tế của \texttt{message}\\ \hline
\end{longtable}
}

Ứng dụng có phiên bản thấp hơn mức tối thiểu của nền tảng nhận 426, mã 4260, trước cả bước xác thực; đây là van an toàn khi một bản đã phát hành có lỗi nghiêm trọng, vì máy chủ không sửa được mã đã cài trên điện thoại. Mọi phản hồi dùng cấu trúc bao (envelope) ở Mã nguồn~\ref{code:c3-json-envelope}: mã HTTP phục vụ hạ tầng và chính sách thử lại, còn \texttt{code} bốn chữ số mang ngữ nghĩa nghiệp vụ để ứng dụng rẽ nhánh.

\apiNeed{3.4cm}\codecaption{Cấu trúc bao của phản hồi thành công và phản hồi lỗi}{code:c3-json-envelope}
\begin{Verbatim}[fontsize=\scriptsize]
HTTP/1.1 200 OK
{ "code": 1000, "message": "Thành công",
  "data": { "items": [ { "id": "0192f3c4...", "name": "Đi chợ tuần này" } ] },
  "meta": { "nextCursor": "eyJkIjoiMjAyNi0xMC0yNiJ9",
            "serverTime": "2026-10-25T03:15:42Z" } }
HTTP/1.1 400 Bad Request
{ "code": 2002, "message": "Số lượng phải là số dương, tối đa hai chữ số...",
  "errors": [ { "field": "items[1].quantity", "rule": "BR14", ... } ],
  "traceId": "4bf92f3577b34da6a3ce929d0e0e4736" }
\end{Verbatim}

Mảng \texttt{errors} liệt kê lỗi theo trường với đường dẫn JSON để hiển thị ngay dưới ô nhập (\nfr{NFR03}), \texttt{rule} ghi mã quy tắc bị vi phạm; với xung đột 4002, phần tử có thêm \texttt{current} là trạng thái hiện tại trên máy chủ. \texttt{traceId} là mã vết OpenTelemetry~\cite{opentelemetry}, \texttt{serverTime} giúp ứng dụng hiệu chỉnh lệch đồng hồ khi tính hạn dùng. Thông điệp được chọn theo \texttt{Accept-Language} từ bộ thông điệp Việt, Anh của Phụ lục~\ref{pl:thong-diep}, nên không có chuỗi viết cứng (\nfr{NFR07}).

Danh sách phân trang theo con trỏ (keyset): \texttt{limit} mặc định 20, tối đa 100; \texttt{cursor} mã hóa khóa sắp xếp và \texttt{id} của phần tử cuối, nên không lặp hay sót khi thành viên khác chèn dữ liệu và chi phí không tăng theo số trang như \texttt{OFFSET}. Khóa JSON viết camelCase; \texttt{id} là UUIDv7~\cite{rfc9562}, có thể do thiết bị sinh khi ngoại tuyến, id trùng bản ghi đã có nhận 409, mã 4003. Thời điểm theo RFC 3339 ở UTC; ngày nghiệp vụ dạng \texttt{YYYY-MM-DD} theo múi giờ Asia/Ho\_Chi\_Minh. Số lượng là số JSON tối đa hai chữ số thập phân (\br{BR14}), được đọc nguyên văn thành kiểu thập phân thay vì \texttt{float64}; tiền là số nguyên tính bằng đồng.

Mọi \texttt{POST} ghi của người dùng đã đăng nhập mang \texttt{Idempotency-Key}~\cite{idempotencykey}, một UUID cho mỗi thao tác, giữ nguyên qua mọi lần gửi lại. API công khai A01--A05, A07, A08 không dùng khóa này: chúng đã được giới hạn tần suất theo IP, được bảo vệ bằng OTP hoặc cửa sổ ân hạn của A05, và việc lưu phản hồi của chúng đồng nghĩa lưu mã phiên ở dạng rõ. Câu lệnh đầu tiên của giao dịch ghi là \texttt{INSERT} vào \path{sync.idempotency_keys} một dòng \texttt{status = 'in\_progress'} kèm băm yêu cầu, với \texttt{ON CONFLICT DO NOTHING}. Nếu khóa đã tồn tại, yêu cầu khác băm nhận 400, mã 2002; khóa \texttt{completed} được phát lại đúng phản hồi đã lưu kèm \texttt{Idempotent-Replayed: true}; khóa \texttt{in\_progress}, tức lần gửi trước chưa kết thúc, nhận 409, mã 4002 kèm \texttt{Retry-After} để thử lại sau. Ngay trước \texttt{COMMIT}, dòng khóa được cập nhật thành \texttt{completed} cùng phản hồi; lỗi nghiệp vụ dẫn tới \texttt{ROLLBACK} nên khóa biến mất theo giao dịch, nghĩa là chỉ kết quả thành công được lưu (7 ngày).

Kiểm soát đồng thời dùng khóa lạc quan theo \texttt{version} qua \texttt{ETag} và \texttt{If-Match}, chỉ cho \texttt{PUT}, \texttt{PATCH} trên thực thể có cột \texttt{version} và cho lệnh M06; R06, A13 đặt trạng thái riêng của người gọi còn D09 lưu giá trị cũ, mới trong nhật ký quản trị, nên không cần. Khi không khớp, \texttt{PATCH} thuộc tính vẫn ghi theo hợp nhất từng trường, lệnh đến sau thắng (LWW), và trả \texttt{overwrittenFields}; \texttt{PUT} vốn thay thế toàn bộ như phân công S10 và lệnh phụ thuộc dữ liệu đã xem như M06 nhận 409, mã 4002 (mục~\ref{sec:thiet-ke-chi-tiet}). Lượng còn lại của lô trong P04 là ngoại lệ của LWW: ghi đè sẽ xóa âm thầm một lần dùng thực phẩm của thành viên khác, nên lệch \texttt{If-Match} khi thân có \texttt{remainingQty} trả 409, mã 4002. Nhóm dùng 409 thay cho 412 để mọi xung đột chỉ có một nhánh xử lý ở ứng dụng. Yêu cầu vượt giới hạn tần suất (mục~\ref{sec:bao-mat}) nhận 429, mã 4290, kèm \texttt{Retry-After} và \msg{MSG36}; riêng gửi lại OTP dùng \msg{MSG30}.

Mọi phép so sánh ngày dùng ngày hiện tại theo múi giờ Asia/Ho\_Chi\_Minh do \texttt{platform/clock} cấp và truyền vào câu SQL như tham số \texttt{\$today}, không dùng \texttt{current\_date} của phiên CSDL chạy UTC, vốn lệch một ngày từ 0 giờ đến 7 giờ sáng. Lệnh tạo, sửa phát lại từ hàng đợi ngoại tuyến mang thêm \texttt{X-Client-Time}: nếu giá trị nằm trong 7 ngày và không sau giờ máy chủ, các quy tắc phụ thuộc ngày hiện tại (\br{BR08}, \br{BR11}, \br{BR18}) được xét theo ngày thao tác thật trên thiết bị, để thao tác hợp lệ lúc ngoại tuyến không bị từ chối chỉ vì được đồng bộ sang ngày hôm sau; ngoài khoảng đó máy chủ trả 422, mã 4004 và ứng dụng giữ bản nháp cho người dùng sửa.

% =====================================================================
\subsection{Danh sách mã phản hồi}\label{ssec:ma-phan-hoi}

Bảng~\ref{tab:c3-ma-phan-hoi} liệt kê bộ mã phản hồi dùng chung, mỗi mã gắn cố định với một mã HTTP và với thông điệp ở Phụ lục~\ref{pl:thong-diep}. Với mã không có thông điệp, ứng dụng tự xử lý và \texttt{message} chỉ mang mô tả kỹ thuật cho nhật ký.

{\scriptsize\renewcommand{\arraystretch}{1.05}
\begin{longtable}{|C{1.1cm}|C{1.3cm}|L{\dimexpr\textwidth-8\tabcolsep-5\arrayrulewidth-1.1cm-1.3cm-5.6cm\relax}|L{5.6cm}|}
\caption{Danh sách mã phản hồi của API}\label{tab:c3-ma-phan-hoi}\\
\hline
\hd{Mã} & \hd{HTTP} & \hd{Ý nghĩa} & \hd{Thông điệp}\\ \hline
\endfirsthead
\multicolumn{4}{l}{\footnotesize\itshape Bảng \thetable\ (tiếp theo)}\\ \hline
\hd{Mã} & \hd{HTTP} & \hd{Ý nghĩa} & \hd{Thông điệp}\\ \hline
\endhead
\multicolumn{4}{|l|}{\hs{1xxx -- Thành công}}\\ \hline
1000 & 200, 201 & Thành công & MSG10 (thao tác ghi); MSG11, MSG15, MSG26, MSG31 (thông tin)\\ \hline
1001 & 202 & Đã tiếp nhận, xử lý nền (xuất báo cáo) & --\\ \hline
\multicolumn{4}{|l|}{\hs{2xxx -- Dữ liệu vào}}\\ \hline
2001 & 400 & Thiếu trường bắt buộc & MSG01\\ \hline
2002 & 400, 413 & Giá trị không hợp lệ; thân yêu cầu quá 1~MB (413) & MSG02, MSG14, MSG16, MSG25, MSG27--MSG29\\ \hline
2003 & 422 & Lượng gửi lên vượt lượng còn lại & MSG17\\ \hline
\multicolumn{4}{|l|}{\hs{3xxx -- Xác thực và phân quyền}}\\ \hline
3001 & 401 & Chưa xác thực hoặc mã không hợp lệ & --\\ \hline
3002 & 401 & Mã truy cập hết hạn, ứng dụng làm mới phiên & --\\ \hline
3003 & 401 & Sai email hoặc mật khẩu & MSG05\\ \hline
3004 & 423 & Tài khoản bị khóa tạm do đăng nhập sai & MSG06\\ \hline
3005 & 403 & Tài khoản bị quản trị viên khóa & MSG21\\ \hline
3006 & 400 & OTP sai hoặc hết hạn & MSG04\\ \hline
3007 & 403 & Không có quyền thực hiện & MSG13\\ \hline
3008 & 401 & Mã làm mới bị dùng lại, đã thu hồi họ phiên & --\\ \hline
3009 & 403 & Tài khoản chưa xác thực email; máy chủ gửi lại OTP trong giới hạn BR02 & MSG34\\ \hline
\multicolumn{4}{|l|}{\hs{4xxx -- Tài nguyên và quy tắc nghiệp vụ}}\\ \hline
4001 & 404 & Không tìm thấy tài nguyên & --\\ \hline
4002 & 409 & Xung đột phiên bản, đã được người khác xử lý, hoặc khóa lũy đẳng đang xử lý & MSG35\\ \hline
4003 & 409 & Dữ liệu đã tồn tại & MSG02, MSG03\\ \hline
4004 & 422 & Vi phạm quy tắc nghiệp vụ, kèm mã BR & MSG07--MSG09, MSG12, MSG22, MSG24, MSG32\\ \hline
4100 & 410 & Con trỏ đồng bộ quá cũ, cần tải lại toàn bộ & --\\ \hline
4260 & 426 & Phiên bản ứng dụng quá cũ & --\\ \hline
4290 & 429 & Quá nhiều yêu cầu & MSG36; MSG30 (gửi lại OTP)\\ \hline
\multicolumn{4}{|l|}{\hs{5xxx -- Hệ thống}}\\ \hline
5000 & 500 & Lỗi hệ thống & MSG33\\ \hline
5003 & 503 & Dịch vụ tạm thời không khả dụng, kể cả khi hết lượt thử lại do khóa chết & MSG33\\ \hline
5004 & 504 & Dịch vụ bên ngoài quá thời gian & MSG23 (tra mã vạch)\\ \hline
\end{longtable}
}

Chữ số đầu xác định nhóm và bên cần hành động: 1xxx thành công, trong đó 1001 dành cho xuất báo cáo (B04 trả \texttt{exportId}, ứng dụng hỏi B05 đến khi tệp sẵn sàng), còn các thông điệp thông tin như \msg{MSG15} hay \msg{MSG31} đi cùng mã 1000 vì không phải lỗi; 2xxx là lỗi người dùng sửa được bằng cách đổi dữ liệu vào; 3xxx là danh tính và quyền, phần lớn được ứng dụng xử lý tự động; 4xxx là dữ liệu đúng nhưng không thực hiện được ở trạng thái hiện tại, với 4004 thì \texttt{rule} chỉ rõ quy tắc bị vi phạm; 5xxx là lỗi phía hệ thống, riêng 5004 chỉ dùng cho P08 khi Open Food Facts không trả lời trong 3 giây hoặc bộ ngắt mạch đang mở (sau 5 lỗi liên tiếp, mở 30 giây).

Hai ranh giới được giữ nhất quán. Ranh giới thứ nhất phân 2xxx với 4xxx theo cách người dùng sửa lỗi. Mã 2001, 2002 phát hiện được chỉ từ nội dung yêu cầu; mã 2003 phải đọc lượng còn lại trong cơ sở dữ liệu nhưng vẫn thuộc nhóm 2xxx, vì lượng gửi lên vượt trạng thái mà người dùng đang thấy và chỉ cần nhập lượng nhỏ hơn trên cùng biểu mẫu (P05, P06). M06 cũng có thể thiếu lượng nhưng trả 4002, vì lượng xác nhận được tính từ bản xem trước và tồn kho đã bị thành viên khác đổi sau đó, nên người dùng phải xem lại cả bản xem trước. Ranh giới thứ hai phân 409 với 422: 409 báo yêu cầu đụng độ một thay đổi khác và có thể hợp lệ sau khi tải lại (\msg{MSG35} nêu tên người vừa thay đổi), còn 422 báo yêu cầu sai quy tắc dù gửi bao nhiêu lần.

Tầng dịch vụ kiểm tra trước để trả lỗi theo từng trường, nhưng khi hai yêu cầu tranh nhau thì ràng buộc của PostgreSQL mới là chốt chặn cuối. Vì vậy mọi ràng buộc được đặt tên, và repository dịch lỗi \texttt{pgconn.PgError} theo SQLSTATE và tên ràng buộc thành lỗi miền có mã cố định như Bảng~\ref{tab:c3-anh-xa-loi}. Nhờ đó cùng một vi phạm luôn cho cùng một phản hồi, dù bị chặn ở bước kiểm tra trước hay ở ràng buộc; lỗi CSDL không có trong bảng trả 500, mã 5000, kèm \texttt{traceId} để đối chiếu nhật ký.

{\scriptsize\renewcommand{\arraystretch}{1.05}
\begin{longtable}{|L{4.45cm}|C{1.6cm}|C{1.45cm}|L{2.15cm}|L{\dimexpr\textwidth-10\tabcolsep-6\arrayrulewidth-4.45cm-1.6cm-1.45cm-2.15cm\relax}|}
\caption{Ánh xạ vi phạm ràng buộc cơ sở dữ liệu sang mã phản hồi}\label{tab:c3-anh-xa-loi}\\
\hline
\hd{Ràng buộc, tình huống} & \hd{SQLSTATE} & \hd{HTTP, mã} & \hd{Thông điệp, quy tắc} & \hd{API điển hình}\\ \hline
\endfirsthead
\multicolumn{5}{l}{\footnotesize\itshape Bảng \thetable\ (tiếp theo)}\\ \hline
\hd{Ràng buộc, tình huống} & \hd{SQLSTATE} & \hd{HTTP, mã} & \hd{Thông điệp, quy tắc} & \hd{API điển hình}\\ \hline
\endhead
\texttt{uq\_users\_email}: email đã được đăng ký & 23505 & 409, 4003 & MSG03, BR01 & A01\\ \hline
Khóa chính \texttt{id} do thiết bị sinh đã tồn tại & 23505 & 409, 4003 & -- & P02, S02, S07, M02\\ \hline
Tên trùng trong danh mục (loại, ký hiệu đơn vị, tên thực phẩm trong cùng loại), mã vạch trùng trong cùng phạm vi & 23505 & 409, 4003 & MSG02, BR22 & D03--D06, D11\\ \hline
\texttt{uq\_member\_family}: đã thuộc một nhóm gia đình & 23505 & 422, 4004 & MSG08, BR04 & G02, G07\\ \hline
\texttt{uq\_space\_leader}: nhóm đã có trưởng nhóm khác & 23505 & 409, 4002 & MSG35, BR05 & G09\\ \hline
\texttt{ck\_lot\_dates}: ngày hết hạn trước ngày mua & 23514 & 400, 2002 & MSG14, BR11 & P02, P04, S14\\ \hline
\texttt{ck\_lot\_remaining}: lượng còn lại âm, tức lượng dùng hoặc bỏ vượt lượng còn & 23514 & 422, 2003 & MSG17, BR14 & P05, P06\\ \hline
Khóa ngoại tới mục danh mục đang được tham chiếu khi xóa & 23503 & 422, 4004 & MSG22, BR22 & D05, D06\\ \hline
Chờ khóa quá \texttt{lock\_timeout} hoặc khóa chết, sau ba lần chạy lại giao dịch & 55P03, 40P01 & 503, 5003 & MSG33 & P05, S14, M06\\ \hline
\end{longtable}
}

Bảng~\ref{tab:c3-anh-xa-loi} cho thấy cùng mã 23505 có thể thành ba mã phản hồi tùy ý nghĩa nghiệp vụ. Trùng email hay trùng tên là dữ liệu đã tồn tại (4003); trùng chỉ mục một nhóm gia đình là vi phạm \br{BR04} mà gửi lại bao nhiêu lần vẫn sai (4004); còn trùng chỉ mục một trưởng nhóm chỉ xảy ra khi hai yêu cầu chuyển quyền chạy đồng thời, nên là xung đột (4002). Khóa chính trùng thường do thiết bị gửi lại một lệnh tạo mà mất khóa lũy đẳng, nên ứng dụng xử lý 4003 trên lệnh tạo như đã thành công và tải lại thực thể.

Phía ứng dụng, bộ chặn (interceptor) của dio xử lý mã theo một bảng quyết định. Gặp 3002, ứng dụng gọi A05 (các yêu cầu song song chờ chung một lần làm mới) rồi gửi lại yêu cầu gốc đúng một lần; nhận tiếp 3002, 3001 hoặc 3008 thì xóa phiên; gặp 3009 ở A04 thì chuyển sang màn hình nhập OTP (\uc{UC02}). Gặp 4290, 5003, 5004, 4002 của khóa lũy đẳng đang xử lý hay lỗi mạng, yêu cầu được thử lại theo \texttt{Retry-After} hoặc độ trễ lũy thừa có nhiễu, giữ nguyên \texttt{Idempotency-Key}. Gặp 4002 khác khi phát lại hàng đợi ngoại tuyến, ứng dụng áp trạng thái \texttt{current} và báo người dùng bằng \msg{MSG35}; mã 2xxx và 4004 không bao giờ được thử lại tự động.

% =====================================================================
\subsection{Danh sách API}\label{ssec:danh-sach-api}

Backend cung cấp 89 mã API chia thành mười nhóm theo module; một số mã quản trị gom nhiều phương thức trên cùng tài nguyên. Mã API được ghi vào tài liệu OpenAPI và làm tiền tố mã test case. Bảng~\ref{tab:c3-danh-sach-api} liệt kê toàn bộ API. Đường dẫn ghi sau tiền tố \path{/it4788/v1}; dấu “\ldots” đứng đầu thay cho \path{/spaces/{sid}}, với \texttt{\{sid\}} là định danh không gian; \texttt{[/\{id\}]} là phần dùng khi thao tác trên một phần tử. Cột UC ghi use case mà API hiện thực; năm API hỗ trợ không gắn use case cụ thể (A05, A16, G01, Y01, Y02) ghi quy tắc hoặc yêu cầu phi chức năng mà chúng phục vụ.

{\scriptsize\setlength{\tabcolsep}{4pt}\renewcommand{\arraystretch}{1.0}
\begin{longtable}{|C{0.8cm}|C{1.8cm}|L{5.4cm}|L{\dimexpr\textwidth-10\tabcolsep-6\arrayrulewidth-0.8cm-1.8cm-5.4cm-1.85cm\relax}|C{1.85cm}|}
\caption{Danh sách API của hệ thống}\label{tab:c3-danh-sach-api}\\
\hline
\hd{Mã} & \hd{Method} & \hd{Đường dẫn} & \hd{Chức năng} & \hd{UC}\\ \hline
\endfirsthead
\multicolumn{5}{l}{\footnotesize\itshape Bảng \thetable\ (tiếp theo)}\\ \hline
\hd{Mã} & \hd{Method} & \hd{Đường dẫn} & \hd{Chức năng} & \hd{UC}\\ \hline
\endhead
\multicolumn{5}{|l|}{\hs{Xác thực và tài khoản (module identity)}}\\ \hline
A01 & POST & \apiPath{/auth/register} & Đăng ký, gửi OTP & UC01\\ \hline
A02 & POST & \apiPath{/auth/otp/verify} & Xác thực OTP; trả \texttt{resetToken} khi khôi phục & UC02, UC04\\ \hline
A03 & POST & \apiPath{/auth/otp/resend} & Gửi lại OTP & UC02\\ \hline
A04 & POST & \apiPath{/auth/login} & Đăng nhập, cấp phiên & UC03\\ \hline
A05 & POST & \apiPath{/auth/refresh} & Làm mới phiên & BR03\\ \hline
A06 & POST & \apiPath{/auth/logout} & Đăng xuất, thu hồi phiên & UC05\\ \hline
A07 & POST & \apiPath{/auth/password/forgot} & Yêu cầu OTP khôi phục & UC04\\ \hline
A08 & POST & \apiPath{/auth/password/reset} & Đặt lại mật khẩu bằng \texttt{resetToken} & UC04\\ \hline
A09 & GET & \apiPath{/me} & Xem hồ sơ & UC06\\ \hline
A10 & PATCH & \apiPath{/me} & Cập nhật hồ sơ & UC06\\ \hline
A11 & PUT & \apiPath{/me/password} & Đổi mật khẩu, cấp phiên mới & UC07\\ \hline
A12 & GET & \apiPath{/me/settings} & Xem thiết lập, khung giờ nhắc & UC08\\ \hline
A13 & PUT & \apiPath{/me/settings} & Lưu thiết lập nhắc nhở & UC08\\ \hline
A14 & POST & \apiPath{/me/devices} & Đăng ký mã thiết bị FCM & UC31\\ \hline
A15 & DELETE & \apiPath{/me/devices/{deviceId}} & Hủy đăng ký thiết bị & UC05\\ \hline
A16 & DELETE & \apiPath{/me} & Xóa tài khoản và dữ liệu & NFR18\\ \hline
A17 & GET & \apiPath{/me/notifications} & Hộp thư thông báo & UC31\\ \hline
A18 & POST & \apiPath{/me/notifications/read} & Đánh dấu đã đọc & UC31\\ \hline
\multicolumn{5}{|l|}{\hs{Không gian và nhóm gia đình (module household)}}\\ \hline
G01 & GET & \apiPath{/spaces} & Các không gian của người dùng & BR04\\ \hline
G02 & POST & \apiPath{/spaces} & Tạo nhóm, sao chép dữ liệu cá nhân & UC09\\ \hline
G03 & GET & \ldots & Nhóm, thành viên, hoạt động gần đây & UC12\\ \hline
G04 & PATCH & \ldots & Đổi tên, ảnh nhóm & UC12\\ \hline
G05 & POST & \ldots\apiPath{/invites} & Tạo mã mời mới & UC10\\ \hline
G06 & DELETE & \ldots\apiPath{/invites/current} & Thu hồi mã mời & UC10\\ \hline
G07 & POST & \apiPath{/spaces/join} & Tham gia nhóm bằng mã mời & UC11\\ \hline
G08 & DELETE & \ldots\apiPath{/members/{userId}} & Xóa thành viên & UC13\\ \hline
G09 & POST & \ldots\apiPath{/leader-transfer} & Chuyển quyền trưởng nhóm & UC14\\ \hline
G10 & POST & \ldots\apiPath{/leave} & Rời nhóm & UC15\\ \hline
\multicolumn{5}{|l|}{\hs{Mua sắm (module shopping)}}\\ \hline
S01 & GET & \ldots\apiPath{/shopping-lists} & Danh sách, lịch sử theo tháng & UC16, UC22, UC23\\ \hline
S02 & POST & \ldots\apiPath{/shopping-lists} & Tạo danh sách & UC16\\ \hline
S03 & GET & \ldots\apiPath{/shopping-lists/{listId}} & Chi tiết danh sách & UC16, UC23\\ \hline
S04 & PATCH & \ldots\apiPath{/shopping-lists/{listId}} & Sửa tên, ngày, ghi chú & UC16\\ \hline
S05 & DELETE & \ldots\apiPath{/shopping-lists/{listId}} & Xóa danh sách & UC16\\ \hline
S06 & POST & \ldots\apiPath{/shopping-lists/{listId}/copy} & Sao chép danh sách & UC18\\ \hline
S07 & POST & \ldots\apiPath{/shopping-lists/{listId}/items} & Thêm mặt hàng, gộp trùng & UC17\\ \hline
S08 & PATCH & \ldots\apiPath{/shopping-items/{itemId}} & Sửa mặt hàng & UC17\\ \hline
S09 & DELETE & \ldots\apiPath{/shopping-items/{itemId}} & Xóa mặt hàng & UC17\\ \hline
S10 & PUT & \ldots\apiPath{/shopping-items/{itemId}/assignee} & Phân công người mua & UC19\\ \hline
S11 & DELETE & \ldots\apiPath{/shopping-items/{itemId}/assignee} & Bỏ phân công & UC19\\ \hline
S12 & POST & \ldots\apiPath{/shopping-items/purchase} & Đánh dấu đã mua & UC20\\ \hline
S13 & POST & \ldots\apiPath{/shopping-items/{itemId}/undo-purchase} & Hoàn tác đã mua & UC20\\ \hline
S14 & POST & \ldots\apiPath{/shopping-items/stock} & Chuyển vào tủ lạnh & UC21\\ \hline
S15 & POST & \ldots\apiPath{/shopping-lists/generate} & Sinh danh sách từ kế hoạch & UC22\\ \hline
\multicolumn{5}{|l|}{\hs{Tủ lạnh (module pantry, catalog)}}\\ \hline
P01 & GET & \ldots\apiPath{/lots} & Danh sách lô, bộ lọc, đếm theo loại & UC24, UC43\\ \hline
P02 & POST & \ldots\apiPath{/lots} & Thêm thực phẩm & UC25\\ \hline
P03 & GET & \ldots\apiPath{/lots/{lotId}} & Chi tiết lô, biến động & UC24\\ \hline
P04 & PATCH & \ldots\apiPath{/lots/{lotId}} & Cập nhật lô & UC27\\ \hline
P05 & POST & \ldots\apiPath{/lots/{lotId}/consume} & Ghi nhận sử dụng (FEFO) & UC28\\ \hline
P06 & POST & \ldots\apiPath{/lots/{lotId}/discard} & Loại bỏ thực phẩm & UC29\\ \hline
P07 & GET & \ldots\apiPath{/lots/expiring} & Lô sắp hoặc đã hết hạn & UC30\\ \hline
P08 & GET & \apiPath{/barcodes/{code}} & Tra cứu mã vạch & UC26\\ \hline
\multicolumn{5}{|l|}{\hs{Công thức nấu ăn (module recipe)}}\\ \hline
R01 & GET & \ldots\apiPath{/recipes} & Bộ sưu tập, công thức mẫu & UC34, UC37\\ \hline
R02 & POST & \ldots\apiPath{/recipes} & Thêm công thức & UC32\\ \hline
R03 & GET & \ldots\apiPath{/recipes/{recipeId}} & Chi tiết, đối chiếu tủ lạnh & UC34\\ \hline
R04 & PUT & \ldots\apiPath{/recipes/{recipeId}} & Cập nhật công thức & UC33\\ \hline
R05 & DELETE & \ldots\apiPath{/recipes/{recipeId}} & Xóa công thức & UC33\\ \hline
R06 & PUT & \apiPath{/me/favorite-recipes/{recipeId}} & Đánh dấu yêu thích & UC35\\ \hline
R07 & DELETE & \apiPath{/me/favorite-recipes/{recipeId}} & Bỏ yêu thích & UC35\\ \hline
R08 & POST & \ldots\apiPath{/recipes/{recipeId}/missing-to-list} & Bổ sung nguyên liệu thiếu & UC36\\ \hline
R09 & POST & \ldots\apiPath{/recipes/{recipeId}/copy} & Sao chép công thức mẫu vào không gian & UC32, UC52\\ \hline
\multicolumn{5}{|l|}{\hs{Kế hoạch bữa ăn (module mealplan)}}\\ \hline
M01 & GET & \ldots\apiPath{/meals} & Lịch bữa ăn theo khoảng ngày & UC38\\ \hline
M02 & POST & \ldots\apiPath{/meals} & Thêm món vào bữa & UC37\\ \hline
M03 & PATCH & \ldots\apiPath{/meals/{mealId}} & Sửa món trong lịch & UC39\\ \hline
M04 & DELETE & \ldots\apiPath{/meals/{mealId}} & Xóa món khỏi lịch & UC39\\ \hline
M05 & GET & \ldots\apiPath{/recipe-suggestions} & Gợi ý món ăn & UC30, UC40\\ \hline
M06 & POST & \ldots\apiPath{/meals/{mealId}/cook} & Xác nhận đã nấu & UC41\\ \hline
\multicolumn{5}{|l|}{\hs{Tìm kiếm và danh mục (module search, catalog)}}\\ \hline
Q01 & GET & \ldots\apiPath{/search} & Tìm thực phẩm, món ăn, bộ lọc & UC42, UC43\\ \hline
Q02 & GET & \ldots\apiPath{/foods/suggest} & Gợi ý tên khi gõ & UC25, UC42\\ \hline
Q03 & GET & \apiPath{/catalog} & Loại, đơn vị, vị trí lưu trữ & UC25, UC43\\ \hline
\multicolumn{5}{|l|}{\hs{Báo cáo (module report)}}\\ \hline
B01 & GET & \ldots\apiPath{/reports/spending} & Báo cáo chi tiêu & UC44\\ \hline
B02 & GET & \ldots\apiPath{/reports/consumption} & Báo cáo tiêu thụ & UC45\\ \hline
B03 & GET & \ldots\apiPath{/reports/waste} & Báo cáo lãng phí & UC46\\ \hline
B04 & POST & \ldots\apiPath{/reports/exports} & Yêu cầu xuất PDF, CSV (202) & UC47\\ \hline
B05 & GET & \ldots\apiPath{/reports/exports/{exportId}} & Trạng thái, URL ký sẵn của tệp & UC47\\ \hline
\multicolumn{5}{|l|}{\hs{Đồng bộ và tệp (module sync, media)}}\\ \hline
Y01 & GET & \ldots\apiPath{/changes} & Kéo thay đổi theo \texttt{seq} & NFR11\\ \hline
Y02 & GET & \apiPath{/ws} & WebSocket gợi ý thay đổi & NFR10\\ \hline
Y03 & POST & \apiPath{/media/upload-urls} & Cấp URL ký sẵn tải ảnh & UC06, UC25, UC32\\ \hline
Y04 & POST & \apiPath{/media/{objectId}/complete} & Xác nhận tải ảnh xong & UC06, UC25, UC32\\ \hline
\multicolumn{5}{|l|}{\hs{Quản trị hệ thống (module admin, catalog, recipe)}}\\ \hline
D01 & GET & \apiPath{/admin/users} & Tra cứu tài khoản & UC48\\ \hline
D02 & PATCH & \apiPath{/admin/users/{userId}} & Khóa, mở khóa; cấp, thu hồi quyền quản trị & UC48\\ \hline
D03 & GET, POST & \apiPath{/admin/foods} & Xem, thêm thực phẩm & UC49\\ \hline
D04 & PUT, DELETE & \apiPath{/admin/foods/{foodId}} & Sửa, ngừng thực phẩm & UC49\\ \hline
D05 & GET, POST, PUT, DELETE & \apiPath{/admin/food-categories[/{id}]} & Quản lý loại mặt hàng & UC50\\ \hline
D06 & GET, POST, PUT, DELETE & \apiPath{/admin/units[/{id}]} & Quản lý đơn vị đo & UC51\\ \hline
D07 & GET, POST, PUT, DELETE & \apiPath{/admin/recipe-templates[/{id}]} & Quản lý công thức mẫu & UC52\\ \hline
D08 & GET & \apiPath{/admin/configs} & Xem tham số hệ thống & UC53\\ \hline
D09 & PUT & \apiPath{/admin/configs} & Cập nhật nhiều tham số trong một giao dịch & UC53\\ \hline
D10 & GET & \apiPath{/admin/audit-logs} & Nhật ký quản trị & UC48, UC53\\ \hline
D11 & POST & \apiPath{/admin/foods/{foodId}/review} & Duyệt hoặc gộp thực phẩm riêng & UC49\\ \hline
\end{longtable}
}

Quyền truy cập được gắn theo nhóm đường dẫn ngay ở bộ định tuyến, với năm mức. API công khai A01--A05, A07, A08 chỉ chịu giới hạn tần suất theo địa chỉ IP. API của người dùng đã đăng nhập gồm phần còn lại của nhóm A cùng G01, G02, G07, R06, R07, P08, Q03, Y02--Y04, chỉ chạm tới dữ liệu của chính người gọi hoặc danh mục dùng chung; R06 và Y03 kiểm tra thêm quyền trên đối tượng đích. API của thành viên không gian là mọi đường dẫn chứa \texttt{\{sid\}}, kể cả Q02 vì gợi ý gồm thực phẩm riêng của không gian; người không phải thành viên nhận 404, mã 4001, thay vì 403 để không xác nhận không gian có tồn tại (\br{BR23}). API của trưởng nhóm gồm G04--G06, G08, G09, S10, S11, còn R04, R05 cho phép người tạo công thức hoặc trưởng nhóm; người thiếu vai trò nhận 403, mã 3007 (\msg{MSG13}). API quản trị D01--D11 yêu cầu \texttt{system\_role = admin}. Quản trị viên vẫn là người dùng nên dùng được không gian của chính mình, nhưng vai trò hệ thống không thay cho tư cách thành viên nên không đọc được dữ liệu của nhóm khác. Y02 nâng cấp lên WebSocket sau khi xác thực; máy chủ gửi \texttt{ping} mỗi 25 giây, đóng kết nối nếu không có \texttt{pong} sau 10 giây và đóng với mã 4401 khi mã truy cập hết hạn. Các điểm cuối vận hành \path{/healthz}, \path{/readyz}, \path{/metrics} chỉ mở trong mạng nội bộ Docker.


\subsection{Đặc tả chi tiết API}\label{ssec:dac-ta-api}

Mục này đặc tả chi tiết hai API tiêu biểu theo đúng mẫu của học phần: A04 Đăng nhập đại diện cho nhóm API an toàn phiên, M06 Xác nhận đã nấu đại diện cho nhóm giao dịch trải nhiều module với trừ kho theo FEFO, lũy đẳng và kiểm soát đồng thời. Mỗi API có một bảng gồm phần đặc tả và phần test case dự kiến, kèm ví dụ yêu cầu, phản hồi. Đặc tả chi tiết của các API có rủi ro thiết kế cao khác, gồm A05, G07, S12 và S14, S15, P05, M05, Y01, cùng bảng tham số và hành vi bổ sung của các API còn lại được trình bày ở Phụ lục~\ref{pl:dac-ta-api}.

% ---------------------------------------------------------------------
\subsubsection{A04 – Đăng nhập}\label{sssec:api-a04}

A04 hiện thực \uc{UC03} (Bảng~\ref{tab:c3-api-a04}, Mã nguồn~\ref{code:c3-api-a04}). Email không tồn tại vẫn được xác minh trên một chuỗi băm giả, còn trạng thái bị quản trị viên khóa (\texttt{disabled}) hay chưa xác thực (\texttt{pending}) chỉ lộ ra sau khi mật khẩu đúng, nên người không biết mật khẩu không dò được tình trạng tài khoản. Ngược lại, khóa tạm do sai mật khẩu (\texttt{locked}) trả 423 trước cả bước xác minh: nếu vẫn so mật khẩu, khác biệt giữa 401 và 423 sẽ cho kẻ dò biết mật khẩu đúng ngay trong thời gian khóa. Mỗi lần đăng nhập thu hồi họ phiên cũ của cùng thiết bị, nên đăng nhập lại sau khi mất phản hồi là vô hại.

\begin{apispec}{Đặc tả và test case dự kiến của API A04 -- Đăng nhập}{tab:c3-api-a04}{tab:c3-tc-a04}
\apirow{URL, quyền}{\texttt{POST} \path{/it4788/v1/auth/login}; công khai, 10 yêu cầu/phút/IP; không dùng \texttt{Idempotency-Key}.}
\apirow{Yêu cầu}{\texttt{email}: chuỗi, bắt buộc, đúng định dạng, $\le$ 100 ký tự; \texttt{password}: chuỗi, bắt buộc, 1--128 ký tự; \texttt{deviceId}: UUID, bắt buộc, sinh khi cài ứng dụng.}
\apirow{Phản hồi}{200, 1000: \texttt{accessToken}, \texttt{tokenType}, \texttt{expiresIn = 1800}, \texttt{refreshToken}, \texttt{refreshExpiresAt}, \texttt{user}, \texttt{lastSpaceId} (không gian mở khi vào ứng dụng).}
\apirow{Mã lỗi}{2001 (MSG01); 2002 (MSG02); 3003 sai email hoặc mật khẩu (MSG05); 3004 với HTTP 423 khi \texttt{locked}, kèm \texttt{Retry-After} (MSG06); 3005 khi \texttt{disabled} (MSG21); 3009 khi \texttt{pending} (MSG34), máy chủ gửi OTP mới nếu còn trong giới hạn của BR02.}
\tchead
\tcrow{TC-A04-01}{Đăng nhập hợp lệ}{Tài khoản \texttt{active}, bộ đếm sai bằng 2, email khác hoa thường}{200, 1000; bộ đếm về 0; thêm một phiên với \texttt{family\_id} mới}
\tcrow{TC-A04-02}{Thiếu trường}{Không có \texttt{password}}{400, 2001, \texttt{field = "password"}; CSDL không đổi}
\tcrow{TC-A04-03}{Biên ngưỡng khóa}{Bộ đếm bằng 3 rồi 4, mật khẩu sai; ngay sau đó mật khẩu đúng}{Lần thứ tư: 401, 3003. Lần thứ năm: 423, 3004, \texttt{locked\_until} được đặt. Mật khẩu đúng sau đó vẫn nhận 423, không tạo phiên}
\tcrow{TC-A04-04}{Không lộ thông tin}{(a) Email chưa đăng ký; (b) tài khoản \texttt{disabled} hoặc \texttt{pending}, mật khẩu sai}{Cả hai: 401, 3003, thân phản hồi giống hệt ca sai mật khẩu thông thường}
\tcrow{TC-A04-05}{Trạng thái tài khoản}{Mật khẩu đúng; (a) \texttt{locked}, \texttt{locked\_until} còn 10 phút; (b) \texttt{disabled}; (c) \texttt{pending}}{(a) 423, 3004 (MSG06), \texttt{Retry-After} 600; (b) 403, 3005 (MSG21); (c) 403, 3009 (MSG34), thêm một dòng \texttt{otp\_challenges}; cả ba không tạo phiên}
\end{apispec}

\apiNeed{3.4cm}\codecaption{Ví dụ yêu cầu và phản hồi của A04}{code:c3-api-a04}
\begin{Verbatim}[fontsize=\scriptsize]
POST /it4788/v1/auth/login                               X-Platform: android
{ "email": "lan.nguyen@example.com", "password": "Bep@2026an",
  "deviceId": "0192f3a1-7c2e-7b10-9a4f-1d2e3f405162" }
HTTP/1.1 200 OK
{ "code": 1000, "message": "Thành công",
  "data": { "accessToken": "eyJhbGciOiJFZERTQSIs...", "expiresIn": 1800,
    "refreshToken": "q3v9Xb7tLk2...wYc8", "lastSpaceId": "0192f0ab-...",
    "user": { "id": "0192f0aa-...", "fullName": "Nguyễn Thị Lan" } } }
\end{Verbatim}

% ---------------------------------------------------------------------
\subsubsection{M06 – Xác nhận đã nấu}\label{sssec:api-m06}

M06 hiện thực \uc{UC41} theo Hình~\ref{fig:seq-xac-nhan-nau} (Bảng~\ref{tab:c3-api-m06}). Bản xem trước do ứng dụng tính từ bản sao cục bộ, nên yêu cầu mang lượng trừ đã xác nhận cho từng nguyên liệu: 0 là lấy ngoài tủ, \texttt{lotIds} chọn rõ lô quá hạn. Nếu tồn kho đã giảm, máy chủ trả 409, mã 4002, kèm lượng mới chứ không phải 2003 như P05, vì lượng xác nhận được tính từ bản xem trước mà tồn kho đã bị người khác đổi sau đó (Mã nguồn~\ref{code:c3-api-m06}). Món đã \texttt{cooked} được xét trước \texttt{If-Match} và nhận 200 với \texttt{alreadyCooked = true}, không trừ kho lần hai. Giao dịch theo thứ tự khóa toàn cục: khóa lũy đẳng, món, các lô theo \texttt{food\_id} rồi \texttt{(expires\_on, id)}, cuối cùng mới tăng \texttt{change\_seq}. Giao dịch đặt \texttt{lock\_timeout} 3 giây và tự chạy lại tối đa ba lần khi gặp khóa chết (40P01) hoặc chờ khóa quá hạn (55P03), hết lượt thì trả 503, mã 5003. Mức cô lập là READ COMMITTED với khóa dòng tường minh, nên lỗi tuần tự hóa 40001 không phát sinh.

\begin{apispec}{Đặc tả và test case dự kiến của API M06 -- Xác nhận đã nấu}{tab:c3-api-m06}{tab:c3-tc-m06}
\apirow{URL, quyền}{\texttt{POST} \path{/it4788/v1/spaces/{sid}/meals/{mealId}/cook}; thành viên không gian; \texttt{Idempotency-Key}, \texttt{If-Match: "<version>"}.}
\apirow{Yêu cầu}{\texttt{servings}: 1--20; \texttt{skipStock}: mặc định \texttt{false}; \texttt{lines} (bắt buộc khi trừ kho), mỗi dòng một nguyên liệu: \texttt{foodId}, \texttt{quantity} ($\ge$ 0, tối đa hai chữ số thập phân), \texttt{unitId}, \texttt{lotIds} (tùy chọn).}
\apirow{Phản hồi}{200, 1000: \texttt{meal} (\texttt{status = cooked}, \texttt{cookedAt}, \texttt{cookedBy}, \texttt{version}), \texttt{deductions} (\texttt{foodId}, \texttt{lotId}, \texttt{quantity}), \texttt{alreadyCooked}.}
\apirow{Mã lỗi}{2001; 2002 (MSG16); 4001; 4002 \texttt{If-Match} không khớp hoặc tồn kho không đủ (MSG35, \texttt{current} chứa lượng tồn mới); 4004 \texttt{rule = "BR18"} khi món chưa đến ngày nấu hoặc không gắn công thức; 5003 hết ba lần chạy lại.}
\tchead
\tcrow{TC-M06-01}{Hợp lệ, một nguyên liệu lấy ngoài}{4 khẩu phần, công thức chuẩn 2; lượng theo bản xem trước, me bằng 0}{200; mỗi lô bị trừ có một dòng \texttt{cook}, không có dòng nào cho me; món \texttt{cooked}, \texttt{version} tăng 1}
\tcrow{TC-M06-02}{Dữ liệu sai, biên ngày}{(a) \texttt{quantity} bằng $-$0,5 hoặc 0,125; (b) \texttt{meal\_date} là ngày mai, đối chứng hôm nay}{(a) 400, 2002; (b) 422, 4004, \texttt{rule = "BR18"}; đối chứng: 200}
\tcrow{TC-M06-03}{Tồn kho giảm đồng thời}{Giữa xem trước và xác nhận, thành viên khác dùng bớt cá lóc}{409, 4002 kèm \texttt{current}; không lô nào bị trừ, món vẫn \texttt{planned}}
\tcrow{TC-M06-04}{Bỏ qua trừ kho}{\texttt{skipStock = true}}{200; món \texttt{cooked}; \texttt{food\_lots}, \texttt{stock\_movements} không đổi}
\tcrow{TC-M06-05}{Món đã thay đổi}{(a) Đã \texttt{cooked} trên thiết bị khác, khóa lũy đẳng mới; (b) \texttt{If-Match: "3"}, món đã ở \texttt{version} 4}{(a) 200, \texttt{alreadyCooked = true}, không trừ kho lần hai; (b) 409, 4002, không ghi}
\end{apispec}

\apiNeed{3.4cm}\codecaption{Ví dụ yêu cầu và phản hồi xung đột của M06}{code:c3-api-m06}
\begin{Verbatim}[fontsize=\scriptsize]
POST /it4788/v1/spaces/0192e8b1-.../meals/0192f5b0-.../cook    If-Match: "3"
{ "servings": 4, "skipStock": false, "lines": [
    { "foodId": "0192a0c2-...", "quantity": 0.8, "unitId": "<kg>" },
    { "foodId": "0192a0f1-...", "quantity": 0, "unitId": "<g>" } ] }
HTTP/1.1 409 Conflict
{ "code": 4002,
  "message": "Dữ liệu vừa được Minh thay đổi. Vui lòng xem lại trước...",
  "errors": [ { "field": "lines[0].quantity", "rule": "BR14",
    "message": "Tồn kho hiện chỉ còn 0,5 kg.",
    "current": { "available": 0.5, "unitId": "<kg>" } } ] }
\end{Verbatim}

\subsection{Tài liệu OpenAPI và sinh mã client}\label{ssec:openapi}

Tệp \path{api/openapi.yaml} theo OpenAPI 3.1~\cite{openapi} là nguồn sự thật của hợp đồng, được viết và duyệt trước khi có mã nguồn (contract-first). Mã máy chủ, mã ứng dụng và tài liệu đều sinh từ tệp này, nên hai phía không thể lệch nhau âm thầm, và phía ứng dụng làm việc được với máy chủ giả lập (mock) sinh từ cùng tệp ngay từ đầu. Mỗi thao tác có \texttt{operationId} mô tả hành động, phần mở rộng \texttt{x-api-id} ghi mã API của Bảng~\ref{tab:c3-danh-sach-api}, và tham chiếu thành phần dùng chung như cấu trúc bao, header \texttt{Idempotency-Key}, kiểu \texttt{Quantity}, như đoạn đặc tả P05 ở Mã nguồn~\ref{code:c3-openapi}.

\apiNeed{3.4cm}\codecaption{Đoạn đặc tả OpenAPI của API P05}{code:c3-openapi}
\begin{Verbatim}[fontsize=\scriptsize]
/spaces/{sid}/lots/{lotId}/consume:
  post:
    operationId: consumeLot
    x-api-id: P05                      # UC28; BR14, BR15, BR16
    security: [ { bearerAuth: [] } ]
    parameters:
      - $ref: '#/components/parameters/SpaceId'
      - $ref: '#/components/parameters/LotId'
      - $ref: '#/components/parameters/IdempotencyKey'
    requestBody:
      required: true
      content:
        application/json:
          schema:
            type: object
            required: [quantity, unitId]
            additionalProperties: false
            properties:
              quantity: { $ref: '#/components/schemas/Quantity' }
              unitId: { type: string, format: uuid }
              strategy: { type: string, enum: [fefo, lot], default: fefo }
    responses:
      '200': { $ref: '#/components/responses/ConsumeResult' }
      '4XX': { $ref: '#/components/responses/Error' }
\end{Verbatim}

Công cụ oapi-codegen sinh cho Go kiểu dữ liệu và giao diện máy chủ dạng chặt (strict) gắn với chi, nên thiếu một API hay trả sai kiểu là lỗi biên dịch; \texttt{additionalProperties: false} chặn gán hàng loạt các trường như \texttt{spaceId}. Phía ứng dụng, openapi-generator với bộ sinh \texttt{dart-dio} tạo gói Dart cho ứng dụng Flutter, trên đó các interceptor ở mục~\ref{ssec:ma-phan-hoi} được cấu hình một lần. Ở stage verify của GitLab CI (mục~\ref{sec:trien-khai}), mã được sinh lại và so với bản đã commit; Spectral kiểm tra các luật riêng như mọi \texttt{POST} ghi ngoài nhóm công khai khai báo \texttt{Idempotency-Key}, mọi \texttt{If-Match} chỉ xuất hiện trên thao tác có trong danh sách thực thể có \texttt{version}, mọi đường dẫn chứa \texttt{\{sid\}} khai báo 404, mọi lỗi dùng cấu trúc bao chung; tệp còn được so với nhánh chính để chặn thay đổi phá vỡ hợp đồng trong \texttt{v1}. Vì các bộ sinh mã chưa hỗ trợ đều mọi cấu trúc riêng của OpenAPI 3.1, nhóm chỉ dùng tập cấu trúc mà cả hai bộ sinh xử lý được.

Khi \texttt{v2} phát hành, \texttt{v1} được giữ song song ít nhất sáu tháng: phản hồi của \texttt{v1} mang \texttt{Deprecation} và \texttt{Sunset}, giám sát thống kê lưu lượng theo phiên bản API và \texttt{X-App-Version}, rồi mức phiên bản ứng dụng tối thiểu được nâng để bản cũ nhận 426 trước ngày gỡ \texttt{v1}. Kiểm kê đầy đủ phiên bản đang mở cũng là biện pháp cho rủi ro API9 của OWASP~\cite{owaspapi}. Tài liệu tương tác Swagger UI sinh từ cùng tệp chỉ mở trên staging trong mạng nội bộ của nhóm, không công bố ở production; như vậy một tệp duy nhất vừa là hợp đồng, vừa sinh mã hai phía và tài liệu, đáp ứng \nfr{NFR20}.

% ---------------------------------------------------------------------
%  3.7. Thiết kế bảo mật
% ---------------------------------------------------------------------
\section{Thiết kế bảo mật}\label{sec:bao-mat}

Dữ liệu của hệ thống mang tính riêng tư của từng gia đình, và một lỗi phân quyền có thể làm lộ dữ liệu của nhiều nhóm cùng lúc. Thiết kế theo nguyên tắc phòng thủ nhiều lớp (defense in depth), lấy OWASP MASVS~\cite{masvs} làm chuẩn phía thiết bị và OWASP API Security Top 10 (2023)~\cite{owaspapi} làm danh mục rủi ro phía máy chủ, đáp ứng \nfr{NFR14}--\nfr{NFR18}.

\subsection{Xác thực và quản lý phiên}\label{ssec:xac-thuc-phien}

Mật khẩu được băm bằng Argon2id~\cite{rfc9106} với tham số OWASP khuyến nghị~\cite{owasppassword} ($m = 19$~MiB, $t = 2$, $p = 1$); chuỗi băm định dạng PHC chứa cả tham số nên có thể nâng tham số bằng cách băm lại ở lần đăng nhập kế tiếp. Email không tồn tại vẫn được xác minh trên một chuỗi băm giả để thời gian phản hồi không lộ tài khoản, khớp thông điệp chung \msg{MSG05}.

Mã truy cập (access token) là JWT~\cite{rfc7519} ký EdDSA trên Ed25519~\cite{rfc8032}, sống 30 phút (\br{BR03}), mang \texttt{sub}, \texttt{sid} (họ phiên), \texttt{ver}, \texttt{role}, \texttt{exp}, \texttt{jti} và \texttt{kid} để xoay khóa; chữ ký bất đối xứng cho phép dịch vụ tách ra sau này xác minh bằng khóa công khai. Mã làm mới (refresh token) là chuỗi ngẫu nhiên 256 bit, chỉ lưu SHA-256 trong \texttt{identity.refresh\_sessions}, sống 30 ngày và được xoay vòng mỗi lần dùng trong một giao dịch khóa dòng phiên. Mã đã xoay bị gửi lại là dấu hiệu bị đánh cắp nên cả họ phiên bị thu hồi (mã 3008); ngoại lệ duy nhất là cửa sổ ân hạn 30 giây cho cùng \texttt{device\_id} khi phản hồi trước bị mất trên mạng di động, lúc đó máy chủ chỉ thu hồi phiên kế nhiệm chưa dùng và cấp cặp mã mới trong cùng họ.

Đăng nhập (A04, \uc{UC02}) kiểm tra trạng thái tài khoản theo một thứ tự có chủ đích. Tài khoản đang khóa tạm (\texttt{locked\_until} còn hiệu lực) bị từ chối với HTTP 423, mã 3004 (\msg{MSG06}) trước khi xác minh mật khẩu, nên kẻ dò không thể thử tiếp trong 15 phút khóa. Sai mật khẩu tăng \texttt{failed\_login\_count}; lần sai thứ năm đặt \texttt{locked\_until} (\br{BR03}). Chỉ khi mật khẩu đúng, máy chủ mới cho biết tài khoản bị quản trị viên khóa (403, mã 3005, \msg{MSG21}) hay chưa xác thực email (403, mã 3009, \msg{MSG34}, kèm một OTP mới nếu còn trong giới hạn của \br{BR02}), vì vậy người không biết mật khẩu không dò được trạng thái của một email. Đăng nhập thành công đặt lại bộ đếm, ghi \texttt{last\_login\_at} và tạo phiên trong cùng một giao dịch.

Thu hồi phiên được phân biệt theo ngữ cảnh. Đăng xuất (A06, \uc{UC02}) thu hồi họ phiên của thiết bị hiện tại. Đổi mật khẩu (A11, \uc{UC03}) đòi mật khẩu cũ, thu hồi mọi phiên khác, tăng \texttt{users.token\_version} rồi cấp cặp mã mới cho chính phiên đang dùng, nên người dùng không phải đăng nhập lại trên thiết bị hiện tại. Khôi phục mật khẩu (\uc{UC03}) dùng chung luồng OTP với mục đích \texttt{reset\_password}: A07 luôn trả cùng một phản hồi dù email có tồn tại hay không; A02 xác minh OTP rồi trả \texttt{resetToken} ngẫu nhiên 256 bit, chỉ lưu băm trong \texttt{otp\_challenges}, sống 10 phút và dùng một lần; A08 nhận \texttt{email}, \texttt{resetToken}, \texttt{newPassword}, đặt lại \texttt{failed\_login\_count}, \texttt{locked\_until}, thu hồi mọi phiên và tăng \texttt{token\_version}. Khóa tài khoản bởi quản trị viên (\uc{UC24}) và phát hiện dùng lại cũng tăng \texttt{token\_version}. Middleware so claim \texttt{ver} với bản đệm \path{auth:ver:<userId>} và trả 3002 cho mã mang phiên bản cũ, nên mã bị lộ mất hiệu lực ngay, còn thiết bị hợp lệ chỉ việc làm mới (\nfr{NFR15}). Trên thiết bị, mã làm mới nằm trong Keychain hoặc Keystore qua \texttt{flutter\_secure\_storage}, mã truy cập chỉ ở bộ nhớ. OTP (\br{BR02}) được lưu dạng HMAC-SHA256 với khóa bí mật (pepper) ngoài cơ sở dữ liệu, vì $10^6$ giá trị quá dễ vét cạn nếu chỉ băm. Hình~\ref{fig:seq-dang-nhap-lam-moi} minh họa đăng nhập và làm mới phiên.

\diagram[\textwidth]{c3-25-seq-dang-nhap-lam-moi}{Biểu đồ tuần tự đăng nhập và làm mới phiên xoay vòng}{fig:seq-dang-nhap-lam-moi}

Ba lối thoát của đăng nhập (bước 3, 5--6, 7) đều kết thúc trước khi chạm vào bảng phiên, nên một lần đăng nhập thất bại không để lại phiên dở dang; việc đặt lại bộ đếm, ghi thời điểm đăng nhập và chèn phiên ở bước 8 nằm trong một giao dịch. Ở bước 13--14, câu \texttt{SELECT \ldots{} FOR UPDATE} tuần tự hóa hai lần làm mới đồng thời với cùng một mã: lần sau chờ lần trước commit rồi thấy \texttt{rotated\_at} đã có, và được xử lý theo cửa sổ ân hạn hoặc như một lần dùng lại. Hình~\ref{fig:seq-dung-lai-ma-lam-moi} trình bày trường hợp dùng lại ngoài ân hạn.

\diagram[\textwidth]{c3-25b-seq-dung-lai-ma-lam-moi}{Biểu đồ tuần tự phát hiện dùng lại mã làm mới}{fig:seq-dung-lai-ma-lam-moi}

Ở bước 2--3, RT1 đã xoay bị dùng lại ngoài cửa sổ ân hạn nên họ phiên F bị thu hồi và \texttt{token\_version} tăng trong cùng giao dịch; bước 4 cập nhật bản đệm để mọi bản API từ chối AT2 ngay ở bước 7--8 thay vì chờ 30 phút. Đây là đánh đổi có chủ ý: máy chủ không phân biệt được chủ thật giữa hai bên cùng giữ mã của một họ phiên nên chấm dứt cả hai, và người dùng thật chỉ mất một lần đăng nhập lại.

\subsection{Phân quyền}\label{ssec:phan-quyen}

Phân quyền được kiểm tra ở máy chủ cho mọi API (\nfr{NFR16}) qua ba lớp. Lớp thứ nhất là vai trò hệ thống (RBAC): nhóm \path{/admin/*} chỉ cho \texttt{system\_role = admin} và bị chặn ngay ở bộ định tuyến. Lớp thứ hai là phạm vi tài nguyên: mọi đường dẫn chứa \texttt{\{sid\}}, kể cả gợi ý tên Q02 vì kết quả gồm thực phẩm riêng của không gian, buộc người gọi là thành viên; người ngoài nhận 404, mã 4001, để không lộ sự tồn tại của không gian. Lớp thứ ba là chính sách nhóm ở tầng ứng dụng: mời, xóa thành viên, chuyển quyền, đổi tên nhóm và phân công thuộc về trưởng nhóm, còn sửa hoặc xóa công thức (R04, R05) thuộc về người tạo hoặc trưởng nhóm; người thiếu vai trò nhận 403, mã 3007 (\msg{MSG13}), và công thức mẫu chỉ được sửa qua API quản trị.

Để chống truy cập đối tượng trái phép (BOLA), mọi truy vấn repository đều lọc \texttt{space\_id} lấy từ ngữ cảnh đã kiểm tra, kể cả khi truy vấn theo khóa chính, và mọi truy vấn Elasticsearch luôn kèm bộ lọc \texttt{scope = global} hoặc \texttt{owner\_space\_id} bằng không gian hiện tại, vì chỉ mục chứa cả thực phẩm riêng và công thức của các không gian. Các API không có \texttt{\{sid\}} trong đường dẫn nhưng tác động lên đối tượng của một không gian phải tự tra chủ sở hữu rồi kiểm tra như trên: R06 chỉ nhận công thức mẫu đã công bố hoặc công thức của không gian mà người gọi là thành viên, ngược lại trả 404, mã 4001; Y03 nhận \texttt{purpose} (ảnh đại diện, không gian, công thức, thực phẩm) cùng \texttt{targetId} và chỉ cấp URL ký sẵn khi người gọi có quyền sửa đối tượng đích.

Quản trị viên kế thừa quyền của người dùng: tài khoản quản trị vẫn có không gian cá nhân và có thể tham gia một nhóm như mọi người, với quyền theo vai trò của mình trong không gian đó. Vai trò \texttt{admin} chỉ mở thêm nhóm \path{/admin/*} (tài khoản, danh mục, công thức mẫu, cấu hình, nhật ký), không thay cho tư cách thành viên, nên quản trị viên không đọc được dữ liệu của nhóm khác (\br{BR23}). Bảng~\ref{tab:c3-ma-tran-quyen} tổng hợp ma trận quyền; với các thao tác không gắn không gian, cột Thành viên đại diện cho mọi người dùng đã đăng nhập.

{\small
\newcommand{\bmok}{\tikz[baseline=-0.3ex,x=1ex,y=1ex,line width=0.16ex]\draw(0,0.1)--(0.4,-0.35)--(1.1,0.55);}
\newcommand{\bmown}{\bmok\textsuperscript{*}}
\begin{longtable}{|L{\dimexpr\textwidth-10\tabcolsep-6\arrayrulewidth-6.6cm\relax}|C{1.45cm}|C{1.65cm}|C{1.75cm}|C{1.75cm}|}
\caption{Ma trận quyền theo nhóm thao tác}\label{tab:c3-ma-tran-quyen}\\
\hline
\hd{Nhóm thao tác (mã API)} & \hd{Khách} & \hd{Thành viên} & \hd{Trưởng nhóm} & \hd{Quản trị viên}\\ \hline
\endfirsthead
\multicolumn{5}{l}{\footnotesize\itshape Bảng \thetable\ (tiếp theo)}\\ \hline
\hd{Nhóm thao tác (mã API)} & \hd{Khách} & \hd{Thành viên} & \hd{Trưởng nhóm} & \hd{Quản trị viên}\\ \hline
\endhead
Đăng ký, OTP, đăng nhập, làm mới, khôi phục (A01--A05, A07, A08) & \bmok & \bmok & \bmok & \bmok\\ \hline
Tài khoản của mình, danh sách không gian, tạo hoặc tham gia nhóm, yêu thích (A06, A09--A18, G01, G02, G07, R06, R07) & — & \bmok & \bmok & \bmok\\ \hline
Xem không gian, rời nhóm (G03, G10; \br{BR07}) & — & \bmok & \bmok & \bmown\\ \hline
Đổi tên nhóm, mã mời, xóa thành viên, chuyển quyền (G04--G06, G08, G09) & — & — & \bmok & \bmown\\ \hline
Dữ liệu không gian: mua sắm, tủ lạnh, công thức, kế hoạch, tìm kiếm, gợi ý tên, báo cáo, đồng bộ (các API S, P, R, M, B có \texttt{\{sid\}} trừ S10, S11, R04, R05; Q01, Q02, Y01) & — & \bmok & \bmok & \bmown\\ \hline
Sửa, xóa công thức của không gian (R04, R05) & — & Người tạo & \bmok & \bmown\\ \hline
Phân công mặt hàng (S10, S11; \br{BR10}) & — & — & \bmok & \bmown\\ \hline
Danh mục, tải ảnh, WebSocket (Q03, Y02--Y04) & — & \bmok & \bmok & \bmok\\ \hline
Quản trị tài khoản, danh mục, duyệt thực phẩm, công thức mẫu, cấu hình, nhật ký (D01--D11) & — & — & — & \bmok\\ \hline
\multicolumn{5}{|p{\dimexpr\textwidth-2\tabcolsep-2\arrayrulewidth\relax}|}{\footnotesize \bmown\ Chỉ trong không gian mà quản trị viên là thành viên, theo vai trò của họ ở đó; không gian của người khác trả 404, mã 4001 (\br{BR23}).}\\ \hline
\end{longtable}
}

\subsection{Bảo vệ API}\label{ssec:bao-ve-api}

Giới hạn tần suất (\nfr{NFR17}) dùng thuật toán GCRA chạy nguyên tử bằng Lua trên Redis, chủ thể đếm là IP do Caddy ghi hoặc \texttt{userId}; vượt ngưỡng nhận 429, mã 4290, kèm \texttt{Retry-After} và \msg{MSG36}, riêng gửi lại OTP dùng \msg{MSG30} (Bảng~\ref{tab:c3-gioi-han-tan-suat}). Trạng thái khóa tài khoản lưu trong PostgreSQL để khởi động lại Redis không xóa được; khi Redis lỗi, bộ đếm tạm chuyển vào bộ nhớ từng tiến trình.

{\small
\begin{longtable}{|L{3.9cm}|L{2.9cm}|L{4.3cm}|L{\dimexpr\textwidth-8\tabcolsep-5\arrayrulewidth-11.1cm\relax}|}
\caption{Các luật giới hạn tần suất}\label{tab:c3-gioi-han-tan-suat}\\
\hline
\hd{Đối tượng} & \hd{Giới hạn} & \hd{Khóa đếm} & \hd{Phản hồi}\\ \hline
\endfirsthead
\multicolumn{4}{l}{\footnotesize\itshape Bảng \thetable\ (tiếp theo)}\\ \hline
\hd{Đối tượng} & \hd{Giới hạn} & \hd{Khóa đếm} & \hd{Phản hồi}\\ \hline
\endhead
Đăng nhập theo IP (A04) & 10/phút/IP & \path{rl:login:<ip>} & 429/4290, MSG36\\ \hline
Đăng nhập sai (\br{BR03}) & 5 lần, khóa 15 phút & \texttt{users.locked\_until} & 423/3004, MSG06\\ \hline
Gửi lại OTP (A03, A07; \br{BR02}) & 1 lần/60 s, 5 lần/giờ & \path{otp:cooldown:<email>:<purpose>}, \path{otp:hourly:<email>} & 429/4290, MSG30\\ \hline
Nhập OTP sai (A02; \br{BR02}) & 5 lần mỗi mã & \texttt{failed\_attempts} & 400/3006, MSG04\\ \hline
Đăng ký (A01) & 5/giờ/IP & \path{rl:register:<ip>} & 429/4290, MSG36\\ \hline
Nhập mã mời (G07) & 10/giờ/người & \path{rl:join:<userId>} & 429/4290, MSG36\\ \hline
API chung & 120/phút/người & \path{rl:api:<userId>} & 429/4290, MSG36\\ \hline
Gợi ý tên (Q02) & 20/giây/người & \path{rl:suggest:<userId>} & 429/4290, MSG36\\ \hline
\end{longtable}
}

Khóa lũy đẳng cũng là một bề mặt tấn công nên được ràng buộc chặt. Khóa chính \texttt{(user\_id, key)} gắn khóa với người gửi, nên đoán được khóa của người khác cũng không nhận được phản hồi của họ; băm thân yêu cầu \texttt{request\_hash} chặn việc dùng lại một khóa cho nội dung khác (400, mã 2002). Khóa được chèn ở câu đầu tiên của giao dịch với trạng thái \texttt{in\_progress}, nên hai lần gửi đồng thời không cùng chạy: lần sau nhận 409, mã 4002 và thử lại sau. Chỉ kết quả thành công được lưu; lỗi làm giao dịch rollback và khóa biến mất, nên không thể “ghim” một phản hồi lỗi. API công khai A01--A05, A07, A08 không nhận \texttt{Idempotency-Key}, vì đã có giới hạn tần suất, OTP và cửa sổ ân hạn của A05, đồng thời tránh lưu phản hồi chứa mã phiên. Bảng~\ref{tab:c3-owasp} đối chiếu thiết kế với từng rủi ro của OWASP API Security Top 10 (2023).

{\small
\begin{longtable}{|L{3.3cm}|L{\dimexpr\textwidth-4\tabcolsep-3\arrayrulewidth-3.3cm\relax}|}
\caption{Ánh xạ OWASP API Security Top 10 (2023) sang biện pháp thiết kế}\label{tab:c3-owasp}\\
\hline
\hd{Rủi ro} & \hd{Biện pháp trong thiết kế}\\ \hline
\endfirsthead
\multicolumn{2}{l}{\footnotesize\itshape Bảng \thetable\ (tiếp theo)}\\ \hline
\hd{Rủi ro} & \hd{Biện pháp trong thiết kế}\\ \hline
\endhead
API1 BOLA & Kiểm tra thành viên theo \texttt{\{sid\}}, người ngoài nhận 404; mọi truy vấn lọc \texttt{space\_id}, truy vấn Elasticsearch lọc theo \texttt{scope}, \texttt{owner\_space\_id}; R06, Y03 tra chủ sở hữu đối tượng đích; khóa lũy đẳng gắn \texttt{user\_id}; kiểm thử truy cập chéo không gian.\\ \hline
API2 Lỗi xác thực & Argon2id, băm giả khi không có email; JWT cố định \texttt{EdDSA}; mã làm mới xoay vòng, phát hiện dùng lại; trạng thái tài khoản chỉ lộ sau khi mật khẩu đúng; \texttt{resetToken} dùng một lần, sống 10 phút; khóa tạm \br{BR03}.\\ \hline
API3 BOPLA & DTO tường minh sinh từ OpenAPI (\texttt{additionalProperties: false}) chống gán hàng loạt \texttt{system\_role}, \texttt{space\_id}; DTO ra không có \texttt{password\_hash}.\\ \hline
API4 Tiêu thụ tài nguyên & Thân yêu cầu $\le$ 1 MB (vượt: 413, mã 2002); \texttt{limit} $\le$ 100; ảnh qua URL ký sẵn (5 phút, $\le$ 5 MB, JPEG/PNG); báo cáo $\le$ 12 tháng; \texttt{statement\_timeout}.\\ \hline
API5 BFLA & \path{/admin/*} kiểm tra \texttt{system\_role} ở bộ định tuyến; vai trò quản trị không mở quyền vào không gian khác (\br{BR23}); chính sách trưởng nhóm, người tạo công thức ở tầng ứng dụng.\\ \hline
API6 Luồng nhạy cảm & Giới hạn riêng cho đăng ký, gửi OTP, mã mời ($32^8$ tổ hợp, 10 lần/giờ); cùng khóa lũy đẳng khác thân yêu cầu bị từ chối.\\ \hline
API7 SSRF & Danh sách host cho phép (FCM, SMTP); không tải URL do người dùng gửi.\\ \hline
API8 Cấu hình sai & TLS $\ge$ 1.2, HSTS; distroless, không chạy root; cổng vận hành nội bộ; lỗi 5000 chỉ trả \texttt{traceId}.\\ \hline
API9 Kiểm kê API & OpenAPI 3.1 là nguồn duy nhất, CI so tuyến đã đăng ký với đặc tả; phiên bản \texttt{/v1}.\\ \hline
API10 API bên thứ ba & Phản hồi \texttt{invalid\_token} của FCM làm xóa thiết bị.\\ \hline
\end{longtable}
}

\subsection{Bảo vệ dữ liệu và quyền riêng tư}\label{ssec:bao-ve-du-lieu}

Caddy~\cite{caddy} kết thúc TLS từ phiên bản 1.2 kèm HSTS (\nfr{NFR14}). Bí mật (khóa ký, pepper, mật khẩu cơ sở dữ liệu, khóa MinIO, FCM, SMTP) được cấp qua Docker secrets và biến được che của GitLab, không nằm trong Git hay image; ứng dụng dùng vai trò PostgreSQL không có quyền đổi lược đồ. Ổ đĩa mã hóa ở mức khối, kho pgBackRest~\cite{pgbackrest} mã hóa AES-256, bucket ảnh ở chế độ riêng tư. Về quyền riêng tư (\nfr{NFR18}), hệ thống chỉ thu email, họ tên, ảnh tùy chọn và mã thiết bị; log che mật khẩu, mã phiên, OTP, \texttt{resetToken} và header \texttt{Authorization}. Mọi thao tác quản trị được ghi vào bảng chỉ ghi thêm \texttt{admin.audit\_logs}.

Xóa tài khoản (A16) đòi nhập lại mật khẩu và chỉ được thực hiện khi người dùng không là trưởng nhóm của một nhóm còn thành viên khác; nếu có, máy chủ trả 422, mã 4004 (\br{BR07}, \msg{MSG24}) để người dùng chuyển quyền trước. Ngay khi xác nhận, trong một giao dịch, máy chủ thu hồi mọi phiên, tăng \texttt{token\_version}, đặt \texttt{deleted\_at} để chặn đăng nhập và xóa các thiết bị để tài khoản không còn nhận thông báo đẩy. Sau 30 ngày chờ, tác vụ nền thay email, họ tên bằng giá trị ẩn danh, xóa ảnh trên MinIO, xóa dữ liệu của không gian cá nhân cùng tài liệu Elasticsearch liên quan. Dữ liệu đã đóng góp trong nhóm như mặt hàng đã mua, lô đã nhập hay nhật ký thay đổi được giữ lại để lịch sử chung không gãy tham chiếu, và người thực hiện hiển thị là “Người dùng đã xóa”. Tương tự, khi một nhóm giải tán, tác vụ \texttt{PurgeSpace} xóa mềm dữ liệu chung và xóa tài liệu Elasticsearch của nhóm, nên dữ liệu riêng tư không còn tìm thấy được qua chỉ mục tìm kiếm.

% ---------------------------------------------------------------------
%  3.8. Triển khai, tính sẵn sàng và khả năng mở rộng
% ---------------------------------------------------------------------
\section{Triển khai, tính sẵn sàng và khả năng mở rộng}\label{sec:trien-khai}

Mục này mô tả hạ tầng giai đoạn đầu, tính sẵn sàng, cách quan sát hệ thống và lộ trình mở rộng, theo nguyên tắc mười hai yếu tố (twelve-factor app)~\cite{twelvefactor}: một image chạy ở mọi môi trường, khác biệt chỉ nằm ở biến môi trường và bí mật.

\subsection{Kiến trúc triển khai}\label{ssec:kien-truc-trien-khai}

Với đỉnh thiết kế 50 yêu cầu mỗi giây (mục~\ref{ssec:uoc-luong-tai}), giai đoạn 1 dùng một VPS Ubuntu 24.04 LTS chạy Docker Compose~\cite{docker}; Kubernetes chưa đáng chi phí vận hành ở quy mô này. Hình~\ref{fig:trien-khai} là sơ đồ triển khai UML: nút thiết bị (\texttt{«device»}) VPS chứa môi trường thực thi (\texttt{«executionEnvironment»}) Docker Engine, trong đó mỗi container là một nút \texttt{«container»}.

\diagram[\textwidth]{c3-26-trien-khai}{Sơ đồ triển khai giai đoạn 1 trên một VPS}{fig:trien-khai}

Container được chia vào ba mạng Docker theo vùng tin cậy, vẽ thành ba dải từ ngoài vào trong; mỗi container nằm ở dải của mạng ngoài cùng mà nó tham gia, dòng chữ nghiêng ghi mạng thứ hai. Mạng \texttt{edge} chỉ có Caddy, nơi duy nhất mở cổng 80, 443 ra Internet. Mạng \texttt{app} nối Caddy với \texttt{api} và MinIO, cùng \texttt{worker} và cụm giám sát. Mạng \texttt{data} là mạng \texttt{internal} của Docker, không có đường ra Internet, chứa PostgreSQL, Redis, Elasticsearch, pgBackRest và được \texttt{api}, \texttt{worker}, MinIO, cụm giám sát cùng tham gia. Caddy định tuyến theo tên miền: \path{api.dicho.example} tới hai bản \texttt{api}; \path{media.dicho.example} tới MinIO, chỉ chấp nhận URL ký sẵn; còn tên miền gốc \path{dicho.example} phục vụ tĩnh hai tệp xác minh \path{/.well-known/apple-app-site-association} và \path{/.well-known/assetlinks.json} (kiểu \texttt{application/json}, không chuyển hướng) để liên kết mời \path{https://dicho.example/join/{code}} của \uc{UC10} mở thẳng ứng dụng qua Universal Links và App Links, kèm trang dự phòng cho thiết bị chưa cài ứng dụng (mục~\ref{sec:tich-hop-di-dong}). \texttt{api} và \texttt{worker} chung một image distroless; cổng vận hành 8081 (\path{/healthz}, \path{/readyz}, \path{/metrics}) không đi qua Caddy, chỉ Prometheus trong mạng nội bộ thu chỉ số qua cổng này. Bảng~\ref{tab:c3-container-trien-khai} đặc tả các container.

{\small
\setlength{\tabcolsep}{3pt}
\begin{longtable}{|L{1.8cm}|L{2.75cm}|C{0.95cm}|L{1.95cm}|L{2.3cm}|L{2.55cm}|L{\dimexpr\textwidth-14\tabcolsep-8\arrayrulewidth-12.3cm\relax}|}
\caption{Đặc tả container triển khai giai đoạn 1}\label{tab:c3-container-trien-khai}\\
\hline
\hd{Dịch vụ} & \hd{Image} & \hd{Số bản} & \hd{Giới hạn tài nguyên} & \hd{Mạng, cổng} & \hd{Kiểm tra sức khỏe} & \hd{Volume}\\ \hline
\endfirsthead
\multicolumn{7}{l}{\footnotesize\itshape Bảng \thetable\ (tiếp theo)}\\ \hline
\hd{Dịch vụ} & \hd{Image} & \hd{Số bản} & \hd{Giới hạn tài nguyên} & \hd{Mạng, cổng} & \hd{Kiểm tra sức khỏe} & \hd{Volume}\\ \hline
\endhead
caddy & caddy 2 & 1 & 256 MB & edge, app; 80, 443 & Cổng quản trị 2019 & caddy-data; tệp xác minh liên kết (chỉ đọc)\\ \hline
api & dicho-server, lệnh \texttt{api} & 2 & 512 MB, 1 vCPU & app, data; 8080, 8081 & \texttt{/readyz} & Không\\ \hline
worker & dicho-server, lệnh \texttt{worker} & 1 & 512 MB, 1 vCPU & app, data; 8081 & \texttt{/healthz} & Không\\ \hline
postgres & postgres 16, có pgBackRest & 1 & 4 GB, 2 vCPU & data; 5432 & \texttt{pg\_isready} & pgdata\\ \hline
redis & redis 7 & 1 & 512 MB & data; 6379 & \texttt{redis-cli ping} & Không\\ \hline
elastic\-search & elasticsearch 8.15, bật xpack & 1 & Heap 1 GB, trần 2 GB & data; 9200 & \path{/_cluster/health} & esdata\\ \hline
minio & minio & 1 & 1 GB & app, data; 9000 & \path{/minio/health/live} & minio-data\\ \hline
pgbackrest & Tự dựng, có pgBackRest & 1 & 256 MB & data & \texttt{pgbackrest check} & pgdata (chỉ đọc)\\ \hline
Giám sát & prometheus, grafana, loki, alloy, hai exporter & 1 mỗi loại & Tổng 2 GB & app, data; 9090, 3000, 3100 & Điểm \texttt{ready} & prom-data, loki-data\\ \hline
\end{longtable}
}

Tổng giới hạn bộ nhớ khoảng 11,5 GB nên VPS đề xuất là 4 vCPU, 16 GB RAM, 200 GB SSD, phần dư dành cho bộ đệm trang của PostgreSQL; nếu tách cụm giám sát sang máy khác có thể dùng 8 GB RAM.

\subsection{Môi trường và quy trình CI/CD}\label{ssec:moi-truong-cicd}

Hệ thống có ba môi trường: local (Compose, Mailpit thay SMTP), staging (cùng cấu trúc với production, dùng khóa ký và dự án Firebase riêng) và production. Di trú lược đồ bằng goose chạy trước khi đổi phiên bản theo chiến lược mở rộng -- thu hẹp (expand--contract): bản phát hành đầu chỉ thêm cột, bảng, chỉ mục, cột cũ chỉ bị xóa ở bản sau, nên bản cũ và mới cùng chạy được khi cuốn chiếu và quay lui chỉ cần image trước. Pipeline chạy trên GitLab CI~\cite{gitlabci} như Hình~\ref{fig:cicd}.

\diagram[0.65\textwidth]{c3-27-cicd}{Biểu đồ hoạt động của pipeline CI/CD trên GitLab}{fig:cicd}

Mọi merge request qua bốn stage của vùng tích hợp; có job lỗi hoặc lỗ hổng mức High, Critical thì pipeline dừng và báo lỗi trên merge request. Các hành động xử lý thất bại được tô đỏ và kết thúc bằng nút kết thúc hoạt động, nên mỗi nhánh dừng đều thấy rõ trên sơ đồ. Image được gắn nhãn bằng SHA của commit để truy ngược về mã nguồn. Chỉ nhánh \texttt{main} đi tiếp tới \texttt{publish} và triển khai; staging lỗi smoke test thì chạy lại image trước, production cần phê duyệt thủ công và quay lui về image trước nếu smoke test thất bại, trong khi lược đồ đã mở rộng được giữ nguyên nhờ chiến lược expand--contract. Ứng dụng di động phát hành riêng qua kho ứng dụng, nên máy chủ luôn phát hành trước và tương thích ngược: \texttt{v1} chỉ thêm trường tùy chọn hoặc API mới, ứng dụng dưới phiên bản tối thiểu nhận 426 (mã 4260).

\subsection{Tính sẵn sàng và khả năng phục hồi}\label{ssec:san-sang-phuc-hoi}

\nfr{NFR13} đòi hỏi sẵn sàng 99\% trong khung 06:00--23:00, tức ngân sách lỗi khoảng 5 giờ mỗi tháng, đủ để chấp nhận VPS là điểm lỗi đơn nếu phục hồi nhanh. \texttt{/healthz} báo tiến trình còn sống, \texttt{/readyz} kiểm tra PostgreSQL và Redis để Caddy chỉ gửi yêu cầu tới bản sẵn sàng; riêng mất Redis sau khi khởi động chỉ báo suy giảm, tránh loại mọi bản API vì một phụ thuộc dùng chung. Docker tự khởi động lại container lỗi, tiến trình tắt mềm trong 20 giây, hai bản API cập nhật cuốn chiếu, River tự bầu lãnh đạo để tác vụ định kỳ chỉ chạy một lần. pgBackRest sao lưu đầy đủ hằng ngày, giữ 7 bản, cùng WAL liên tục ra ngoài máy chủ, cho mục tiêu RPO $\le$ 5 phút, RTO $\le$ 1 giờ, được xác nhận qua diễn tập khôi phục hằng tháng; Elasticsearch dựng lại được từ PostgreSQL. Bảng~\ref{tab:c3-che-do-loi} phân tích chế độ lỗi.

{\small
\begin{longtable}{|L{2.3cm}|L{3.5cm}|L{3.1cm}|L{\dimexpr\textwidth-8\tabcolsep-5\arrayrulewidth-8.9cm\relax}|}
\caption{Phân tích chế độ lỗi của các thành phần}\label{tab:c3-che-do-loi}\\
\hline
\hd{Thành phần} & \hd{Ảnh hưởng} & \hd{Phát hiện} & \hd{Giảm thiểu, phục hồi}\\ \hline
\endfirsthead
\multicolumn{4}{l}{\footnotesize\itshape Bảng \thetable\ (tiếp theo)}\\ \hline
\hd{Thành phần} & \hd{Ảnh hưởng} & \hd{Phát hiện} & \hd{Giảm thiểu, phục hồi}\\ \hline
\endhead
Một bản API & Yêu cầu đang chạy lỗi & Caddy kiểm tra \texttt{/readyz} & Dồn tải sang bản còn lại; ứng dụng thử lại cùng \texttt{Idempotency-Key}: giao dịch cũ đã rollback thì chạy lại, đã commit thì nhận phản hồi đã lưu\\ \hline
Worker & Thông báo, chỉ mục bị trễ & Hàng đợi River dài ra & Tác vụ nằm trong PostgreSQL (outbox), chạy tiếp khi worker sống lại\\ \hline
PostgreSQL & Mọi API lỗi (5003) & \texttt{/readyz}, exporter & Khởi động lại hoặc khôi phục theo thời điểm; ứng dụng chạy ngoại tuyến\\ \hline
Redis & Mất bộ đệm, Pub/Sub & Lỗi kết nối trong chỉ số & Đọc thẳng PostgreSQL; ứng dụng kéo \texttt{/changes}, nhận FCM\\ \hline
Elastic\-search & Tìm kiếm kém chịu lỗi gõ & Bộ ngắt mạch mở & Dự phòng \texttt{pg\_trgm}; dựng lại chỉ mục từ PostgreSQL\\ \hline
MinIO & Không tải, xem ảnh mới & \path{/minio/health/live} & Ảnh giữ chỗ, ảnh đệm; khôi phục từ bucket sao lưu\\ \hline
FCM, APNs & Thông báo đẩy trễ & Tỷ lệ gửi lỗi & Thử lại 3 lần (\nfr{NFR12}); thông báo vẫn có trong ứng dụng\\ \hline
SMTP relay & OTP không đến & Lỗi ở hàng \texttt{notify} & Thử lại; người dùng gửi lại mã (\br{BR02})\\ \hline
Open Food Facts & Không nhận diện mã mới & Quá thời gian chờ & Bộ đệm, danh mục nội bộ; mã 5004, nhập tay (\msg{MSG23})\\ \hline
Toàn bộ VPS & Gián đoạn toàn hệ thống & Đầu dò bên ngoài & Hàng đợi ngoại tuyến giữ thay đổi; dựng máy mới theo quy trình RTO\\ \hline
\end{longtable}
}

\subsection{Quan sát hệ thống}\label{ssec:quan-sat-he-thong}

Prometheus~\cite{prometheus} thu chỉ số RED (tốc độ, lỗi, độ trễ) cho mỗi tuyến API và hàng đợi River, USE (mức dùng, bão hòa, lỗi) cho CPU, bộ nhớ, đĩa và kết nối cơ sở dữ liệu. Log JSON mang \texttt{request\_id}, \texttt{trace\_id} được Grafana Alloy gửi tới Loki; vết OpenTelemetry~\cite{opentelemetry} theo một yêu cầu qua cả tác vụ River nó sinh ra. Theo SRE~\cite{sre}, ngưỡng NFR được chuyển thành chỉ báo (SLI) và mục tiêu mức dịch vụ (SLO) ở Bảng~\ref{tab:c3-slo}; cảnh báo sẵn sàng dựa trên tốc độ đốt ngân sách lỗi, kèm đầu dò bên ngoài.

{\small
\begin{longtable}{|L{6.4cm}|L{3.3cm}|L{\dimexpr\textwidth-6\tabcolsep-4\arrayrulewidth-9.7cm\relax}|}
\caption{Chỉ báo, mục tiêu mức dịch vụ và ngưỡng cảnh báo}\label{tab:c3-slo}\\
\hline
\hd{Chỉ báo SLI} & \hd{Mục tiêu SLO} & \hd{Ngưỡng cảnh báo}\\ \hline
\endfirsthead
\multicolumn{3}{l}{\footnotesize\itshape Bảng \thetable\ (tiếp theo)}\\ \hline
\hd{Chỉ báo SLI} & \hd{Mục tiêu SLO} & \hd{Ngưỡng cảnh báo}\\ \hline
\endhead
Yêu cầu API hoàn tất trong 1 s (\nfr{NFR09}) & $\ge$ 95\% & p95 > 1 s trong 10 phút\\ \hline
Độ trễ gợi ý món M05 (\nfr{NFR09}) & p95 $\le$ 2 s & p95 > 2 s trong 15 phút\\ \hline
Độ trễ gợi ý tên Q02 (\nfr{NFR09}) & p95 $\le$ 500 ms & p95 > 500 ms trong 10 phút\\ \hline
Commit tới khi gửi qua WebSocket (\nfr{NFR10}) & p95 $\le$ 1 s phía máy chủ & p95 > 1 s; mất Pub/Sub\\ \hline
Thông báo nhắc được tạo trong 15 phút sau \texttt{reminder.runAt} (\nfr{NFR12}) & $\ge$ 99\% & Còn tác vụ \texttt{UserReminder} chưa xong sau 15 phút\\ \hline
Thông báo đẩy được FCM chấp nhận sau tối đa 3 lần thử, trừ mã thiết bị không hợp lệ (\nfr{NFR12}) & $\ge$ 99\%/ngày & Tỷ lệ \texttt{deliveries} lỗi > 1\% trong 1 giờ\\ \hline
Yêu cầu không lỗi 5xx, 06:00--23:00 (\nfr{NFR13}) & $\ge$ 99\%/tháng & Đốt ngân sách > 14 lần trong 1 giờ\\ \hline
Tuổi WAL lưu trữ, bản sao lưu gần nhất (\nfr{NFR13}) & $\le$ 5 phút; $\le$ 26 giờ & Vượt ngưỡng hoặc sao lưu lỗi\\ \hline
\end{longtable}
}

\subsection{Khả năng mở rộng}\label{ssec:lo-trinh-mo-rong}

Ở đỉnh thiết kế, tiến trình Go còn dư nhiều; điểm nghẽn đầu tiên khi tải tăng là PostgreSQL (số kết nối, truy vấn báo cáo, khóa dòng khi trừ kho FEFO), đã được giảm bằng bảng tổng hợp báo cáo, bộ đệm Redis, Elasticsearch và việc đánh số thay đổi theo từng không gian; đợt tải ở giờ nhắc được cô lập bằng hàng đợi \texttt{notify}. Mở rộng ngang đã sẵn sàng vì \texttt{api} không giữ trạng thái phiên, WebSocket phát tán qua Redis Pub/Sub nên không cần phiên dính (sticky session), River cho thêm worker và lưu trữ dùng giao thức S3. Khi CPU vượt 70\% kéo dài ở giờ cao điểm, p95 tiến gần 1 giây hoặc số người dùng tiến tới 50.000, hệ thống chuyển sang giai đoạn 2 như Hình~\ref{fig:mo-rong}.

\diagram[0.9\textwidth]{c3-28-mo-rong}{Sơ đồ triển khai giai đoạn 2 với nhiều nút}{fig:mo-rong}

Giai đoạn 2 hướng tới khoảng 50.000 người dùng: cơ sở dữ liệu tách sang máy riêng với PostgreSQL primary và bản sao chỉ đọc (read replica), ít nhất ba bản API sau bộ cân bằng tải, Redis Sentinel, Elasticsearch ba nút và MinIO phân tán. Worker chạy trên hai nút ứng dụng; River tự bầu leader qua bảng \texttt{river\_leader} nên các tác vụ định kỳ như \texttt{ReminderDispatcher} vẫn chỉ được lập một lần, còn tác vụ thường được mọi worker chia nhau lấy bằng \texttt{SKIP LOCKED}. pgBackRest sao lưu từ primary cả WAL lẫn bản đầy đủ ra kho ngoài máy chủ. Bản sao chỉ đọc phục vụ báo cáo \uc{UC44}--\uc{UC47} vốn chấp nhận trễ vài giây, còn giao dịch ghi và \texttt{GET /changes} đọc primary để giữ thứ tự \texttt{seq}. Giai đoạn 3 chuyển lên Kubernetes với tự động co giãn (Horizontal Pod Autoscaler) và tách module notification thành dịch vụ riêng, vì không module nào gọi trực tiếp nó, tải dồn theo đợt và chủ yếu chờ FCM, SMTP nên cần co giãn, cô lập lỗi riêng. Đánh đổi là ranh giới giao dịch: outbox vẫn ghi trong cơ sở dữ liệu của khối nguyên khối, một tiến trình chuyển tiếp đẩy sự kiện sang message broker, và dịch vụ mới phải xử lý lũy đẳng theo \texttt{dedupe\_key}.

% ---------------------------------------------------------------------
%  3.9. Truy vết yêu cầu và kế hoạch kiểm chứng thiết kế
% ---------------------------------------------------------------------
\section{Truy vết yêu cầu và kế hoạch kiểm chứng thiết kế}\label{sec:truy-vet-thiet-ke}

Bảng~\ref{tab:c3-truy-vet} truy vết \nfr{NFR09}--\nfr{NFR23} tới cơ chế thiết kế và cách kiểm chứng; \nfr{NFR01}--\nfr{NFR08} thuộc giao diện nên được truy vết ở chương tiếp theo. Kế hoạch gồm kiểm thử đơn vị, kiểm thử tích hợp với container PostgreSQL, Redis, Elasticsearch thật (vì khóa dòng hay ràng buộc duy nhất chỉ đúng trên cơ sở dữ liệu thật), kiểm thử hợp đồng theo OpenAPI, kiểm thử tải k6, kiểm thử trên thiết bị thật và quét bảo mật; các ca phụ thuộc thời gian như nhắc hạn dùng hay xóa tài khoản dùng đồng hồ giả do \texttt{platform/clock} cấp. Kết quả sẽ được báo cáo ở chương kiểm thử.

{\small
\begin{longtable}{|L{1.35cm}|L{5.3cm}|L{2.3cm}|L{\dimexpr\textwidth-8\tabcolsep-5\arrayrulewidth-8.95cm\relax}|}
\caption{Truy vết yêu cầu phi chức năng tới thiết kế và kế hoạch kiểm chứng}\label{tab:c3-truy-vet}\\
\hline
\hd{NFR} & \hd{Quyết định, cơ chế thiết kế} & \hd{Mục trình bày} & \hd{Cách kiểm chứng}\\ \hline
\endfirsthead
\multicolumn{4}{l}{\footnotesize\itshape Bảng \thetable\ (tiếp theo)}\\ \hline
\hd{NFR} & \hd{Quyết định, cơ chế thiết kế} & \hd{Mục trình bày} & \hd{Cách kiểm chứng}\\ \hline
\endhead
NFR09 & pgx, sqlc; chỉ mục FEFO, phân trang keyset; Elasticsearch cho gợi ý tên, dự phòng \texttt{pg\_trgm}; bộ đệm \texttt{suggest:\{sid\}}; bảng tổng hợp báo cáo & \ref{sec:csdl}; Hình~\ref{fig:seq-tim-kiem-chi-muc}, \ref{fig:seq-goi-y-mon} & k6 ở 50 yêu cầu/giây với 500 công thức, 200 lô; đo p95 theo Bảng~\ref{tab:c3-slo}; tắt Elasticsearch để đo nhánh dự phòng\\ \hline
NFR10 & WebSocket, Redis Pub/Sub phát sau \texttt{COMMIT}; FCM khi ứng dụng chạy nền & Hình~\ref{fig:seq-thoi-gian-thuc}; \ref{sec:tich-hop-di-dong} & Tích hợp hai bản API, hai máy khách; đo từ lúc commit tới khi thiết bị kia nhận gợi ý thay đổi\\ \hline
NFR11, NFR23 & UUIDv7 do máy khách sinh; \texttt{Idempotency-Key} chèn ở câu đầu giao dịch; \texttt{If-Match}; \texttt{seq} theo không gian; \texttt{X-Client-Time} để kiểm \br{BR08}, \br{BR11}, \br{BR18} theo ngày thao tác trong 7 ngày; SQLite trên thiết bị & Hình~\ref{fig:seq-dong-bo-ngoai-tuyen}; \ref{sec:api}; \ref{sec:tich-hop-di-dong} & Tích hợp: phát lại cùng khóa nhận \texttt{Idempotent-Replayed}; cùng khóa khác thân nhận 2002; hai lệnh mua đồng thời cho đúng một 4002; \texttt{X-Client-Time} quá 7 ngày nhận 4004; con trỏ quá 90 ngày nhận 410\\ \hline
NFR12 & \texttt{ReminderDispatcher} chạy mỗi phút, duy nhất theo ngày; \texttt{UserReminder} duy nhất theo người và ngày; \texttt{dedupe\_key}; ngày nghiệp vụ theo giờ Việt Nam; thử lại tối đa 3 lần; tách SLI tạo và gửi & Hình~\ref{fig:seq-nhac-han-dung-be}; Bảng~\ref{tab:c3-slo} & Tích hợp với đồng hồ giả: chạy lại trong ngày và người thuộc hai không gian vẫn chỉ nhận một thông báo; khung 00:00--07:00 giờ Việt Nam tính đúng ngày; FCM giả lỗi tạm được thử lại 3 lần, \texttt{invalid\_token} làm xóa thiết bị\\ \hline
NFR13 & Giao dịch ACID, thứ tự khóa toàn cục, outbox River; hai bản API cuốn chiếu; pgBackRest & Hình~\ref{fig:seq-xac-nhan-nau}; Bảng~\ref{tab:c3-thu-tu-khoa}; \ref{ssec:san-sang-phuc-hoi} & \uc{UC41} chạy đồng thời không làm âm kho, không bế tắc; k6 trong lúc cuốn chiếu; diễn tập khôi phục hằng tháng\\ \hline
NFR14, NFR15 & Argon2id, TLS $\ge$ 1.2; JWT EdDSA; mã làm mới xoay vòng, phát hiện dùng lại; \texttt{token\_version}; A11 giữ phiên hiện tại, A08 thu hồi mọi phiên & \ref{ssec:xac-thuc-phien}; Hình~\ref{fig:seq-dang-nhap-lam-moi}, \ref{fig:seq-dung-lai-ma-lam-moi} & Đơn vị cho băm và ký mã; tích hợp luồng dùng lại (3008) và đổi mật khẩu thu hồi phiên khác; quét TLS; quét log tìm bí mật\\ \hline
NFR16 & Ba lớp phân quyền; \texttt{space\_id} ở mọi truy vấn; lọc không gian trên Elasticsearch; vai trò quản trị không mở quyền vào nhóm khác & \ref{ssec:phan-quyen}; Hình~\ref{fig:seq-roi-nhom-chuyen-quyen}, \ref{fig:seq-duyet-thuc-pham} & Ca truy cập chéo sinh từ OpenAPI theo Bảng~\ref{tab:c3-ma-tran-quyen}; tài khoản quản trị gọi API của nhóm khác nhận 404; OWASP ZAP\\ \hline
NFR17 & GCRA trên Redis; khóa tạm \br{BR03}; OTP lưu HMAC, giới hạn \br{BR02} & \ref{ssec:bao-ve-api}; Hình~\ref{fig:seq-dang-ky-otp} & Vượt ngưỡng nhận 429 (\msg{MSG36}, \msg{MSG30}); lần sai thứ năm nhận 423; OTP sai năm lần bị vô hiệu\\ \hline
NFR18 & Tối thiểu dữ liệu, che log; A16 nhập lại mật khẩu, thu hồi phiên ngay, ẩn danh sau 30 ngày; \texttt{PurgeSpace} khi giải tán nhóm & \ref{ssec:bao-ve-du-lieu} & Tích hợp với đồng hồ giả: đăng nhập bị chặn ngay; sau 30 ngày email, họ tên được ẩn danh, ảnh và tài liệu Elasticsearch bị xóa, dữ liệu nhóm hiển thị “Người dùng đã xóa”; trưởng nhóm còn thành viên nhận 4004\\ \hline
NFR19 & Một API FCM HTTP v1 cho cả hai nền tảng, iOS chuyển tiếp qua APNs; header \texttt{X-Platform}; khóa gộp \texttt{apns-collapse-id}, \texttt{collapse\_key}; WebSocket ở tiền cảnh, FCM khi nền; mã 426 buộc cập nhật & \ref{sec:tich-hop-di-dong}; Hình~\ref{fig:seq-thoi-gian-thuc} & Trên thiết bị thật Android 8 và iOS 13: nhận đẩy khi ứng dụng ở nền và đã tắt, thông báo cùng khóa được gộp, bản dưới phiên bản tối thiểu nhận 426\\ \hline
NFR20 & Khối nguyên khối chia module, lục giác; tham số nghiệp vụ đọc lúc chạy, đổi qua \texttt{PUT /admin/configs}; OpenAPI làm trước & \ref{sec:kien-truc}, \ref{sec:thiet-ke-chi-tiet}; Bảng~\ref{tab:c3-tham-so}; Hình~\ref{fig:seq-quan-tri-cau-hinh} & Lint phụ thuộc module; bao phủ tầng nghiệp vụ $\ge$ 70\% là cổng của CI; đổi tham số có hiệu lực mà không triển khai lại\\ \hline
NFR21 & REST/JSON tại \texttt{/it4788/v1}; bao phản hồi và bộ mã thống nhất; ánh xạ ràng buộc cơ sở dữ liệu sang mã & \ref{sec:api}; Bảng~\ref{tab:c3-anh-xa-loi}; \ref{ssec:dac-ta-api} & Kiểm thử hợp đồng mọi phản hồi theo OpenAPI; mỗi ràng buộc bị vi phạm trả đúng mã theo bảng ánh xạ; chạy các test case dự kiến\\ \hline
NFR22 & PostgreSQL, MinIO, FCM HTTP v1; River làm bộ lập lịch phía máy chủ & \ref{ssec:kien-truc-trien-khai} & Tích hợp với container thật trong CI; với hai worker, tác vụ định kỳ chỉ chạy một lần\\ \hline
\end{longtable}
}

Tóm lại, chương này đã thiết kế phần máy chủ theo kiến trúc khối nguyên khối chia module viết bằng Go, lấy PostgreSQL làm nguồn sự thật duy nhất, kết hợp outbox River, đồng bộ ngoại tuyến lũy đẳng, thời gian thực qua WebSocket và FCM, đối chiếu với kiến trúc tham chiếu của Android, iOS và Flutter, bảo mật nhiều lớp và hạ tầng có sao lưu, quan sát, lộ trình mở rộng. Mỗi yêu cầu phi chức năng đều có cơ chế thiết kế và cách kiểm chứng tương ứng. Chương tiếp theo trình bày thiết kế giao diện, nơi các API và cơ chế đồng bộ này trở thành màn hình cụ thể.

\endgroup
